% =====================================================================
%  HIDRÁULICA APLICADA -- CIV-346
%  Apuntes transcritos de las clases (Unidad I, II y Drenaje superficial)
%  Compilar con: pdflatex main.tex  (dos veces, para el índice)
% =====================================================================
\documentclass[11pt,a4paper]{report}

% ---------------------------------------------------------------------
%  Paquetes
% ---------------------------------------------------------------------
\usepackage[utf8]{inputenc}
\usepackage[T1]{fontenc}
\usepackage[spanish,es-noshorthands,es-tabla]{babel}
\usepackage{lmodern}
\usepackage[margin=2.3cm]{geometry}
\usepackage{amsmath,amssymb,amsfonts}
\usepackage{mathtools}
\usepackage{siunitx}
\sisetup{output-decimal-marker={.}}
\usepackage{graphicx}
\usepackage{booktabs,array,multirow,tabularx}
\usepackage[table]{xcolor}
\usepackage{enumitem}
\usepackage{caption,subcaption}
\usepackage{float}
\usepackage{cancel}
\usepackage[most]{tcolorbox}
\usepackage{tikz}
\usetikzlibrary{arrows.meta,calc,patterns,decorations.pathmorphing,
  decorations.markings,shapes.geometric,positioning,fadings,shadings,
  intersections,backgrounds}
\usepackage{pgfplots}
\pgfplotsset{compat=1.16}
\usepgfplotslibrary{fillbetween}
\usepackage{xurl}
\usepackage{hyperref}
\hypersetup{colorlinks=true,linkcolor=azulusm,urlcolor=azulusm}

% ---------------------------------------------------------------------
%  Colores y estilos
% ---------------------------------------------------------------------
\definecolor{azulusm}{RGB}{0,69,137}
\definecolor{amarillo}{RGB}{247,176,0}
\definecolor{agua}{RGB}{110,165,215}
\definecolor{aguaclara}{RGB}{190,215,240}
\definecolor{suelo}{RGB}{200,160,70}
\definecolor{sueloclaro}{RGB}{240,215,150}
\definecolor{roca}{RGB}{190,190,195}
\definecolor{rojo}{RGB}{210,20,30}

% Caja amarilla para fórmulas importantes (como en las diapositivas)
\newtcolorbox{formula}[1][]{colback=amarillo!35,colframe=amarillo,
  arc=3mm,boxrule=0.8pt,left=4mm,right=4mm,top=1mm,bottom=1mm,
  before upper={\centering},#1}
% Caja celeste para desarrollos / notas
\newtcolorbox{nota}[1][]{colback=aguaclara!40,colframe=agua,
  arc=2mm,boxrule=0.6pt,fonttitle=\bfseries,#1}

\setlist[itemize,1]{label=\textcolor{rojo}{$\odot$}}
\setlist[itemize,2]{label=\textcolor{rojo}{$\square$}}

% Estilos TikZ reutilizables
\tikzset{
  nivelagua/.pic={%  símbolo de nivel de agua ▽
    \draw[fill=azulusm] (-0.12,0.18)--(0.12,0.18)--(0,0)--cycle;
    \draw (-0.15,-0.05)--(0.15,-0.05) (-0.10,-0.10)--(0.10,-0.10)
          (-0.05,-0.15)--(0.05,-0.15);},
  impermeable/.style={pattern=north east lines,pattern color=black!60},
  flecha/.style={-{Stealth[length=2.5mm]},thick},
  cota/.style={{Stealth[length=2mm]}-{Stealth[length=2mm]}},
}

% \paragraph como título en línea propia (evita cajas desbordadas)
\makeatletter
\renewcommand\paragraph{\@startsection{paragraph}{4}{\z@}%
  {1.5ex \@plus 0.5ex \@minus 0.2ex}{0.5ex}{\normalfont\normalsize\bfseries}}
\makeatother

\newcommand{\dd}{\mathrm{d}}
\newcommand{\clase}[2]{\chapter[Clase #1: #2]{#2}\noindent\textit{Clase #1}\par\medskip}

\title{\textbf{Hidráulica Aplicada -- CIV-346}\\[4pt]
  \large Apuntes de clases: Hidrología subterránea, Aguas subterráneas\\
  y Drenaje superficial}
\author{}
\date{}

\begin{document}
\maketitle
\tableofcontents

\part{Hidrología y aguas subterráneas}
% =====================================================================
%  CLASE 01 -- Definiciones, propiedades de suelos, Ley de Darcy
% =====================================================================
\clase{01}{Hidrología subterránea: definiciones, propiedades de suelos y Ley de Darcy}

\section{Ocurrencia de agua subterránea}

Al estudiar el agua subterránea dentro del ciclo hidrológico surgen tres
preguntas fundamentales:
\begin{itemize}
  \item ¿Cómo llega el agua al suelo?
  \item ¿Cómo se transmite y almacena?
  \item ¿De qué factores dependerá?
\end{itemize}

\begin{figure}[H]
\centering
\begin{tikzpicture}[scale=0.9,font=\scriptsize]
  % terreno
  \fill[sueloclaro] (0,0) -- (0,2.2) -- (1.5,2.4) -- (3,4.2) -- (4.3,2.6)
        -- (6.5,2.0) -- (8.5,1.6) -- (10,1.2) -- (10,0) -- cycle;
  \fill[roca!70] (0,0) rectangle (10,0.6);
  \node[white] at (5,0.3) {Estratos impermeables};
  % oceano
  \fill[agua] (8.5,1.6) -- (10,1.2) -- (10,0.6) -- (8.3,0.6) -- cycle;
  \node at (9.3,1.0) {Océano};
  % nivel freático
  \draw[dashed,azulusm,thick] (0,1.6) .. controls (3,2.4) and (5,1.6) .. (8.5,1.3);
  \node[azulusm] at (4.6,1.95) {Nivel freático};
  % nubes y precipitación
  \foreach \x in {2,5.5,8.5}{
    \fill[gray!40] (\x,5.2) circle (0.35) (\x+0.4,5.3) circle (0.4) (\x+0.8,5.2) circle (0.35);
    \foreach \d in {0,0.3,0.6,0.9}
      \draw[agua,thick] (\x+\d-0.1,4.7) -- (\x+\d-0.25,4.35);}
  \node at (2.4,5.9) {Precipitación nival};
  \node at (5.9,5.9) {Precipitación pluvial};
  \node at (9.3,5.9) {Precipitación oceánica};
  % evaporación
  \draw[flecha,agua] (7.2,2.2) -- (7.2,3.4) node[right,black,align=left]{Evaporación y\\evapotranspiración};
  \draw[flecha,agua] (9.4,1.5) -- (9.4,3.9);
  % escorrentía y flujo subterráneo
  \draw[flecha,azulusm] (3.3,3.7) .. controls (4,2.8) .. (6,2.1) node[midway,above,sloped,black]{Escorrentía superficial};
  \draw[flecha,azulusm,dashed] (2,1.2) -- (7.5,1.0) node[midway,below,black]{Flujo subterráneo};
  \draw[flecha,azulusm] (1,2.3) -- (1,1.7) node[below,black]{Percolación};
  \node at (5,6.4) {\textbf{CICLO HIDROLÓGICO}};
\end{tikzpicture}
\caption{Esquema del ciclo hidrológico y ocurrencia del agua subterránea.}
\end{figure}

\section{Definiciones}

\subsection{Aguas subterráneas}
Agua existente bajo la superficie terrestre, que satura los intersticios
(``vacíos'') del terreno (``suelo'').

\begin{itemize}
  \item \textbf{Origen atmosférico}
  \item Proveniente de la \textbf{infiltración directa} de lluvias y nieves.
  \item \textbf{Aportes indirectos} de ríos y lagos.
  \item \textbf{Espesor de la zona capilar}:
  \begin{itemize}
    \item En suelos gruesos, arenas y gravas, del orden de mm o a lo más cm.
    \item En suelos finos, limos y principalmente arcillas, puede superar la
          decena de metros.
  \end{itemize}
\end{itemize}

\begin{figure}[H]
\centering
\begin{tikzpicture}[font=\small]
  % columna izquierda
  \fill[white] (0,0) rectangle (2.4,4);
  \fill[agua!25] (0,0) rectangle (2.4,1.6);
  \fill[agua] (2.4,0) rectangle (4.4,1.6);
  \fill[agua!25] (4.4,0) rectangle (6.8,1.6);
  \fill[gray!40] (0,-0.2) rectangle (6.8,0);
  \draw[dashed] (2.4,0) -- (2.4,4) (4.4,0) -- (4.4,4);
  \draw (0,4) -- (2.4,4) (4.4,4) -- (6.8,4);
  \draw[dashed] (0,2.6) -- (6.8,2.6);
  \draw[thick,azulusm] (0,1.6) -- (6.8,1.6);
  \foreach \i in {1,...,70}{
    \pgfmathsetmacro{\x}{rnd*2.3+0.05}\pgfmathsetmacro{\y}{rnd*1.5+0.05}
    \fill[azulusm] (\x,\y) circle (0.6pt);
    \fill[azulusm] (\x+4.4,\y) circle (0.6pt);}
  \foreach \i in {1,...,30}{
    \pgfmathsetmacro{\x}{rnd*2.3+0.05}\pgfmathsetmacro{\y}{rnd*2.3+1.65}
    \fill[black!60] (\x,\y) circle (0.5pt);
    \fill[black!60] (\x+4.4,\y) circle (0.5pt);}
  \node[align=center,font=\scriptsize] at (1.2,3.35) {Zona no aireación\\(agua y aire en poros)};
  \node[font=\scriptsize] at (1.2,2.1) {Zona capilar};
  \node[align=center,font=\scriptsize] at (1.2,0.8) {Zona saturada\\(sólo agua)};
  \node[font=\scriptsize] at (5.6,4.2) {Superficie};
  \node[font=\scriptsize] at (5.6,3.3) {Agua en el suelo};
  \node[font=\scriptsize] at (5.5,1.85) {Nivel freático};
  \pic at (6.55,1.62) {nivelagua};
  \node[font=\scriptsize] at (5.6,0.8) {Agua subterránea};
  \draw[decorate,decoration={brace,amplitude=5pt}] (-0.1,1.65) -- (-0.1,4)
     node[midway,left=6pt,align=center,font=\scriptsize]{Zona\\no saturada};
\end{tikzpicture}
\caption{Distribución vertical del agua en el subsuelo.}
\end{figure}

\subsection{Porosidad y porosidad efectiva}

\begin{itemize}
  \item \textbf{Porosidad:} máxima fracción volumétrica de agua en el suelo.
  \item \textbf{Porosidad efectiva:} fracción volumétrica que participa en el flujo.
\end{itemize}

\begin{formula}
$\displaystyle n=\frac{V_{\text{vacíos}}}{V_{\text{total}}}
\qquad\qquad
n_e=\frac{V_{\text{drenado}}}{V_{\text{total}}}$
\end{formula}

Valores típicos de porosidad según la estructura del suelo:
\begin{center}
\begin{tabular}{lc@{\hspace{1.5cm}}lc}
\toprule
Arena bien graduada & $n\approx 0.4$ & Arena mal graduada & $n\approx 0.2$\\
Arcilla sin compactar & $n\approx 0.7$ & Arcilla compactada & $n\approx 0.5$\\
\bottomrule
\end{tabular}
\end{center}

\begin{figure}[H]
\centering
\begin{minipage}{0.55\textwidth}
\centering
\begin{tikzpicture}
\begin{axis}[width=\linewidth,height=5.5cm,xmin=0,xmax=10,ymin=0,ymax=60,
  xticklabels={},xtick=\empty,ylabel={Volumen (\%)},xlabel={Tamaño del grano $\longrightarrow$},
  axis lines=left,font=\small,clip=false]
  \addplot[thick,black,smooth] coordinates {(0,52)(3,49)(6,45)(9,37)};
  \addplot[thick,black,smooth] coordinates {(0,47)(2,40)(4,25)(6,14)(9,6)};
  \addplot[thick,agua!80!black,smooth] coordinates {(0,4)(2,12)(4,24)(6,31)(7.5,33)(9,31)};
  \node[rojo,anchor=west,font=\scriptsize] at (axis cs:9.1,37){Porosidad};
  \node[rojo,anchor=west,font=\scriptsize,align=left] at (axis cs:9.1,30){Agua disponible\\para el flujo};
  \node[rojo,anchor=west,font=\scriptsize,align=left] at (axis cs:9.1,6){Agua no disponible\\para flujo};
\end{axis}
\end{tikzpicture}
\end{minipage}\hfill
\begin{minipage}{0.35\textwidth}
\centering
\begin{tabular}{lcc}
\toprule
Tipo de suelo & $n\,[\%]$ & $n_e\,[\%]$\\
\midrule
Arcillas & 50 & 4\\
Limos & 48 & 8\\
Arenas & 40 & 33\\
Gravas & 25 & 22\\
Rocas & $1\sim 8$ & $0.5\sim 4$\\
\bottomrule
\end{tabular}
\end{minipage}
\caption{Porosidad total, agua retenida (no disponible) y agua disponible
para el flujo en función del tamaño del grano.}
\end{figure}

\subsection{Clasificación de formaciones geológicas}
\begin{itemize}
  \item \textbf{Acuíferos}
  \begin{itemize}
    \item Formación geológica (gravas, arenas) que puede almacenar y transmitir
          agua (drenaje alto).
  \end{itemize}
  \item \textbf{Acuitardos}
  \begin{itemize}
    \item Formación geológica (limos, limos arcillosos, arcillas limosas) que
          puede almacenar, pero con mayor dificultad para transmitir agua
          (drenaje medio/bajo).
    \item Bajo condiciones especiales permiten una recarga vertical de otros acuíferos.
  \end{itemize}
  \item \textbf{Acuicludos}
  \begin{itemize}
    \item Formación geológica (arcillas, arcillas plásticas, limos arcillosos)
          con gran capacidad de almacenar, pero que no tiene posibilidad de
          transmitir (drenaje dificultoso).
  \end{itemize}
  \item \textbf{Acuífugos}
  \begin{itemize}
    \item Formación geológica impermeable (rocas), que no puede almacenar ni
          transmitir agua, salvo que existan fracturas que permitan flujos.
  \end{itemize}
\end{itemize}

\begin{table}[H]
\centering
\rowcolors{2}{azulusm!8}{white}
\begin{tabularx}{\textwidth}{l c c c X}
\toprule
\textbf{Tipo} & \textbf{Cap. almacenar} & \textbf{Cap. drenar} &
\textbf{Cap. transmitir} & \textbf{Formaciones características}\\
\midrule
Acuíferos  & Alta & Alta       & Alta & Gravas, arenas\\
Acuitardos & Alta & Media/Baja & Baja & Limos, limos arcillosos, arcillas limosas\\
Acuicludos & Alta & Muy baja   & Nula & Arcillas\\
Acuífugos  & Nula & Nula       & Nula & Rocas\\
\bottomrule
\end{tabularx}
\caption{Resumen de formaciones geológicas.}
\end{table}

Los que nos interesan son los \textbf{acuíferos}, caracterizados por su
\emph{extensión} (km$^2$) y su \emph{potencia} (m).

\begin{formula}
\textbf{Permeabilidad:} medida de la capacidad de un material poroso para transmitir agua.
\end{formula}

\section{Carga hidráulica y potencial piezométrico}

Antes de definir los tipos de acuíferos es necesario recordar la definición
de carga hidráulica. El nivel de energía en un punto se define utilizando la
ley de Bernoulli:
\begin{formula}
$\displaystyle H = z + \frac{p}{\gamma} + \frac{v^2}{2g}$
\end{formula}
Debido a que las velocidades en flujo subterráneo son bajas,
$\frac{v^2}{2g}\approx 0$, y por lo tanto:
\[
H = z + \frac{p}{\gamma} + \cancelto{\approx 0}{\frac{v^2}{2g}}
\quad\Longrightarrow\quad
\boxed{h = z + \frac{p}{\gamma}}\qquad h:\ \text{carga hidráulica}
\]

\begin{figure}[H]
\centering
\begin{tikzpicture}[font=\small]
  \draw[thick] (-1.5,4) -- (1.5,4);
  \fill[gray!20] (-1.5,0) rectangle (1.5,-0.2);
  \draw (-1.5,0) -- (1.5,0) node[right]{$z=0$};
  \draw (-0.15,1.2) rectangle (0.15,4.6);
  \fill[agua] (-0.15,1.2) rectangle (0.15,3.2);
  \draw[cota] (-0.8,0) -- (-0.8,3.2) node[midway,left]{$h$};
  \draw[cota] (0.7,0) -- (0.7,1.2) node[midway,right]{$z$};
  \draw[cota] (0.7,1.2) -- (0.7,3.2) node[midway,right]{$\dfrac{p}{\rho g}$};
  \draw[dashed] (-0.8,3.2) -- (0.7,3.2);
\end{tikzpicture}
\caption{Carga hidráulica medida en un piezómetro.}
\end{figure}

Para que exista flujo en un medio poroso se requiere la existencia de un
\textbf{gradiente de energía o potencial piezométrico}. El agua fluye desde
puntos de mayor carga $h_1$ hacia puntos de menor carga $h_2$.

\begin{figure}[H]
\centering
\begin{tikzpicture}[font=\small]
  \fill[pattern=crosshatch,pattern color=gray!50] (0,-0.3) rectangle (9,0);
  \draw (0,0) -- (9,0) node[right]{$z=0$};
  \draw (0,3.2) -- (4,3.2) (5,3.2) -- (9,3.2);
  \fill[agua] (4,0) rectangle (5,1.8);
  \draw (4,0) -- (4,3.2) (5,0) -- (5,3.2);
  \draw[dashed,azulusm,thick] (0,3.0) .. controls (2,2.9) and (3.5,2.3) .. (4,1.8);
  \draw[dashed,azulusm,thick] (5,1.8) .. controls (5.5,2.3) and (7,2.9) .. (9,3.0);
  \foreach \x/\h/\n in {0.8/2.95/1,2.4/2.65/2}{
    \draw (\x-0.12,1.0) rectangle (\x+0.12,3.5);
    \fill[agua] (\x-0.12,1.0) rectangle (\x+0.12,\h);
    \draw[cota] (\x-0.4,0) -- (\x-0.4,\h) node[midway,left]{$h_\n$};}
  \draw[flecha,azulusm] (1.4,0.8) -- (2.0,0.8);
\end{tikzpicture}
\caption{El flujo se produce en la dirección de disminución de la carga hidráulica ($h_1>h_2$).}
\end{figure}

\section{Tipos de acuíferos}

\begin{itemize}
  \item \textbf{Acuífero libre:} cercanos a la superficie terrestre; su límite
        superior es el nivel freático (superficie con presión atmosférica).
  \item \textbf{Acuífero artesiano o confinado:} confinados entre estratos
        impermeables, por ende, bajo presión.
  \item \textbf{Acuífero semi-confinado:}
  \begin{itemize}
    \item Confinados entre estratos semi-permeables.
    \item Bajo presión, pero una o ambas capas permiten alguna filtración.
  \end{itemize}
\end{itemize}

\begin{figure}[H]
\centering
\begin{tikzpicture}[font=\scriptsize,scale=1.0]
  % capa impermeable inferior
  \fill[suelo!80!black] (0,0) rectangle (11,0.8);
  \node at (7,0.4) {Capa impermeable};
  % acuífero cautivo
  \fill[gray!45] (0,0.8) -- (11,0.8) -- (11,1.9) .. controls (7,1.9) and (4,1.5) .. (1.5,4.0) -- (0,4.0) -- cycle;
  \node at (6.5,1.3) {Acuífero cautivo (confinado)};
  % capa confinante
  \fill[suelo] (1.5,4.0) .. controls (4,1.5) and (7,1.9) .. (11,1.9) -- (11,2.2)
      .. controls (7,2.2) and (4.2,1.9) .. (2.0,4.3) -- cycle;
  \node at (4.2,2.05) {Capa confinante};
  % acuífero libre
  \fill[agua!40] (2.0,4.3) .. controls (4.2,1.9) and (7,2.2) .. (11,2.2) -- (11,3.3) -- (3.3,3.3) -- cycle;
  \node at (9.2,2.7) {Acuífero libre};
  % zona no saturada
  \fill[sueloclaro] (3.3,3.3) -- (11,3.3) -- (11,4.0) .. controls (8,4.2) and (6,3.8) .. (4,4.4) -- (2.0,4.3) -- cycle;
  \draw[dashed] (0,3.5) -- (11,3.5);
  \node[align=center] at (0.9,3.7) {Nivel de agua\\subterránea};
  \node at (9.8,3.55) {Nivel freático};
  % zona de recarga
  \fill[sueloclaro] (0,4.0) -- (1.5,4.0) -- (2.0,4.3) -- (0,4.9) -- cycle;
  \foreach \x in {0.3,0.6,0.9,1.2}
    \draw[flecha,rojo] (\x,5.7) -- (\x,4.9);
  \node at (0.8,5.95) {Zona de recarga};
  % pozos
  \draw[line width=2pt,azulusm] (4.8,4.1) -- (4.8,1.2);
  \foreach \a in {60,90,120} \draw[azulusm] (4.8,4.1) -- ++(\a:0.4);
  \node[align=center] at (4.8,4.85) {Pozo artesiano\\surgente};
  \draw[line width=2pt,azulusm] (8.0,4.1) -- (8.0,1.2);
  \node[align=center] at (8.0,4.85) {Pozo artesiano\\no surgente};
\end{tikzpicture}
\caption{Acuífero libre y acuífero confinado (cautivo). El nivel piezométrico
del acuífero confinado puede quedar sobre el terreno (pozo surgente).}
\end{figure}

\begin{figure}[H]
\centering
\begin{tikzpicture}[font=\small]
  \fill[pattern=crosshatch,pattern color=black!60] (0,0) rectangle (8,0.3);
  \node[fill=white,inner sep=1pt] at (7,0.15) {Impermeable};
  \fill[sueloclaro] (0,0.3) rectangle (8,1.6);
  \node at (1.8,0.95) {Acuífero semiconfinado};
  \fill[gray!35] (0,1.6) rectangle (8,2.4);
  \node at (1.8,2.0) {Acuitardo};
  \fill[agua!25] (0,2.4) rectangle (8,3.4);
  \node at (1.8,2.9) {Acuífero libre};
  \fill[sueloclaro!60] (0,3.4) rectangle (8,4.2);
  \draw[thick,azulusm] (0,3.4) -- (8,3.4) node[right,align=left,font=\scriptsize,yshift=-6pt]{Superficie freática\\(acuífero libre superior)};
  \draw[thick,rojo,dashed] (0,3.7) -- (8,3.7) node[right,align=left,font=\scriptsize,yshift=10pt]{Superficie piezométrica\\(acuífero semiconfinado)};
  \foreach \x in {3.5,5,6.5} \draw[flecha,azulusm] (\x,2.35) -- (\x,1.65);
\end{tikzpicture}
\caption{Acuífero semiconfinado: la recarga vertical ocurre a través del acuitardo.}
\end{figure}

\section{Explotación de acuíferos}

Para poder aprovechar los recursos de aguas subterráneas es necesario que
ésta aflore a la superficie. Para esto se presentan dos alternativas:
\begin{itemize}
  \item \textbf{Vertientes:} lugares donde, debido a la configuración geológica
        y geomorfológica de los suelos, el agua aflora naturalmente a la superficie.
  \item \textbf{Pozos o drenes.}
\end{itemize}

\subsection{Tipos de vertientes}
\begin{itemize}
  \item \textbf{Vertientes de depresión:} generalmente cajones de ríos.
  \item \textbf{Vertientes de contacto.}
  \item \textbf{Vertientes de fractura:} normalmente en zonas cordilleranas.
  \item \textbf{Vertientes de dique.}
\end{itemize}

\begin{figure}[H]
\centering
\begin{subfigure}{0.48\textwidth}\centering
\begin{tikzpicture}[scale=0.75,font=\scriptsize]
  \fill[top color=agua!30,bottom color=azulusm!80] (0,0) -- (0,3) -- (2.8,3) -- (3.6,2.4) -- (4.4,2.4) -- (5.2,3) -- (8,3) -- (8,0) -- cycle;
  \fill[gray!25] (0,3) -- (0,3.4) -- (2.2,3.4) -- (3.4,2.55) -- (2.8,3) -- cycle;
  \fill[gray!25] (8,3) -- (8,3.4) -- (5.8,3.4) -- (4.6,2.55) -- (5.2,3) -- cycle;
  \draw (0,3.4) -- (2.2,3.4) -- (3.4,2.55) .. controls (3.8,2.3) and (4.2,2.3) .. (4.6,2.55) -- (5.8,3.4) -- (8,3.4);
  \foreach \x in {0.3,0.8,...,2.0} \draw[flecha] (\x,4.1) -- (\x,3.45);
  \foreach \x in {6.0,6.5,...,7.8} \draw[flecha] (\x,4.1) -- (\x,3.45);
  \draw[impermeable] (0,-0.3) rectangle (8,0);
  \draw[<-] (4.4,2.45) -- (6,2.2) node[right]{Recuperación};
\end{tikzpicture}
\caption{Vertiente de depresión.}
\end{subfigure}\hfill
\begin{subfigure}{0.48\textwidth}\centering
\begin{tikzpicture}[scale=0.75,font=\scriptsize]
  \fill[gray!30] (0,0) -- (1.5,0) -- (2.5,1.3) -- (7,1.3) -- (7,0) -- cycle;
  \fill[gray!10,draw=black] (2.5,1.3) -- (7,1.3) -- (7,1.8) -- (2.2,1.6) -- cycle;
  \fill[agua] (2.2,1.6) -- (7,1.8) -- (7,2.4) -- (4,2.2) -- cycle;
  \fill[gray!30] (2.2,1.6) -- (4,2.2) -- (7,2.4) -- (7,3.6) -- (5,3.4) -- cycle;
  \draw (1.5,0) -- (2.5,1.3) (2.2,1.6) -- (5,3.4) -- (7,3.6) -- (7,0);
  \fill[agua] (2.2,1.6) -- (1.7,1.75) -- (2.0,1.5) -- cycle;
  \draw[<-] (2.0,1.7) -- (1.3,2.6) node[above]{Afloramiento de agua};
  \draw[<-] (7,1.55) -- (7.6,1.2) node[below]{Material impermeable};
  \draw[<-] (5.8,3.5) -- (6.2,4.1) node[above]{Superficie del suelo};
\end{tikzpicture}
\caption{Vertiente de contacto.}
\end{subfigure}

\medskip
\begin{subfigure}{0.48\textwidth}\centering
\begin{tikzpicture}[scale=0.75,font=\scriptsize]
  \fill[azulusm!70] (1.2,0) -- (1.2,0.4) -- (8,2.5) -- (8,0) -- cycle;
  \fill[gray!25] (1.2,0.4) -- (8,2.5) -- (8,4) -- (5,3.3) -- (3,2.6) -- (1.3,0.6) -- cycle;
  \draw[thick,gray] (1.3,0.6) -- (3,2.6) -- (5,3.3) -- (8,4);
  \foreach \a/\b in {3.0/2.5,4.5/3.1,6/3.6,7.2/3.9}
     \draw[azulusm] (\a,\b) -- (\a+0.7,\b-1.6);
  \fill[agua] (1.2,0.5) -- (0.6,0.8) -- (0.8,0.4) -- cycle;
  \node[above] at (0.8,1.2) {Afloramiento};
  \node at (4.5,4.0) {Fracturas};
  \node[right] at (5.8,4.5) {Material impermeable};
\end{tikzpicture}
\caption{Vertiente de fractura.}
\end{subfigure}\hfill
\begin{subfigure}{0.48\textwidth}\centering
\begin{tikzpicture}[scale=0.75,font=\scriptsize]
  \fill[gray!40] (0,0) rectangle (8,1.2);
  \fill[agua] (0.3,1.8) .. controls (1.5,0.4) and (2.6,0.4) .. (3.3,1.8) -- cycle;
  \fill[agua] (5.2,1.8) .. controls (5.8,0.3) and (7,0.3) .. (7.7,1.8) -- cycle;
  \fill[sueloclaro] (0.3,1.8) -- (3.3,1.8) -- (3.5,2.1) -- (0.2,2.0) -- cycle;
  \fill[sueloclaro] (5.2,1.8) -- (7.7,1.8) -- (7.8,2.2) -- (5,2.1) -- cycle;
  \fill[gray!40] (0,1.2) -- (0,1.9) .. controls (1.5,0.2) and (2.8,0.2) .. (3.4,1.8)
      -- (4.2,4) -- (5.1,1.8) .. controls (5.8,0) and (7,0) .. (8,2.2) -- (8,1.2) -- cycle;
  \node at (4.2,1.8) {Roca basal};
  \fill[agua] (0,1.9) -- (-0.5,2.1) -- (-0.3,1.8) -- cycle;
  \node[above] at (-0.2,2.2) {Afloramiento};
  \node at (1.3,0.6) {``Ojo de agua''};
  \node[above] at (6.5,2.3) {Material permeable};
\end{tikzpicture}
\caption{Vertiente de dique.}
\end{subfigure}
\caption{Tipos de vertientes.}
\end{figure}

\section{Ley de Darcy (1856)}

Darcy observó experimentalmente, en un tubo lleno de arena, que el caudal
$Q$ que atraviesa la muestra cumple:
\[
Q\propto A,\qquad Q\propto \Delta h,\qquad Q\propto\frac{1}{L}
\]
de donde:
\begin{formula}
$\displaystyle Q=-K\cdot A\cdot\left(\frac{\dd h}{\dd l}\right)=-K\cdot A\cdot i$\\[4pt]
\textbf{Ley de Darcy} --- válida para flujo laminar ($Re<4$)
\end{formula}

\begin{figure}[H]
\centering
\begin{minipage}{0.5\textwidth}\centering
\begin{tikzpicture}[scale=0.85,font=\small]
  \draw (-0.5,0) -- (5.5,0) node[right]{$z=0$};
  % tubo inclinado
  \begin{scope}
    \fill[agua!60] (0.9,4.2) -- (1.3,4.6) -- (4.1,1.8) -- (3.7,1.4) -- cycle;
    \draw (0.9,4.2) -- (3.7,1.4) (1.3,4.6) -- (4.1,1.8);
    \foreach \i in {1,...,60}{\pgfmathsetmacro{\t}{rnd}\pgfmathsetmacro{\s}{rnd}
      \fill[azulusm] ({0.9+2.8*\t+0.4*\s*0.7},{4.2-2.8*\t+0.4*\s*0.7}) circle (0.6pt);}
  \end{scope}
  % piezómetros
  \fill[agua] (1.05,4.35) rectangle (1.35,5.6);
  \draw (1.05,4.35) -- (1.05,6.0) (1.35,4.5) -- (1.35,6.0);
  \fill[agua] (3.85,1.65) rectangle (4.15,5.2);
  \draw (3.85,1.8) -- (3.85,6.0) (4.15,1.9) -- (4.15,6.0);
  \draw[dashed] (0,5.6) -- (5,5.6) (3.6,5.2) -- (5,5.2);
  \draw[cota] (0.2,0) -- (0.2,5.6) node[midway,left]{$h_1$};
  \draw[cota] (0.6,0) -- (0.6,4.4) node[midway,left]{$z_1$};
  \draw[cota] (4.9,0) -- (4.9,5.2) node[midway,right]{$h_2$};
  \draw[cota] (4.5,0) -- (4.5,1.8) node[midway,right]{$z_2$};
  \draw[cota] (5.3,5.2) -- (5.3,5.6) node[midway,right]{$\Delta h$};
  \draw[cota] (0.6,3.6) -- (3.3,0.9) node[midway,below left]{$L$};
  \draw[flecha,azulusm] (1.4,3.9) -- (2.2,3.1);
\end{tikzpicture}
\end{minipage}\hfill
\begin{minipage}{0.45\textwidth}\centering
\begin{tikzpicture}
\begin{axis}[width=\linewidth,height=5.5cm,xmin=0,xmax=15,ymin=0,ymax=200,
  xlabel={Gradiente hidráulico, $-\dd h/\dd l$},ylabel={Descarga específica $v$},font=\scriptsize]
  \addplot[only marks,mark=o] coordinates {(2.5,45)(2.5,42)(5,85)(5,90)(6,100)(7.5,125)(9,155)(11,185)(11.5,190)(12,195)};
  \addplot[only marks,mark=*] coordinates {(2,27)(4,42)(7,64)(9,88)(9.5,92)(11,108)};
  \addplot[domain=0:12,samples=2] {16.5*x};
  \addplot[domain=0:15,samples=2] {9.6*x};
  \draw (axis cs:10,96) -- (axis cs:12.5,96) -- (axis cs:12.5,120);
  \node[font=\tiny] at (axis cs:11.5,85) {pendiente $=K$};
\end{axis}
\end{tikzpicture}
\end{minipage}
\caption{Experimento de Darcy: la descarga específica es proporcional al
gradiente hidráulico, siendo la pendiente la conductividad hidráulica $K$.}
\end{figure}

\subsection{Generalización de la Ley de Darcy}

Generalmente, en un suelo estratificado la resistencia al flujo no es igual
en todas direcciones $\Rightarrow$ \textbf{medio anisotrópico}:
\begin{formula}
$\vec v=-\mathbf{K}\,\nabla h$
\end{formula}
con
\[
\vec v=v_x\hat\imath+v_y\hat\jmath+v_z\hat k,\qquad
\nabla h=\frac{\partial h}{\partial x}\hat\imath+\frac{\partial h}{\partial y}\hat\jmath+\frac{\partial h}{\partial z}\hat k,\qquad
\mathbf K=\begin{pmatrix}K_{xx}&K_{xy}&K_{xz}\\K_{yx}&K_{yy}&K_{yz}\\K_{zx}&K_{zy}&K_{zz}\end{pmatrix}
\]
\begin{itemize}
  \item Si la dirección principal de anisotropía coincide con $x,y,z$:
  \[K_{xy}=K_{yx}=K_{xz}=K_{zx}=K_{yz}=K_{zy}=0\]
  \item Si no hay dirección de flujo preferente (\textbf{medio isotrópico}):
  \[K_x=K_y=K_z=K\]
\end{itemize}

\section{Conductividad hidráulica, $K$}
\begin{itemize}
  \item También conocida como \textbf{coeficiente de permeabilidad}.
  \item Unidades L/T, comúnmente expresada en cm/s.
  \item Función de las propiedades del medio poroso y del fluido.
\end{itemize}
\begin{formula}
$\displaystyle K=k\,\frac{\rho g}{\mu}=k\,\frac{\gamma}{\mu}$
\qquad $k$: permeabilidad intrínseca [darcy],\quad
1 darcy $=9.87\times10^{-9}$ cm$^2$
\end{formula}

\begin{table}[H]
\centering
\begin{tabular}{lcc}
\toprule
\textbf{Material} & \textbf{Permeabilidad intrínseca $k$ [darcy]} & \textbf{Conductividad $K$ [cm/s]}\\
\midrule
Arcilla & $10^{-6}$ -- $10^{-3}$ & $10^{-9}$ -- $10^{-6}$\\
Limo, limos arenosos, arenas arcillosas & $10^{-3}$ -- $10^{-1}$ & $10^{-6}$ -- $10^{-4}$\\
Arenas limosas, arenas finas & $10^{-2}$ -- $1$ & $10^{-5}$ -- $10^{-3}$\\
Arenas bien distribuidas & $1$ -- $10^{2}$ & $10^{-3}$ -- $10^{-1}$\\
Gravas bien distribuidas & $10$ -- $10^{3}$ & $10^{-2}$ -- $1$\\
\bottomrule
\end{tabular}
\caption{Valores típicos de permeabilidad intrínseca y conductividad hidráulica.}
\end{table}

\subsection{Factores de los suelos que condicionan $K$}
\begin{itemize}
  \item \textbf{Porosidad efectiva.}
  \item \textbf{Forma de los granos:} esférica, tubular o prismática, laminar
        ($K$ aumenta desde laminar hacia esférica).
  \item \textbf{Granulometría:} I.~material homogéneo; II.~material heterogéneo.
        Se cumple $K_I>K_{II}$.
  \item \textbf{Presencia de material fino} (pasante bajo malla 200,
        0.074~mm): si se presenta un 12\% a 15\% de material bajo esta malla,
        el material presenta una mala permeabilidad.
  \item Viscosidad del agua.
  \item Grado de saturación del suelo.
  \item En el caso de las arcillas también influyen:
  \begin{itemize}
    \item La estructura.
    \item La concentración iónica.
    \item Espesor de capas de agua adheridas a las partículas de arcilla.
  \end{itemize}
  \item Para un material arenoso o gravoso, que sea mal graduado indica que
        posee una mayor cantidad de limos o arcillas.
  \item Las arcillas no se consideran bien o mal graduadas.
\end{itemize}

\begin{figure}[H]
\centering
\begin{tikzpicture}[font=\small]
  \draw[flecha] (0,0) -- (6,0) node[right]{$D$};
  \draw[flecha] (0,0) -- (0,3.5) node[left]{\%};
  \draw[azulusm] (0,3) -- (5.5,3);
  \draw[thick] (0,0) .. controls (1.5,0.3) and (2,2.9) .. (3,3) -- (5,3);
  \draw[thick,dashed] (0,0) .. controls (2.5,0.2) and (3.5,2) .. (4.6,3) -- (5,3);
  \node at (1.8,1.8) {I};
  \node at (3.3,1.5) {II};
  \node[rojo] at (0.9,2.4) {Homogéneo};
  \node[rojo] at (5,1.2) {Heterogéneo};
  \node at (2.5,-0.5) {$K_I>K_{II}$};
  \node[azulusm] at (3,3.35) {Curvas granulométricas};
\end{tikzpicture}
\caption{Curvas granulométricas: material homogéneo (I) y heterogéneo (II).}
\end{figure}

\section{Métodos indirectos para determinar $K$}

\subsection{Información bibliográfica o empírica}
Depende del tipo de suelo; existen tablas de rangos de permeabilidad
intrínseca y conductividad hidráulica por tipo de material (p.~ej.\ Freeze
\& Cherry, 1979): desde rocas ígneas no fracturadas ($K\sim10^{-11}$ cm/s)
hasta gravas y calizas kársticas ($K\sim 10^{0}$--$10^{2}$ cm/s).

\subsection{Método de Hazen}
En suelos arenosos, a partir de la curva granulométrica, para arenas con
$d_{10}$ entre 0.1 y 0.3~mm:
\begin{formula}
$K=C\cdot(d_{10})^2\ \ [\text{cm/s}],\qquad d_{10}$ en [cm]
\end{formula}
\begin{table}[H]
\centering
\begin{tabular}{lc}
\toprule
\textbf{Tipo de arena} & $C$\\
\midrule
Arena muy fina, mal distribuida & 40 -- 80\\
Arena fina con una gran cantidad de material fino & 40 -- 80\\
Arena media, bien distribuida & 80 -- 120\\
Arena gruesa, mal distribuida & 80 -- 120\\
Arena gruesa, bien distribuida, limpia & 120 -- 150\\
\bottomrule
\end{tabular}
\end{table}

\subsection{Método de Slichter}
Basado en curva granulométrica y porosidad:
\[
K=\frac{\gamma}{\mu}\,\frac{10.0219\cdot d_{10}^2}{K_1},\qquad d_{10}\ \text{en [cm]}
\]
\begin{table}[H]
\centering
\begin{tabular}{cc@{\hspace{1cm}}cc}
\toprule
$n$ & $1/K_1$ & $n$ & $1/K_1$\\
\midrule
0.26 & 0.00187 & 0.38 & 0.04154\\
0.28 & 0.01517 & 0.40 & 0.04922\\
0.30 & 0.01905 & 0.42 & 0.05789\\
0.32 & 0.02356 & 0.44 & 0.06776\\
0.34 & 0.02878 & 0.46 & 0.07838\\
0.36 & 0.03473 & & \\
\bottomrule
\end{tabular}
\end{table}

\subsection{Método de Kozeny--Carman}
\[
K=\left(\frac{\rho g}{\mu}\right)\cdot\left[\frac{n^3}{(1-n^2)}\right]\cdot\left(\frac{d_{50}^2}{180}\right)
\]
con $\mu$: viscosidad; $n$: porosidad; $d_{50}$ en [cm].

\section{Métodos directos para determinar $K$}

\subsection{Pruebas de laboratorio}
La conductividad hidráulica de un suelo saturado se mide mediante
\textbf{permeámetros}.

\paragraph{Permeámetro de carga constante}
\begin{itemize}
  \item Sedimentos no cohesivos, arenas, rocas.
  \item Se alcanza una situación de equilibrio.
  \item Utilizar gradientes hidráulicos similares a los existentes en terreno.
  \item $h_1-h_2<50\%\,L$.
\end{itemize}
\begin{formula}
$\displaystyle K=\frac{4\cdot L\cdot Q}{\pi\cdot(h_1-h_2)\cdot D^2}$
\end{formula}

\paragraph{Permeámetro de carga variable}
\begin{itemize}
  \item Sedimentos cohesivos con baja permeabilidad.
\end{itemize}
\begin{formula}
$\displaystyle K=\frac{d_t^2\cdot L}{d_c^2\cdot t}\cdot\ln\!\left(\frac{h_0}{h}\right)$
\qquad $d_t$: diámetro del tubo;\ $d_c$: diámetro de la muestra
\end{formula}

\begin{figure}[H]
\centering
\begin{subfigure}{0.48\textwidth}\centering
\begin{tikzpicture}[scale=0.8,font=\small]
  % estanque superior
  \draw (2,5) rectangle (3,6.2); \fill[agua] (2,5) rectangle (3,5.9);
  \draw[thick] (2.3,6.2) -- (2.3,6.6) -- (3.3,6.6) node[right,rojo]{Entrada};
  \draw[thick] (3,5.9) -- (3.6,5.9) -- (3.6,5.6) node[below,rojo]{Descarga};
  % tubo de conexión
  \draw[line width=3pt,agua] (2.9,5) -- (2.9,0.2) -- (2.2,0.2);
  % muestra
  \draw (1,0) rectangle (2.2,2.5); \fill[agua] (1,0) rectangle (2.2,2.2);
  \fill[pattern=north east lines] (1,0.4) rectangle (2.2,0.55);
  \fill[pattern=north east lines] (1,1.75) rectangle (2.2,1.9);
  \foreach \i in {1,...,30}\fill[azulusm!70!black] ({1.05+1.1*rnd},{0.6+1.1*rnd}) circle (0.7pt);
  \node[font=\scriptsize,align=center] at (1.6,1.15) {Muestra};
  \draw[thick] (1,2.2) -- (0.4,2.2) -- (0.4,1.9);
  \fill[agua] (-0.2,0) rectangle (0.8,0.8); \draw (-0.2,0) rectangle (0.8,1.1);
  \draw[dashed] (-0.6,5.9) -- (2,5.9) node[pos=0,left]{$h_1$};
  \draw[dashed] (-0.6,2.2) -- (0.4,2.2) node[pos=0,left]{$h_2$};
  \draw[cota] (2.5,2.2) -- (2.5,5.9) node[midway,left]{$h_1-h_2$};
  \draw[cota] (0.8,0.55) -- (0.8,1.75) node[midway,left]{$L$};
  \draw[cota] (1,-0.3) -- (2.2,-0.3) node[midway,below]{$D$};
  \node[rojo,align=center] at (4.3,1.7) {Placas\\porosas};
\end{tikzpicture}
\caption{Carga constante.}
\end{subfigure}\hfill
\begin{subfigure}{0.48\textwidth}\centering
\begin{tikzpicture}[scale=0.8,font=\small]
  \draw (1,0) rectangle (2.5,3);
  \fill[agua] (1,0) rectangle (2.5,2.7);
  \fill[pattern=north east lines] (1,0.4) rectangle (2.5,0.55);
  \fill[pattern=north east lines] (1,1.9) rectangle (2.5,2.05);
  \foreach \i in {1,...,30}\fill[azulusm!70!black] ({1.05+1.4*rnd},{0.6+1.25*rnd}) circle (0.7pt);
  \node[font=\scriptsize,align=center] at (1.75,1.2) {Muestra};
  % tubo delgado
  \draw (3.2,0.2) -- (3.2,5.5) (3.4,0.2) -- (3.4,5.5);
  \fill[agua] (3.2,0.2) rectangle (3.4,3.5);
  \draw[line width=2pt,agua] (2.5,0.15) -- (3.3,0.15);
  \draw[thick] (1,2.7) -- (0.2,2.7) -- (0.2,2.4);
  \fill[agua] (-0.3,0) rectangle (0.7,0.9); \draw (-0.3,0) rectangle (0.7,1.2);
  \draw[dashed] (1.7,5) -- (3.2,5) (1.7,2.7) -- (1,2.7) (2.5,3.5) -- (3.2,3.5);
  \draw[cota] (1.7,2.7) -- (1.7,5) node[midway,left]{$h_0$};
  \draw[cota] (2.8,2.7) -- (2.8,3.5) node[midway,left]{$h$};
  \draw[cota] (0.8,0.55) -- (0.8,1.9) node[midway,left]{$L$};
  \draw[cota] (1,-0.3) -- (2.5,-0.3) node[midway,below]{$d_c$};
  \draw[->] (2.8,5.8) -- (3.2,5.8); \draw[<-] (3.4,5.8) -- (3.8,5.8) node[right]{$d_t$};
  \node[rojo,align=center] at (4.6,1.9) {Placas\\porosas};
\end{tikzpicture}
\caption{Carga variable.}
\end{subfigure}
\caption{Permeámetros de laboratorio.}
\end{figure}

\subsection{Mediciones en terreno}
Para obtener ya sea la conductividad hidráulica o la tasa de infiltración.
\begin{itemize}
  \item \textbf{Pruebas de infiltración:} tasas de infiltración y
        conductividad hidráulica saturada (infiltrómetro de doble anillo,
        método Porchet).
\end{itemize}
\paragraph{Método Porchet}
\[
f=\frac{R}{2\cdot(t_2-t_1)}\cdot\ln\!\left(\frac{2h_1+R}{2h_2+R}\right)
\]

\paragraph{Uso de dos pozos} (asume velocidad horizontal):
\begin{formula}
$\displaystyle K\approx\frac{L^2}{\bar t\cdot\Delta h}$
\end{formula}
$L$: distancia entre pozos de observación; $\bar t$: tiempo promedio de
observación; $\Delta h$: desnivel de la carga hidráulica.

\paragraph{Mediante un pozo} (trazador):
\begin{formula}
$\displaystyle \ln\!\left(\frac{C}{C_0}\right)=\frac{4\cdot v}{\pi\cdot D}\cdot t,
\qquad v=K\cdot i$
\end{formula}
$C$: concentración del trazador en el instante $t$; $C_0$: concentración
inicial del trazador; $v$: velocidad de inyección; $D$: diámetro del pozo;
$i$: tasa de infiltración $\Delta h/\Delta t$.

\paragraph{Pruebas de bombeo:} se verán más adelante.

\section{Estratificación (heterogeneidad)}
Cuando las características hidráulicas varían espacialmente, el acuífero es
\textbf{heterogéneo}; caso contrario, es \textbf{homogéneo}.

\subsection{Flujo paralelo a la estratificación}
El caudal total es la suma de los caudales de cada estrato
($q_T=\sum q_i$), y todos los estratos tienen el mismo gradiente:
\begin{formula}
$\displaystyle K_{eq}=\frac{\sum_{i=1}^{N}K_i\cdot m_i}{\sum_{i=1}^{N}m_i}$
\end{formula}
Un ejemplo de flujo paralelo es el flujo de agua que atraviesa el suelo bajo
una presa: dos estratos $(m_1,K_1)$ y $(m_2,K_2)$ pueden reemplazarse por un
único estrato de espesor $m_1+m_2$ y conductividad $K_{eq}$.

\begin{figure}[H]
\centering
\begin{tikzpicture}[font=\small]
  \fill[agua] (0,0) rectangle (1.2,3.2);
  \fill[agua] (6.2,0) rectangle (7.4,2.6);
  \fill[agua!60] (1.2,0) rectangle (6.2,2.0);
  \fill[pattern=crosshatch,pattern color=gray!50] (1.2,2.0) rectangle (6.2,3.6);
  \draw (1.2,0) rectangle (6.2,3.6);
  \foreach \y/\k in {0/1,0.4/2,0.8/3,1.4/4}{
    \draw (1.2,\y) -- (6.2,\y);
    \node at (1.5,\y+0.2) {$K_\k$};
    \draw[flecha,azulusm] (3.3,\y+0.2) -- (4.3,\y+0.2) node[above,font=\scriptsize,black]{$Q_\k$};}
  \draw[dashed,azulusm] (1.2,3.2) -- (6.2,2.6);
  \draw[cota] (0.3,0) -- (0.3,3.2) node[midway,left]{$h_0$};
  \draw[cota] (7.1,0) -- (7.1,2.6) node[midway,right]{$h_L$};
  \fill[gray!40] (0,-0.2) rectangle (7.4,0);
  \node[below] at (1.2,-0.2) {$0$}; \node[below] at (6.2,-0.2) {$L$};
\end{tikzpicture}
\caption{Flujo paralelo a la estratificación.}
\end{figure}

\subsection{Flujo perpendicular a la estratificación}
El caudal (por unidad de área) es el mismo en todos los estratos y la
pérdida de carga total es la suma de las pérdidas en cada estrato
($\Delta h_T=\sum\Delta h_i$):
\begin{formula}
$\displaystyle K_{eq}=\frac{\sum_{i=1}^{N}m_i}{\sum_{i=1}^{N}\dfrac{m_i}{K_i}}$
\end{formula}

\begin{figure}[H]
\centering
\begin{tikzpicture}[font=\small]
  \fill[agua!70] (0,0) rectangle (7,3.2);
  \foreach \y/\k/\d in {0/1/0.8,0.8/2/0.6,1.4/3/1.3,2.7/4/0.5}{
    \draw (0,\y) -- (7,\y);
    \node[left] at (0,\y+\d/2) {$k_\k$};
    \draw[cota] (6.4,\y) -- (6.4,\y+\d) node[midway,right]{$d_\k$};
    \draw[flecha,azulusm,decorate,decoration={snake,amplitude=1pt,segment length=4pt,post length=2mm}] (1.5,\y+0.1) -- (1.5,\y+\d-0.1);}
  \draw (0,3.2) -- (7,3.2);
  % piezómetros
  \foreach \x/\b/\t in {3/0/4.6,3.25/0.8/4.4,3.5/1.4/4.05,3.75/2.7/3.8}{
    \draw (\x,\b) rectangle (\x+0.22,5);
    \fill[azulusm!80] (\x,\b) rectangle (\x+0.22,\t);}
  \node[right,font=\scriptsize] at (4.1,4.5) {$\Delta h_1,\dots,\Delta h_4$};
\end{tikzpicture}
\caption{Flujo perpendicular a la estratificación.}
\end{figure}

\section{Propiedades de los acuíferos}

\paragraph{Transmisibilidad}
\begin{itemize}
  \item Medida de la facilidad con que un acuífero transmite agua.
  \item Unidades de m$^3$/s/m.
  \item Asume que el flujo es horizontal.
\end{itemize}
\begin{formula}
$T=K\times m$\qquad $m$: espesor saturado o potencia
\end{formula}

\paragraph{Coeficiente de almacenamiento}
\begin{itemize}
  \item Volumen de agua, por unidad de área y cambio de altura, que la unidad
        permeable absorbe o libera.
\end{itemize}
\begin{formula}
$V_{\text{drenado}}=-S\cdot A\cdot\Delta h$\qquad $S$: coeficiente de almacenamiento
\end{formula}

% =====================================================================
%  CLASE 02 -- Hidráulica de pozos en régimen permanente
% =====================================================================
\clase{02}{Aguas subterráneas: hidráulica de pozos en régimen permanente}

\section{Recordando\ldots}
\begin{itemize}
  \item \textbf{Acuífero libre:} cercano a la superficie terrestre; el agua
        subterránea tiene superficie abierta a la atmósfera.
  \item \textbf{Acuífero artesiano o confinado:} confinado entre estratos
        impermeables, por ende, bajo presión.
  \item \textbf{Ley de Darcy (1856)}, válida para flujo laminar ($Re<4$):
  \[
  Q=-K\cdot A\cdot\left(\frac{\dd h}{\dd l}\right)=-K\cdot A\cdot i
  \]
\end{itemize}

\section{Propiedades de los acuíferos}

\subsection{Coeficiente de almacenamiento específico $S_s$ [1/L]}
Es la cantidad de agua, por unidad de volumen, que es almacenada o liberada
debido a la compresibilidad del esqueleto mineral y del agua en los poros,
debido a un cambio unitario en el nivel de agua en el acuífero [1/L].
\begin{formula}
$S_s=g\rho(\alpha+\eta\beta)$
\end{formula}
donde $g$: aceleración de gravedad [m/s$^2$]; $\rho$: densidad del agua
[kg/m$^3$]; $\eta$: porosidad del medio; $\alpha$: compresibilidad de la
matriz sólida del acuífero [m$^2$/N]; $\beta$: compresibilidad del agua [m$^2$/N].

\subsection{Coeficiente de almacenamiento $S$}
Volumen de agua, por unidad de área y cambio de altura, que la unidad
permeable absorbe o libera.

\paragraph{Acuífero confinado de espesor $m$:}
\begin{formula}
$S=S_s\,m\qquad\qquad V_{\text{drenado}}=-S\cdot A\cdot\Delta h$\\[2pt]
Acuífero confinado: $10^{-4}<S<0.005$
\end{formula}

\begin{figure}[H]
\centering
\begin{tikzpicture}[font=\small]
  \fill[gray!30] (0,0) rectangle (7,0.8);  \node at (1,0.4) {Acuitardo};
  \fill[cyan!40] (0,0.8) rectangle (7,2.4); \node at (1,1.6) {Acuífero};
  \fill[gray!30] (0,2.4) rectangle (7,3.3); \node at (1,2.85) {Acuitardo};
  \draw[thick] (0,3.6) -- (7,3.4);
  \draw[thick,dashed] (0,3.3) -- (7,3.1);
  \draw[decorate,decoration={brace,amplitude=3pt}] (7.1,3.4) -- (7.1,3.1)
    node[midway,right,align=left,font=\scriptsize]{Descenso unitario de la\\superficie piezométrica};
  % prisma unitario
  \draw[dashed] (4,0.8) rectangle (4.8,2.4);
  \draw[dashed] (4,2.4) -- (4.3,2.65) -- (5.1,2.65) -- (4.8,2.4);
  \draw[dashed] (5.1,2.65) -- (5.1,1.05) -- (4.8,0.8);
  \node[font=\scriptsize] at (4.4,2.85) {1};\node[font=\scriptsize] at (4.95,2.85) {1};
  \draw[cota] (5.6,0.8) -- (5.6,2.4) node[midway,right]{$m$};
\end{tikzpicture}
\caption{Coeficiente de almacenamiento en un acuífero confinado.}
\end{figure}

\paragraph{Acuífero libre:}
\begin{itemize}
  \item En acuíferos libres también se conoce como \textbf{rendimiento
        específico} ($S=S_y$).
  \item Coincide con la porosidad efectiva ($S_y=n_e$).
\end{itemize}
\begin{formula}
$V_{\text{drenado}}=-S\cdot A\cdot\Delta h$\qquad Acuífero libre: $0.05<S<0.2$
\end{formula}

\begin{table}[H]
\centering
\begin{tabular}{lc}
\toprule
\textbf{Material} & $S_y$\\
\midrule
Arcilla & 0.00 -- 0.05\\
Limo & 0.03 -- 0.19\\
Arena fina & 0.10 -- 0.32\\
Arena media & 0.15 -- 0.32\\
Arena gruesa & 0.20 -- 0.35\\
Grava & 0.14 -- 0.35\\
\bottomrule
\end{tabular}
\caption{Valores típicos de rendimiento específico (Johnson, 1967).}
\end{table}

\begin{figure}[H]
\centering
\begin{tikzpicture}[font=\scriptsize]
  \foreach \X/\tit/\lab in {0/{Acuífero confinado}/S, 5.5/{Acuífero libre}/m_e}{
  \begin{scope}[xshift=\X cm]
    \draw (0,0) rectangle (1.6,5);
    \draw (1.6,0) -- (2.1,0.4) -- (2.1,5.4) -- (0.5,5.4) -- (0,5);
    \draw (1.6,5) -- (2.1,5.4);
    \node[rotate=15] at (1.05,5.2) {1 m$^2$};
    \fill[pattern=dots,pattern color=gray] (0,0) rectangle (1.6,2.6);
    \draw[dashed] (0,3.6) -- (1.6,3.6) (0,2.6) -- (1.6,2.6);
    \draw[cota] (1.8,2.6) -- (1.8,3.6) node[midway,right]{1 m};
    \draw (-0.8,4.8) rectangle (-0.3,5.4); \node at (-0.55,5.1) {$\lab$};
    \node[align=center] at (0.8,-0.5) {\tit};
  \end{scope}}
  \node[align=left] at (-0.3,6.2) {Extrayendo un volumen $S$ desciende\\la superficie piezométrica 1 metro};
  \node[align=left] at (5.4,6.2) {Extrayendo un volumen $m_e$ desciende\\la superficie freática 1 metro};
\end{tikzpicture}
\caption{Interpretación física del coeficiente de almacenamiento en acuíferos
confinado y libre.}
\end{figure}

\section{Ecuación de conservación de masa en un medio saturado}

Para un acuífero confinado, el balance de masa sobre un elemento diferencial,
asumiendo fluido incompresible y flujo sólo en la dirección $x$, queda:
\[
\rho\left(v_x-\left(v_x+\frac{\partial v_x}{\partial x}\,\dd x\right)\right)\dd y\,\dd z
=\frac{\partial M}{\partial t}
\quad\Longrightarrow\quad
\boxed{-\rho\,\frac{\partial v_x}{\partial x}\,\dd x\,\dd y\,\dd z=\frac{\partial M}{\partial t}}
\]

\begin{figure}[H]
\centering
\begin{tikzpicture}[font=\small]
  \fill[gray!25] (0,0) rectangle (2.5,2.5);
  \fill[gray!45] (2.5,0) -- (3.3,0.6) -- (3.3,3.1) -- (2.5,2.5) -- cycle;
  \fill[gray!15] (0,2.5) -- (0.8,3.1) -- (3.3,3.1) -- (2.5,2.5) -- cycle;
  \draw[flecha,cyan] (-1.2,1.25) -- (-0.1,1.25) node[pos=0,above right]{$v_x$};
  \draw[flecha,cyan] (2.9,1.6) -- (4,1.6) node[right,black]{$v_x+\dfrac{\partial v_x}{\partial x}\dd x$};
  \node[below] at (0,0) {$x$}; \node[below] at (2.5,0) {$x+\Delta x$};
  \node[below] at (1.25,0) {$\Delta x$};
  \node[right] at (3.3,0.3) {$\Delta y$}; \node[right] at (3.3,2.9) {$\Delta z$};
\end{tikzpicture}
\caption{Elemento diferencial para el balance de masa.}
\end{figure}

Considerando que
\[
M=\rho V_{\text{drenado}}=-\rho S_s\,\dd z\,A\,\Delta h,\qquad A=\dd x\,\dd y
\quad\Longrightarrow\quad
\frac{\partial M}{\partial t}=-\rho S_s\frac{\partial h}{\partial t}\,\dd x\,\dd y\,\dd z
\]
Si se asume flujo en las tres direcciones:
\begin{formula}
$\displaystyle -\left(\frac{\partial v_x}{\partial x}+\frac{\partial v_y}{\partial y}+\frac{\partial v_z}{\partial z}\right)=-S_s\frac{\partial h}{\partial t}$
\end{formula}
\emph{¿Y de dónde obtenemos las velocidades?} Recordando que Darcy (1856)
estableció que las velocidades de escurrimiento en medios porosos están
dadas por:
\[
v_x=-K_x\frac{\partial h}{\partial x},\qquad
v_y=-K_y\frac{\partial h}{\partial y},\qquad
v_z=-K_z\frac{\partial h}{\partial z}
\]
De este modo se deduce la \textbf{ecuación general de escurrimiento en un
medio saturado confinado}:
\begin{formula}
$\displaystyle
\frac{\partial}{\partial x}\left(K_x\frac{\partial h}{\partial x}\right)+
\frac{\partial}{\partial y}\left(K_y\frac{\partial h}{\partial y}\right)+
\frac{\partial}{\partial z}\left(K_z\frac{\partial h}{\partial z}\right)=S_s\frac{\partial h}{\partial t}$\\[4pt]
$S_s$: almacenamiento específico, $S/m$
\end{formula}

\section{Ecuaciones generales}
\subsection{Casos particulares}
\begin{itemize}
  \item Flujo transiente en medio homogéneo e isotrópico:
  \[\frac{\partial^2h}{\partial x^2}+\frac{\partial^2h}{\partial y^2}+\frac{\partial^2h}{\partial z^2}=\frac{S_s}{K}\frac{\partial h}{\partial t}\]
  \item Flujo permanente, anisotrópico, homogéneo:
  \[\frac{\partial}{\partial x}\left(K_x\frac{\partial h}{\partial x}\right)+\frac{\partial}{\partial y}\left(K_y\frac{\partial h}{\partial y}\right)+\frac{\partial}{\partial z}\left(K_z\frac{\partial h}{\partial z}\right)=0\]
  \item Flujo permanente, isotrópico, homogéneo (\textbf{ecuación de Laplace}):
\end{itemize}
\begin{formula}
$\displaystyle\frac{\partial^2h}{\partial x^2}+\frac{\partial^2h}{\partial y^2}+\frac{\partial^2h}{\partial z^2}=0$\qquad\textbf{Laplace}
\end{formula}

\subsection{Ecuación de Laplace en coordenadas cilíndricas}
Flujo permanente, isotrópico y homogéneo:
\begin{formula}
$\displaystyle\frac{1}{r}\frac{\partial h}{\partial r}+\frac{\partial^2h}{\partial r^2}+\frac{1}{r^2}\frac{\partial^2h}{\partial\theta^2}+\frac{\partial^2h}{\partial z^2}=0$
\end{formula}

\subsection{Métodos de solución}
\begin{enumerate}
  \item Solución directa exacta.
  \item Integración analítica simplificada.
  \item Redes de flujo.
  \item Modelos analógicos.
  \item Modelos numéricos.
\end{enumerate}

\paragraph{¿Para qué hacemos esto?} ¿Por qué es importante saber la dinámica
de la altura del nivel freático en el tiempo? Porque al bombear un pozo se
genera un \emph{cono de depresión} que determina cuánto desciende la napa en
los alrededores, y con ello el caudal que es posible extraer.

\section{Hidráulica de pozos}
\subsection{Supuestos básicos}
\begin{itemize}
  \item Los acuíferos se considerarán de propiedades homogéneas e isotrópicas,
        de extensión infinita ($S$, $m$ constantes).
  \item El tipo de acuífero se supone invariante en el espacio y el tiempo,
        i.e., es libre, confinado o semiconfinado.
  \item Se supone que la superficie piezométrica es horizontal antes del
        inicio de operación del pozo.
  \item El flujo será radial concéntrico hacia el pozo, y tendrá un radio de
        influencia por sobre el cual se supondrá despreciable el descenso de la napa.
  \item No existen pérdidas de energía por roce al escurrir agua hacia el pozo.
  \item No se considera (por ahora) superposición de pozos.
\end{itemize}

% ---------------------------------------------------------------------
%  Macro para dibujar un pozo en acuífero (confinado=1 / libre=0)
% ---------------------------------------------------------------------
\newcommand{\dibujopozo}[1]{%
\begin{tikzpicture}[font=\small,scale=0.95]
  % base impermeable
  \draw[thick] (-4,0) -- (4,0);
  \foreach \x in {-3.9,-3.6,...,3.9} \draw (\x,0) -- (\x-0.2,-0.25);
  % pozo
  \draw[thick] (-0.35,0.1) -- (-0.35,4.6) (0.35,0.1) -- (0.35,4.6);
  \draw[thick] (-4,4.6) -- (-0.35,4.6) (0.35,4.6) -- (4,4.6);
  \fill[azulusm] (-0.1,4.5) rectangle (0.1,5.1);
  \fill[azulusm] (-0.25,5.1) -- (0.25,5.1) -- (0,5.45) -- cycle;
  \node at (0.5,5.35) {$Q$};
  % superficie original y cono
  \draw[dashed,cyan!70!blue,thick] (-4,3.9) -- (4,3.9);
  \node[azulusm,font=\scriptsize\bfseries,anchor=south west] at (-4.2,3.95) {Superficie piezométrica original};
  \draw[cyan!70!blue,thick] (-4,3.85) .. controls (-2,3.7) and (-1,3.2) .. (-0.35,2.4) -- (0.35,2.4)
      .. controls (1,3.2) and (2,3.7) .. (4,3.85);
  \pic at (0,2.42) {nivelagua};
  \ifnum#1=1
    % techo impermeable del acuífero confinado
    \draw[thick] (-4,2.0) -- (-0.35,2.0) (0.35,2.0) -- (4,2.0);
    \foreach \x in {-3.9,-3.6,...,-0.5,0.6,0.9,...,3.9} \draw (\x,2.0) -- (\x+0.15,2.25);
    \draw[cota] (-3.5,0) -- (-3.5,2.0) node[midway,left]{$m$};
    \node[azulusm,font=\scriptsize\bfseries,align=center] at (-2.6,3.1) {Superficie\\piezométrica};
    \def\ytop{1.9}
  \else
    \node[azulusm,font=\scriptsize,align=center] at (-2.9,2.9) {$\dfrac{\partial h}{\partial s}\approx\dfrac{\partial h}{\partial r}$\\Dupuit--Forchheimer};
    \def\ytop{2.2}
  \fi
  % flechas de flujo
  \foreach \y in {0.3,0.6,...,\ytop}{
    \draw[->,azulusm] (-1.2,\y) -- (-0.5,\y);
    \draw[->,azulusm] (1.2,\y) -- (0.5,\y);}
  % cotas
  \draw[cota] (-0.15,0) -- (-0.15,2.4) node[pos=0.8,left]{$h_0$};
  \draw[->] (0,1.2) -- (0.35,1.2) node[right,font=\scriptsize]{$r_0$};
  \draw[cota] (2.2,0) -- (2.2,3.62) node[midway,right]{$h$};
  \draw[cota] (3.8,0) -- (3.8,3.9) node[midway,left]{$H$};
  \draw[->] (0,-0.5) -- (2.2,-0.5) node[midway,below]{$r$};
  \draw[->] (0,-1.1) -- (3.9,-1.1) node[midway,below,align=center]{$R$\\[-2pt]\scriptsize\color{azulusm}\bfseries Radio de influencia};
\end{tikzpicture}}

\subsection{Pozo de penetración total en acuífero confinado}
\begin{itemize}
  \item Extensión lateral infinita.
  \item Pozo bombea caudal constante.
  \item Líneas de flujo radiales y convergentes hacia el pozo.
  \item Equipotenciales son círculos concéntricos.
\end{itemize}

\begin{figure}[H]
\centering
\dibujopozo{1}
\caption{Pozo de penetración total en acuífero confinado.}
\end{figure}

\begin{nota}[title={Deducción (ecuación de Thiem)}]
Por continuidad, a través de un cilindro de radio $r$ y altura $m$:
\[
Q=2\pi r\,m\,K\frac{\dd h}{\dd r}\;\Rightarrow\;\int_h^H\dd h=\frac{Q}{2\pi Km}\int_r^R\frac{\dd r}{r}
\]
\end{nota}
\begin{formula}
$\displaystyle\Delta h=H-h=\frac{Q}{2\pi T}\ln\!\left(\frac{R}{r}\right)$\\[4pt]
$R$: radio de influencia;\quad $T=Km$: transmisibilidad
\end{formula}

\subsection{Pozo de penetración total en acuífero libre}
\begin{itemize}
  \item Pozo bombea caudal constante.
  \item Líneas de flujo radiales y convergentes hacia el pozo.
  \item Equipotenciales son círculos concéntricos.
  \item Flujo horizontal.
  \item Gradiente hidráulico no varía con la profundidad.
  \item Se cumple ley hidrostática.
\end{itemize}
Estas son las hipótesis de \textbf{Dupuit--Forchheimer}:
$\dfrac{\partial h}{\partial s}\approx\dfrac{\partial h}{\partial r}$.

\begin{figure}[H]
\centering
\dibujopozo{0}
\caption{Pozo de penetración total en acuífero libre.}
\end{figure}

\begin{nota}[title={Deducción (Dupuit)}]
Ahora el área de flujo es $2\pi r h$:
\[
Q=2\pi r\,h\,K\frac{\dd h}{\dd r}\;\Rightarrow\;\int_h^H2h\,\dd h=\frac{Q}{\pi K}\int_r^R\frac{\dd r}{r}
\]
\end{nota}
\begin{formula}
$\displaystyle H^2-h^2=\frac{Q}{\pi K}\ln\!\left(\frac{R}{r}\right)$\qquad $R$: radio de influencia
\end{formula}

\paragraph{Aproximaciones.} Si el espesor inicial de la napa saturada es muy
grande en comparación con las depresiones cerca del pozo, la ecuación se
simplifica y es similar a la de un acuífero confinado:
\begin{itemize}
  \item Si $\Delta h<0.1\,H$:
  \[\Delta h=\frac{Q}{2\pi KH}\ln\frac{R}{r}\]
  \item Si $\Delta h<0.5\,H$ (\textbf{corrección de Jacob}):
  \[\Delta h'=\Delta h-\frac{(\Delta h)^2}{2H}=\frac{Q}{2\pi KH}\ln\frac{R}{r}\]
\end{itemize}

\section{Acuíferos libres de gran potencia}

\begin{figure}[H]
\centering
\begin{subfigure}{0.48\textwidth}\centering
\begin{tikzpicture}[scale=0.8,font=\small]
  \draw[thick] (-2.5,4) -- (-0.4,4) (0.4,4) -- (2.5,4);
  \foreach \x in {-2.4,-2.0,...,-0.6,0.6,1.0,...,2.4} \draw (\x,4) -- (\x+0.25,4.25);
  \draw[thick] (-0.4,4.6) -- (-0.4,1.3) (0.4,4.6) -- (0.4,1.3);
  \draw[cyan,thick] (-2.5,3.2) -- (2.5,3.2); \pic at (-2.1,3.22) {nivelagua};
  \shade[ball color=azulusm!70] (0,1) circle (0.8);
  \foreach \a in {0,45,...,315} \draw[->,azulusm,thick] (\a:2.2) ++(0,1) -- ++(\a+180:1.1);
  \node at (1.6,4.5) {$H>R$};
\end{tikzpicture}
\caption{Pozo en acuífero libre de gran potencia\\(escurrimiento esférico).}
\end{subfigure}\hfill
\begin{subfigure}{0.48\textwidth}\centering
\begin{tikzpicture}[scale=0.8,font=\small]
  \draw[thick] (-2.5,4) -- (-0.4,4) (0.4,4) -- (2.5,4);
  \foreach \x in {-2.4,-2.0,...,-0.6,0.6,1.0,...,2.4} \draw (\x,4) -- (\x+0.25,4.25);
  \draw[thick] (-0.4,4.6) -- (-0.4,2.6) (0.4,4.6) -- (0.4,2.6);
  \draw[cyan,thick] (-2.5,2.8) -- (2.5,2.8); \pic at (-2.1,2.82) {nivelagua};
  \fill[azulusm!70] (-0.7,2.8) arc (180:360:0.7) -- cycle;
  \foreach \a in {0,-45,-90,-135,-180} \draw[->,azulusm,thick] (\a:2.0) ++(0,2.8) -- ++(\a+180:1.1);
\end{tikzpicture}
\caption{Pozo somero en acuífero libre de gran potencia (escurrimiento semi-esférico).}
\end{subfigure}
\end{figure}

\begin{formula}
Escurrimiento esférico: $\displaystyle\Delta h=H-h=\frac{Q}{4\pi K}\left(\frac1r-\frac1R\right)$\\[6pt]
Escurrimiento semi-esférico: $\displaystyle\Delta h=H-h=\frac{Q}{2\pi K}\left(\frac1r-\frac1R\right)$,
\quad $r_0\approx0.6\sim8$ [m]
\end{formula}

\section{Pozo de penetración parcial en acuífero confinado}
Sea $b$ la longitud de penetración del pozo en un acuífero confinado de
espesor $m$.

\begin{figure}[H]
\centering
\begin{tikzpicture}[font=\small,scale=0.9]
  \draw[thick] (-4,0) -- (4,0);
  \foreach \x in {-3.9,-3.6,...,3.9} \draw (\x,0) -- (\x-0.2,-0.25);
  \draw[thick] (-4,3) -- (-0.35,3) (0.35,3) -- (4,3);
  \foreach \x in {-3.9,-3.6,...,-0.5,0.6,0.9,...,3.9} \draw (\x,3) -- (\x+0.15,3.25);
  \draw[thick] (-0.35,1.2) -- (-0.35,4.6) (0.35,1.2) -- (0.35,4.6);
  \draw[thick] (-4,4.6) -- (-0.35,4.6) (0.35,4.6) -- (4,4.6);
  \fill[azulusm] (-0.1,4.5) rectangle (0.1,5.1);
  \fill[azulusm] (-0.25,5.1) -- (0.25,5.1) -- (0,5.45) -- cycle; \node at (0.5,5.35) {$Q$};
  \draw[dashed,cyan!70!blue,thick] (-4,4.2) .. controls (-2,4.0) and (-1,3.8) .. (-0.35,3.6) -- (0.35,3.6) .. controls (1,3.8) and (2,4.0) .. (4,4.2);
  \foreach \y in {1.5,1.8,...,2.8}{\draw[->,azulusm] (-1.2,\y) -- (-0.5,\y);\draw[->,azulusm] (1.2,\y) -- (0.5,\y);}
  \foreach \a in {-60,-90,-120}\draw[->,azulusm] (\a:0.9) ++(0,1.2) -- ++(\a+180:0.5);
  \draw[cota] (-3.5,0) -- (-3.5,3) node[midway,left]{$m$};
  \draw[cota] (1.6,1.2) -- (1.6,3) node[midway,right]{$b$};
  \draw[->] (0,-0.5) -- (2.5,-0.5) node[midway,below]{$r$};
\end{tikzpicture}
\caption{Pozo de penetración parcial en acuífero confinado.}
\end{figure}

\paragraph{Solución de Nasberg} (válida para $b/m\geq0.3$):
\[
h_1-h_2=\frac{Q}{2\pi Km}\left(\ln\!\left(\frac{r_1}{r_2}\right)+\operatorname{senh}^{-1}\!\left(\frac{m}{r_1}\right)-\operatorname{senh}^{-1}\!\left(\frac{m}{r_2}\right)-\frac{m}{b}\left(\operatorname{senh}^{-1}\!\left(\frac{b}{r_1}\right)-\operatorname{senh}^{-1}\!\left(\frac{b}{r_2}\right)\right)\right)
\]

\paragraph{Solución de Girinsky} (válida para $b/m<0.3$):
\[
h_1-h_2=\frac{Q}{2\pi K}\left(\frac1b\ln\frac{1.6\,b}{r_2}-\operatorname{arcsenh}\frac{b}{r_1}\right)
\]

\section{Estimación del radio de influencia $R$}
\paragraph{En terreno}
\begin{itemize}
  \item Mediciones de depresiones en función de la distancia al pozo.
  \item Determinar $R$ cuando la depresión se hace cero (gráfico $\Delta h$ vs.\ $\ln r$).
\end{itemize}
\begin{figure}[H]
\centering
\begin{tikzpicture}[font=\small]
  \draw[flecha] (0,0) -- (5,0) node[below]{$\ln r$};
  \draw[flecha] (0,0) -- (0,3.2) node[left]{$\Delta h$};
  \draw[dashed] (0,2.6) node[left]{$h_0$} -- (0.8,2.6) -- (0.8,0) node[below]{$r_0$};
  \draw[rojo,thick] (0.8,2.6) -- (3.2,0) node[below,black]{$R$};
  \draw[azulusm,thick] (0.8,2.6) .. controls (1.5,1.6) and (2.8,0.3) .. (3.3,0.05) -- (4.5,0.02);
\end{tikzpicture}
\caption{Estimación de $R$ a partir de mediciones de depresión.}
\end{figure}

\paragraph{Fórmula de Sichardt (1928):} depende de $K$ y, por ende, de la granulometría.
\begin{formula}
$R=3000\,\Delta h_0\sqrt{K}$ \ (m),\qquad $K$ en [m/s],\ $\Delta h_0$ en [m]
\end{formula}
\begin{table}[H]
\centering
\begin{tabular}{lcc}
\toprule
\textbf{Material} & $D$ (mm) & $R$ (m)\\
\midrule
Gravas & mayor a 10 & 1500\\
Gravillas & 2 -- 10 & 500 -- 1500\\
Gravilla-arcilla & 1 -- 2 & 400 -- 500\\
Arena gruesa & 0.5 -- 1 & 200 -- 400\\
Arena media & 0.25 -- 0.5 & 100 -- 200\\
Arena fina & 0.1 -- 0.25 & 50 -- 100\\
Arena muy fina & 0.05 -- 0.1 & 10 -- 50\\
\bottomrule
\end{tabular}
\caption{Valores referenciales del radio de influencia según granulometría.}
\end{table}

\paragraph{Otros métodos empíricos} (ver Louwyck et al., 2022, \emph{The
Radius of Influence Myth}, Water 14(2):149,
\url{https://www.mdpi.com/2073-4441/14/2/149}):
\begin{table}[H]
\centering\small
\begin{tabular}{llll}
\toprule
\textbf{Modelo} & \textbf{Régimen} & \textbf{Borde superior} & \textbf{Radio de influencia $R$}\\
\midrule
Dupuit & Permanente & Nivel freático & Borde exterior (dato)\\
Thiem & Permanente & Impermeable & Borde exterior (dato)\\
de Glee & Permanente & Filtrante & $R=4\sqrt{cKD}$\\
Theis & Transiente & Impermeable & $R=1.5\sqrt{KDt/S}$\\
Hantush--Jacob & Transiente & Filtrante & $R=1.5\sqrt{KDt/S}$ si $t<0.01Sc$;\ $R=4\sqrt{cKD}$ si $t>10Sc$\\
Ernst & Permanente & Drenaje + recarga & $R=\sqrt{Q/(\pi N)}$ o $R=4\sqrt{cKD}$\\
\bottomrule
\end{tabular}
\end{table}
con $KD$: transmisibilidad; $c$: resistencia; $S$: coeficiente de
almacenamiento; $N$: flujo de infiltración; $Q$: caudal de bombeo; $t$: tiempo.

% =====================================================================
%  CLASE 03 -- Interferencia de pozos, método de las imágenes,
%              pozos en acuíferos de varias capas
% =====================================================================
\clase{03}{Interferencia de pozos, método de las imágenes y pozos en acuíferos de varias capas}

\section{Resumen de hidráulica de pozos}

\begin{table}[H]
\centering
\renewcommand{\arraystretch}{2.2}
\begin{tabular}{l>{$\displaystyle}c<{$}}
\toprule
\textbf{Caso} & \textbf{Ecuación}\\
\midrule
Penetración total, acuífero confinado & \Delta h=H-h=\frac{Q}{2\pi T}\ln\left(\frac Rr\right)\\
Penetración total, acuífero libre & H^2-h^2=\frac{Q}{\pi K}\ln\left(\frac Rr\right)\\
Acuífero libre de gran potencia & \Delta h=H-h=\frac{Q}{4\pi K}\left(\frac1r-\frac1R\right)\\
Pozo somero en acuífero libre de gran potencia & \Delta h=H-h=\frac{Q}{2\pi K}\left(\frac1r-\frac1R\right)\\
Penetración parcial, acuífero confinado & h_1-h_2=\frac{Q}{2\pi K}\left(\frac1b\ln\frac{1.6b}{r_2}-\operatorname{arcsenh}\frac{b}{r_1}\right)\\
\bottomrule
\end{tabular}
\end{table}

\begin{center}
\colorbox{azulusm}{\parbox{0.8\textwidth}{\centering\color{white}\large\bfseries
¿Y si hay más de un pozo extrayendo agua del acuífero?}}
\end{center}

\section{Depresión producto de múltiples pozos}

\textbf{Solución: superposición.} Como la ecuación de Laplace es lineal, la
depresión total en un punto es la suma de las depresiones producidas por
cada pozo.

\begin{figure}[H]
\centering
\begin{tikzpicture}[font=\small]
  \coordinate (p) at (0,0);
  \foreach \pt/\lab/\pos in {(-2.3,1.3)/r_1/above, (-2.0,-1.3)/r_2/above, (2.3,1.6)/r_3/above, (2.2,-1.5)/r_4/above}{
    \draw[dashed,azulusm,thick] (p) -- \pt node[midway,\pos]{$\lab$};
    \fill[azulusm] \pt circle (0.18);}
  \fill[orange] (p) circle (0.08) node[above=2pt]{$p$};
\end{tikzpicture}
\caption{Punto $p$ afectado por varios pozos de bombeo.}
\end{figure}

\paragraph{Pozos de penetración total en acuíferos confinados:}
\begin{formula}
$\displaystyle\Delta h_T=\frac{1}{2\pi T}\sum Q_i\ln\left(\frac{R_i}{r_i}\right)$
\end{formula}
\begin{itemize}
  \item El punto $p$ sólo se ve afectado por los pozos ubicados a una distancia
        menor al radio de influencia ($R_i$).
  \item ¡Para acuíferos libres, linealizar primero!
\end{itemize}

\paragraph{Caso particular:} pozos que bombean el mismo caudal y tienen el
mismo radio de influencia:
\begin{formula}
$\displaystyle\Delta h_T=\frac{Q}{2\pi T}\ln\left(\frac{R^n}{\prod_{i=1}^n r_i}\right)$
\end{formula}

\begin{center}
\colorbox{azulusm}{\parbox{0.8\textwidth}{\centering\color{white}\large\bfseries
¿Y si el acuífero no posee extensión lateral infinita?}}
\end{center}
Por ejemplo, un pozo cerca de un arroyo (zona ribereña) o cerca de un
afloramiento de roca impermeable.

\section{Método de las imágenes}

Superposición de soluciones para obtener la respuesta a condiciones de borde
específicas:
\begin{itemize}
  \item Recargas lineales (masas de agua con nivel constante).
  \item Barreras impermeables.
\end{itemize}

\subsection{Recarga lineal}
\begin{itemize}
  \item Fuente que recarga al acuífero (río, lago, mar).
  \item Se debe lograr \textbf{desnivel nulo en el borde}.
  \item Se reemplaza la fuente por un \textbf{pozo imagen} ubicado
        simétricamente respecto al pozo real, que \textbf{inyecta} caudal $Q$.
\end{itemize}

% Macro esquema pozo real + pozo imagen
\newcommand{\esquemaimagen}[3]{% #1 signo imagen, #2 etiqueta borde, #3 textos lados
\begin{tikzpicture}[font=\small]
  \coordinate (R) at (0,0); \coordinate (I) at (6,0); \coordinate (P) at (2,2);
  \draw[dashed,azulusm,thick] (3,-0.8) -- (3,2.6);
  \draw[orange!90!black,{Stealth}-{Stealth},thick] (R) -- (3,0) node[midway,above]{$b$};
  \draw[orange!90!black,{Stealth}-{Stealth},thick] (3,0) -- (I) node[midway,above]{$b$};
  \draw[cyan!80!blue,thick,-{Stealth}] (P) -- (R) node[pos=0.4,left]{$r_r$};
  \draw[cyan!80!blue,thick,-{Stealth}] (P) -- (I) node[pos=0.55,above]{$r_i$};
  \fill[azulusm] (R) circle (0.18) (I) circle (0.18);
  \fill[orange] (P) circle (0.08) node[above]{$p$};
  \node[below=4pt,align=center] at (R) {Pozo\\real};
  \node[below=4pt,align=center] at (I) {Pozo\\imaginario};
  \node[left=6pt] at (R) {$+Q,\Delta h_r$};
  \node[right=6pt] at (I) {$#1Q,\Delta h_i$};
  \node[azulusm,align=center,below] at (3,-0.8) {#2};
  #3
\end{tikzpicture}}

\begin{figure}[H]
\centering
\esquemaimagen{-}{Línea de conservación\\del nivel de carga $h$}{%
  \node at (1.5,-0.4) {\textit{Tierra}}; \node at (4.5,-0.4) {\textit{Agua}};}
\caption{Método de las imágenes para una recarga lineal: el pozo imagen inyecta caudal.}
\end{figure}

Depresiones producidas por los pozos:
\[
\Delta h_r=\frac{Q}{2\pi T}\ln\left(\frac{R}{r_r}\right),\qquad
\Delta h_i=\frac{Q}{2\pi T}\ln\left(\frac{R}{r_i}\right)
\]
La depresión producida en el punto $p$ es $\Delta h=\Delta h_r-\Delta h_i$:
\begin{formula}
$\Delta h=\begin{cases}
\dfrac{Q}{2\pi T}\ln\left(\dfrac{r_i}{r_r}\right) & r_r<R\ \wedge\ r_i<R\\[10pt]
\dfrac{Q}{2\pi T}\ln\left(\dfrac{R}{r_r}\right) & r_r<R\ \wedge\ r_i\geq R\\[10pt]
0 & r_r\geq R\ \wedge\ r_i\geq R
\end{cases}$
\end{formula}

\subsection{Barrera impermeable}
\begin{itemize}
  \item Barrera que impide el flujo.
  \item Se debe lograr \textbf{velocidad nula en el borde}.
  \item Se reemplaza la barrera por un pozo imagen ubicado simétricamente
        respecto al pozo real, que \textbf{extrae} caudal $Q$.
\end{itemize}

\begin{figure}[H]
\centering
\esquemaimagen{+}{Borde impermeable}{}
\caption{Método de las imágenes para una barrera impermeable: el pozo imagen extrae caudal.}
\end{figure}

La depresión producida en el punto $p$ es $\Delta h=\Delta h_r+\Delta h_i$:
\begin{formula}
$\Delta h=\begin{cases}
\dfrac{Q}{2\pi T}\ln\left(\dfrac{R^2}{r_i\cdot r_r}\right) & r_r<R\ \wedge\ r_i<R\\[10pt]
\dfrac{Q}{2\pi T}\ln\left(\dfrac{R}{r_r}\right) & r_r<R\ \wedge\ r_i\geq R\\[10pt]
0 & r_r\geq R\ \wedge\ r_i\geq R
\end{cases}$
\end{formula}

\section{Método de las imágenes: situaciones particulares}

\subsection{Más de un pozo cerca de una fuente de recarga o barrera impermeable}
Cada pozo real tiene su propio pozo imagen, simétrico respecto al borde:
con signo opuesto ($-Q_j$) frente a una recarga y con igual signo ($+Q_j$)
frente a una barrera.

\begin{figure}[H]
\centering
\begin{tikzpicture}[font=\small,scale=0.75]
  \foreach \X/\s/\lab in {0/-/{Línea de conservación\\del nivel de carga $h$}, 9/+/{Borde impermeable}}{
  \begin{scope}[xshift=\X cm]
    \draw[dashed,azulusm,thick] (2.5,-0.6) -- (2.5,3.2);
    \foreach \y/\n/\b in {0/1/b_1, 2/2/b_2}{
      \fill[azulusm] (0.5,\y) circle (0.15) (4.5,\y) circle (0.15);
      \draw[orange!90!black,{Stealth}-{Stealth}] (0.5,\y) -- (2.5,\y) node[midway,above]{$\b$};
      \draw[orange!90!black,{Stealth}-{Stealth}] (2.5,\y) -- (4.5,\y) node[midway,above]{$\b$};
      \node[above left,font=\scriptsize] at (0.5,\y+0.1) {$+Q_\n,\Delta h_{r\n}$};
      \node[above right,font=\scriptsize] at (4.5,\y+0.1) {$\s Q_\n,\Delta h_{i\n}$};}
    \node[azulusm,align=center,below] at (2.5,-0.6) {\lab};
  \end{scope}}
\end{tikzpicture}
\caption{Varios pozos cerca de una recarga (izquierda) y de una barrera impermeable (derecha).}
\end{figure}

\subsection{Borde curvo o esquina de una fuente de recarga}
\begin{itemize}
  \item Se crea el pozo imaginario 1 ($P_{i1}$) para anular la depresión en la línea 1.
  \item Se crea el pozo imaginario 2 ($P_{i2}$) para anular la depresión en la línea 2.
  \item Se crea un tercer pozo imaginario ($P_{i3}$) con un caudal $+Q$ que
        anule la influencia de $P_{i1}$ sobre la línea 2 y de $P_{i2}$ sobre la línea 1.
\end{itemize}

\begin{figure}[H]
\centering
\begin{tikzpicture}[font=\small]
  \fill[aguaclara] (-3,-2.5) rectangle (3,2.5);
  \fill[suelo!80] (0,-2.5) rectangle (3,0) ;
  \fill[suelo!80] (0.4,0) -- (3,0) -- (3,0) ;
  \draw[dashed,azulusm,very thick] (0,-2.5) -- (0,0) -- (3,0);
  \node[azulusm] at (2.3,0.3) {Línea 1}; \node[azulusm] at (0,-2.8) {Línea 2};
  \node at (-1.5,-2.1) {Agua}; \node at (1.5,-2.1) {Tierra};
  \coordinate (Pr) at (1.8,-1.2); \coordinate (Pi1) at (1.8,1.2);
  \coordinate (Pi2) at (-1.8,-1.2); \coordinate (Pi3) at (-1.8,1.2);
  \foreach \c in {Pr,Pi1,Pi2,Pi3} \fill[azulusm] (\c) circle (0.13);
  \node[below right] at (Pr) {Pozo real};
  \node[right] at (Pi1) {$P_{i1}$: $-Q,\Delta h_{i1}$};
  \node[left] at (Pi2) {$P_{i2}$: $-Q,\Delta h_{i2}$};
  \node[left] at (Pi3) {$P_{i3}$: $+Q,\Delta h_{i3}$};
  \draw[yellow!70!orange,{Stealth}-{Stealth}] (Pi2) -- (0,-1.2) node[midway,below]{$b$};
  \draw[yellow!70!orange,{Stealth}-{Stealth}] (0,-1.2) -- (Pr) node[midway,below]{$b$};
  \draw[red,{Stealth}-{Stealth}] (1.8,0) -- (Pr) node[midway,right]{$a$};
  \draw[red,{Stealth}-{Stealth}] (1.8,0) -- (Pi1) node[midway,right]{$a$};
  \draw[yellow!70!orange,{Stealth}-{Stealth}] (Pi3) -- (Pi1) node[midway,above]{$2b$};
  \draw[red,{Stealth}-{Stealth}] (Pi3) -- (Pi2) node[midway,left]{$2a$};
\end{tikzpicture}
\caption{Pozo cerca de la esquina de una fuente de recarga: tres pozos imagen.}
\end{figure}

\subsection{Presencia de una fuente de recarga y un borde impermeable}

\begin{figure}[H]
\centering
\begin{tikzpicture}[font=\small]
  \fill[aguaclara] (-3,0) rectangle (0,2);
  \fill[suelo!70] (0,0) rectangle (3.5,2);
  \fill[roca!60] (3.5,0) rectangle (6.5,2);
  \draw[dashed,azulusm,very thick] (0,-0.3) -- (0,2.3);
  \draw[dashed,azulusm,very thick] (3.5,-0.3) -- (3.5,2.3);
  \coordinate (Pr) at (1.2,1); \coordinate (Pi2) at (-1.2,1); \coordinate (Pi1) at (5.8,1);
  \foreach \c in {Pr,Pi1,Pi2} \fill[azulusm] (\c) circle (0.15);
  \draw[yellow!70!orange,{Stealth}-{Stealth}] (Pi2) -- (0,1) node[midway,below]{$a$};
  \draw[yellow!70!orange,{Stealth}-{Stealth}] (0,1) -- (Pr) node[midway,below]{$a$};
  \draw[orange!90!black,{Stealth}-{Stealth}] (Pr) -- (3.5,1) node[midway,below]{$b$};
  \draw[orange!90!black,{Stealth}-{Stealth}] (3.5,1) -- (Pi1) node[midway,below]{$b$};
  \node[above] at (Pr) {$+Q,\Delta h_{r1}$};
  \node[above left] at (Pi2) {$P_{i2}$: $-Q,\Delta h_{r2}$};
  \node[above right] at (Pi1) {$P_{i1}$: $Q,\Delta h_{i1}$};
  \node at (-1.5,0.3) {Agua}; \node at (1.75,0.3) {Tierra}; \node at (5,0.3) {Roca};
  \node[below,align=center] at (0,-0.3) {Línea de conservación\\del nivel de carga $h$};
  \node[below] at (3.9,-0.3) {Borde impermeable};
\end{tikzpicture}
\caption{Pozo entre una fuente de recarga y un borde impermeable.}
\end{figure}

Puede ocurrir que:
\[
2b+a<R\quad\vee\quad 2a+b<R
\]
\textcolor{rojo}{¡Un pozo imaginario está afectando la otra condición de borde!}
En ese caso deben agregarse nuevas imágenes (de las imágenes) hasta que
todas las imágenes adicionales queden fuera del radio de influencia.

\section{Acuíferos estratificados}

\subsection{Pozo que capta de varias napas (napas artesianas)}
\begin{itemize}
  \item El caudal total proviene de ambas napas.
  \item Cada napa aporta según la diferencia entre su nivel piezométrico y la
        depresión en el pozo.
\end{itemize}

% Macro para dibujar el pozo en dos napas
\newcommand{\pozodosnapas}{%
  \fill[suelo!80] (0,3) rectangle (6,4.6);
  \fill[pattern=north east lines,pattern color=brown] (0,2.7) rectangle (6,3);
  \fill[aguaclara] (0,1.5) rectangle (6,2.7);
  \fill[pattern=north east lines,pattern color=brown] (0,1.2) rectangle (6,1.5);
  \fill[agua!70] (0,0) rectangle (6,1.2);
  \fill[pattern=north east lines,pattern color=brown] (0,-0.3) rectangle (6,0);
  \fill[white] (2.6,0) rectangle (3.4,4.6);
  \fill[agua!50] (2.6,0) rectangle (3.4,3.3);
  \draw[thick] (2.6,0) -- (2.6,4.6) (3.4,0) -- (3.4,4.6);
  \draw[thick] (0,4.6) -- (2.6,4.6) (3.4,4.6) -- (6,4.6);
  \node at (1,2.1) {$m_1,\,K_1$}; \node at (1,0.6) {$m_2,\,K_2$};
  \node at (3.1,0.25) {$r_0$};
  \fill[azulusm] (2.9,4.6) rectangle (3.1,5.1);
  \fill[azulusm] (2.75,5.1) -- (3.25,5.1) -- (3,5.4) -- cycle; \node at (3.4,5.3) {$Q$};}

\begin{figure}[H]
\centering
\begin{tikzpicture}[font=\small,scale=1.25]
  \pozodosnapas
  \draw[dashed,azulusm,thick] (0,4.3) -- (6,4.3) node[pos=0,above right,font=\scriptsize]{Nivel piezométrico 2};
  \draw[dashed,azulusm,thick] (0,3.8) -- (6,3.8) node[pos=0,above right,font=\scriptsize]{Nivel piezométrico 1};
  \draw[cota] (4.2,3.3) -- (4.2,3.8) node[midway,left]{$\Delta h_1$};
  \draw[cota] (5.3,3.3) -- (5.3,4.3) node[midway,left]{$\Delta h_2$};
  \draw[dashed,cyan] (3.4,3.3) -- (6,3.3);
\end{tikzpicture}
\caption{Pozo que capta de dos napas artesianas.}
\end{figure}

Si el valor de $R$ es constante para las distintas napas:
\begin{formula}
$\displaystyle Q=\frac{2\pi}{\ln\left(\frac{R}{r}\right)}\cdot\sum_{i=1}^n\Delta h_i\cdot K_i\cdot m_i$
\end{formula}
\begin{itemize}
  \item Si $\Delta h_1$ llega al acuífero confinado 1, éste funciona como un acuífero libre.
  \item En el caso de una sola napa estratificada, se tiene un único valor $\Delta h$.
\end{itemize}

\subsection{Considerando el instante inicial}
\begin{itemize}
  \item No existe extracción de caudal del pozo.
  \item El nivel estático (N.E.) se encuentra entre los dos niveles piezométricos.
\end{itemize}
\begin{center}
\colorbox{azulusm}{\parbox{0.5\textwidth}{\centering\color{white}\bfseries Una napa alimenta a la otra}}
\end{center}

\begin{figure}[H]
\centering
\begin{tikzpicture}[font=\small,scale=1.25]
  \pozodosnapas
  \draw[dashed,azulusm,thick] (0,4.3) -- (6,4.3) node[pos=0,above right,font=\scriptsize]{Nivel piezométrico 1} node[right]{$t=0$};
  \draw[dashed,azulusm,thick] (0,3.6) -- (6,3.6) node[pos=0,below right,font=\scriptsize]{Nivel piezométrico 2};
  \draw[dashed,cyan] (3.4,4.0) -- (6,4.0) node[right,black]{$t>0$};
  \draw[cota] (4.2,4.0) -- (4.2,4.3) node[midway,right]{$\Delta h_1$};
  \draw[cota] (5.3,3.6) -- (5.3,4.0) node[midway,right]{$\Delta h_2$};
  \draw[cota] (1.3,3.6) -- (1.3,4.3) node[midway,right]{$\Delta h_T$};
  \draw[flecha,azulusm] (1.8,2.4) -- (2.55,2.4); \draw[flecha,azulusm] (4.8,2.4) -- (3.45,2.4) node[pos=0,right]{$Q$};
  \draw[flecha,azulusm] (2.55,0.8) -- (1.8,0.8); \draw[flecha,azulusm] (3.45,0.8) -- (4.8,0.8) node[right]{$Q$};
\end{tikzpicture}
\caption{Condición inicial: la napa 1 alimenta a la napa 2 a través del pozo.}
\end{figure}

\[
\Delta h_1=\frac{Q}{2\pi T_1}\ln\left(\frac{R}{r_0}\right),\qquad
\Delta h_2=\frac{-Q}{2\pi T_2}\ln\left(\frac{R}{r_0}\right)
\quad\Longrightarrow\quad
\frac{\Delta h_T}{\Delta h_1}=\frac{T_1}{T_2}+1
\]
\begin{formula}
$\displaystyle\Delta h_T=\Delta h_1\left(\frac{T_1+T_2}{T_2}\right)$
\end{formula}
Por lo tanto, midiendo $\Delta h_1$ en la perforación se puede calcular
$\Delta h_T$ y $\Delta h_2$.

\subsection{Nivel estático en el estrato de arcilla o roca del segundo acuífero}
Si el nivel estático (N.E.) se encuentra justo en el estrato de arcilla o
roca del segundo acuífero:
\begin{itemize}
  \item El acuífero 1 se convierte en un acuífero libre, aportando un caudal
        constante $Q_0$.
\end{itemize}
\begin{center}
\colorbox{azulusm}{\parbox{0.7\textwidth}{\centering\color{white}\bfseries Cuando comienza el bombeo las napas están en equilibrio ($h_1=h_2$)}}
\end{center}
\begin{formula}
$\displaystyle Q=Q_0+\frac{2\pi K_2\cdot m_2\,\Delta h_2}{\ln\left(\frac{R}{r_0}\right)}$
\end{formula}

\begin{figure}[H]
\centering
\begin{tikzpicture}[font=\small]
  \draw[flecha] (0,0) -- (5,0) node[below]{$Q$};
  \draw[flecha] (0,0) -- (0,4) node[left]{$\Delta h$};
  \draw[orange!90!black,very thick] (0,0) -- (2.2,1.3) -- (3.1,3.8);
  \draw[azulusm,very thick,dashed] (1.7,0) node[below,black]{$Q_0$} -- (2.2,1.3);
  \node[above left] at (2.2,1.3) {$Q_1(t)+Q_2(t)$};
  \node[right] at (2.7,2.8) {$Q_0+Q_2(t)$};
\end{tikzpicture}
\caption{Curva de agotamiento ($\Delta h$ vs.\ $Q$).}
\end{figure}

Si se posee una curva de agotamiento ($\Delta h$ vs.\ $Q$) se pueden calcular
$K_1$ y $K_2$.

% =====================================================================
%  CLASE 04 -- Otros métodos de solución a la ecuación de Laplace
% =====================================================================
\clase{04}{Otros métodos de solución a la ecuación de Laplace}

\section{Ecuaciones generales}
\begin{formula}
$\displaystyle\frac{\partial^2h}{\partial x^2}+\frac{\partial^2h}{\partial y^2}+\frac{\partial^2h}{\partial z^2}=0$
\qquad Flujo permanente, isotrópico, homogéneo (\textbf{Laplace})
\end{formula}
Ecuación de Laplace en coordenadas cilíndricas:
\[
\frac1r\frac{\partial h}{\partial r}+\frac{\partial^2h}{\partial r^2}+\frac{1}{r^2}\frac{\partial^2h}{\partial\theta^2}+\frac{\partial^2h}{\partial z^2}=0
\]
\paragraph{Métodos de solución}
\begin{enumerate}
  \item Solución directa exacta (clases previas).
  \item Integración analítica simplificada (clases previas).
  \item \textbf{Redes de flujo.}
  \item \textbf{Modelos analógicos.}
  \item \textbf{Modelos numéricos.}
\end{enumerate}

\section{Redes de flujo}

Se desea resolver la ecuación (flujo permanente, medio anisotrópico y homogéneo):
\[
\frac{\partial}{\partial x}\left(K_x\frac{\partial h}{\partial x}\right)+\frac{\partial}{\partial y}\left(K_y\frac{\partial h}{\partial y}\right)+\frac{\partial}{\partial z}\left(K_z\frac{\partial h}{\partial z}\right)=0
\]
Si el medio es isotrópico, $K_x=K_y=K_z=K$:
\[
K\frac{\partial^2h}{\partial x^2}+K\frac{\partial^2h}{\partial y^2}+K\frac{\partial^2h}{\partial z^2}=K\nabla^2h=0
\quad\Longrightarrow\quad\boxed{\nabla^2h=0}\ \ \text{(Laplace)}
\]
Luego, si consideramos la \textbf{función potencial} $\phi$:
\begin{formula}
$\phi=Kh\quad\Longrightarrow\quad\nabla^2\phi=0$
\end{formula}
Las superficies $\phi(x,y,z)=\text{cte}$ son \textbf{superficies equipotenciales}, y
\[
\phi_1-\phi_2=K(h_1-h_2)=\text{cte}\quad\Rightarrow\quad\Delta h=\text{cte}\ \ \text{(caída de potencial)}
\]

\begin{figure}[H]
\centering
\begin{tikzpicture}[font=\small]
  \draw[flecha] (0,0) -- (0,3) node[left]{$z$};
  \draw[flecha] (0,0) -- (3,-0.8) node[right]{$y$};
  \draw[flecha] (0,0) -- (-1.2,-1.5) node[left]{$x$};
  \foreach \y/\n in {0.4/2,1.6/1}{
    \fill[gray!40,draw=azulusm] (0.2,\y) .. controls (0.6,\y+0.7) and (1.4,\y+0.3) .. (2.2,\y+0.9)
      .. controls (1.8,\y+0.2) and (1.2,\y-0.2) .. (0.2,\y) -- cycle;
    \node at (1.1,\y+0.4) {$\phi_\n$};}
  \draw[flecha,azulusm] (2.6,1.8) -- (3.3,1.8);
  \node[right,align=center] at (3.3,1.8) {$\phi(x,y,z)=\text{cte}$\\\color{azulusm}Superficies equipotenciales};
\end{tikzpicture}
\caption{Superficies equipotenciales.}
\end{figure}

\subsection{Flujo bidimensional}
En el caso de flujo bidimensional ($v_y=0$):
\[
\frac{\partial^2h}{\partial x^2}+\frac{\partial^2h}{\partial z^2}=0
\]
Esta ecuación representa dos familias ortogonales de curvas que satisfacen
la ecuación de Laplace:
\begin{formula}
$\psi(x,z)$: líneas de flujo\qquad\qquad $\phi(x,z)$: líneas equipotenciales
\end{formula}

\paragraph{Líneas equipotenciales}
\[
\frac{\partial^2\phi}{\partial x^2}+\frac{\partial^2\phi}{\partial z^2}=0
\]
Como $\phi(x,z)=\text{cte}$:
\[
\dd\phi=\frac{\partial\phi}{\partial x}\dd x+\frac{\partial\phi}{\partial z}\dd z=0
\]
Por otro lado, $\phi=Kh$ (tomando la convención de signo $v=K\,\partial h/\partial x$):
\[
\frac{\partial\phi}{\partial x}=K\frac{\partial h}{\partial x}=v_x,\qquad
\frac{\partial\phi}{\partial z}=K\frac{\partial h}{\partial z}=v_z
\quad\Longrightarrow\quad
\frac{\dd z}{\dd x}=-\frac{v_x}{v_z}=m_\phi
\]

\paragraph{Líneas de flujo} (línea a lo largo de la cual viaja una partícula de agua):
\[
\left(\frac{\dd z}{\dd x}\right)_{L.f.}=\frac{v_z}{v_x}\quad\Longrightarrow\quad v_z\,\dd x-v_x\,\dd z=0
\]
Se define la \textbf{función de flujo} $\psi(x,z)$:
\[
\frac{\partial\psi}{\partial x}=v_z,\qquad\frac{\partial\psi}{\partial z}=-v_x
\]
Así,
\[
\dd\psi=\frac{\partial\psi}{\partial x}\dd x+\frac{\partial\psi}{\partial z}\dd z=0\;\Rightarrow\;\psi=\text{cte},
\qquad m_\psi=\frac{v_z}{v_x}\;\Rightarrow\;\boxed{m_\psi\cdot m_\phi=-1}
\]
¡Funciones ortogonales!

\begin{figure}[H]
\centering
\begin{subfigure}{0.45\textwidth}\centering
\begin{tikzpicture}[font=\small]
  \draw[flecha] (0,0) -- (4,0) node[right]{$x$};
  \draw[flecha] (0,0) -- (0,3) node[left]{$z$};
  \draw[azulusm,thick] (0.2,0.8) .. controls (0.8,2.2) and (2,2.7) .. (3.6,2.9) node[pos=0.4,above left]{$\phi_1$};
  \draw[orange!90!black,thick] (0.6,0.2) .. controls (1.1,1.5) and (2.2,2.0) .. (3.6,2.2) node[pos=0.45,above left]{$\phi_2$};
  \node[align=center,font=\scriptsize] at (2.9,1.2) {$\phi(x,z)=\text{cte}$\\\color{azulusm}Líneas equipotenciales};
\end{tikzpicture}
\end{subfigure}
\begin{subfigure}{0.45\textwidth}\centering
\begin{tikzpicture}[font=\small]
  \draw[flecha] (0,0) -- (4,0) node[right]{$x$};
  \draw[flecha] (0,0) -- (0,3) node[left]{$z$};
  \draw[orange!90!black,thick] (0.3,0.2) .. controls (0.8,1.6) and (1.6,2.3) .. (3.6,2.6) node[pos=0.1,left]{$\psi$};
  \draw[flecha] (2.0,2.2) -- (3.3,2.9) node[above]{$\vec v$};
  \draw[dashed] (2.0,2.2) -- (3.3,2.2) node[midway,below]{$v_x$} -- (3.3,2.9) node[midway,right]{$v_z$};
\end{tikzpicture}
\end{subfigure}
\caption{Líneas equipotenciales y línea de flujo con su vector velocidad tangente.}
\end{figure}

\paragraph{Condiciones de Cauchy--Riemann.} Por lo tanto,
\begin{formula}
$\displaystyle v_x=\frac{\partial\phi}{\partial x}=-\frac{\partial\psi}{\partial z},\qquad
v_z=\frac{\partial\phi}{\partial z}=\frac{\partial\psi}{\partial x}$
\end{formula}
De esta forma, $\psi$ también satisface la ecuación de Laplace:
\[
\frac{\partial^2\psi}{\partial x^2}+\frac{\partial^2\psi}{\partial z^2}=\nabla^2\psi=0
\]
Por otro lado, el que $\psi=\text{cte}$ permite definir los \textbf{tubos o
canales de flujo}:
\[
\Delta q=\int_{\psi_1}^{\psi_2}v_x\,\dd z=\int_{\psi_1}^{\psi_2}-\frac{\partial\psi}{\partial z}\dd z
=\int_{\psi_1}^{\psi_2}-\partial\psi=\psi_1-\psi_2=\text{cte}
\]

\subsection{Red de flujo}
Como las familias son ortogonales, se pueden obtener \textbf{redes de flujo}:
malla compuesta por líneas de flujo y líneas equipotenciales.

\begin{figure}[H]
\centering
\begin{tikzpicture}[font=\small,scale=1.1]
  \draw[flecha] (0,0) -- (7,0) node[below]{$x$};
  \draw[flecha] (0,0) -- (0,4) node[left]{$z$};
  % canal de flujo sombreado
  \fill[orange!40] (0.8,3.0) .. controls (3,1.8) and (4.5,1.5) .. (6.3,1.4) -- (6.3,0.95) .. controls (4.5,1.05) and (3,1.3) .. (0.6,2.5) -- cycle;
  \foreach \d in {0,0.45,0.9,1.35,1.8}
    \draw[orange!90!black,thick] (0.8,3.0-\d) .. controls (3,1.8-\d*0.9) and (4.5,1.5-\d*0.8) .. (6.3,1.4-\d*0.7);
  \foreach \x in {1.2,2.0,2.8,3.6,4.4,5.2}
    \draw[azulusm,thick] (\x-0.6,1.0) -- (\x+0.4,3.6-\x*0.28);
  \fill[cyan!30,opacity=0.5] (1.4,1.0) -- (2.4,3.0) -- (3.2,2.8) -- (2.2,1.0) -- cycle;
  \node[right] at (6.3,1.4) {Línea de flujo ($\psi$)};
  \node[right] at (6.3,0.95) {Canal de flujo};
  \node[above] at (4.2,2.9) {Línea equipotencial ($\phi$)};
  \node[align=center,font=\scriptsize] at (1.8,0.5) {Caída equipotencial\\($K\Delta h=\text{cte}$)};
\end{tikzpicture}
\caption{Red de flujo: líneas de flujo, líneas equipotenciales y canal de flujo.}
\end{figure}

\subsection{Caudal a través de una red de flujo}
Para un canal de flujo se tiene, por continuidad:
\[
\Delta Q=\text{cte}\quad\Rightarrow\quad K\frac{\Delta h_1}{l_1}\cdot a_1\cdot b=K\frac{\Delta h_2}{l_2}\cdot a_2\cdot b
\]
Como $b=\text{cte}$ y $\Delta\phi_i=K\Delta h_i$:
\[
\frac{\Delta\phi_1}{l_1}\cdot a_1=\frac{\Delta\phi_2}{l_2}\cdot a_2
\]
Si $\Delta\phi_i=\text{cte}$:
\[
\frac{a_1}{l_1}=\frac{a_2}{l_2}=\text{cte}\quad\Longrightarrow\quad\Delta q=K\Delta h\cdot\frac{a}{l}\quad\text{(caudal por unidad de ancho)}
\]
Considerando $\Delta h=h/N_d$, con $N_d$ = número de caídas de potencial:
\begin{formula}
$\displaystyle\Delta q=K\frac{h}{N_d}\cdot\frac al$
\end{formula}
Para la red de flujo completa, con $N_t$ = número de canales de flujo:
\begin{formula}
$\displaystyle q=\sum\Delta q=N_t\,\Delta q=K\,h\,\frac{N_t}{N_d}\cdot\frac{a}{l}$
\end{formula}
\begin{itemize}
  \item $a/l=1$ $\Rightarrow$ elementos cuadrados.
  \item $a/l\neq1$ $\Rightarrow$ elementos rectangulares.
\end{itemize}

\subsection{Condiciones de borde}
\begin{itemize}
  \item \textbf{Barrera impermeable:} la velocidad es tangente a la barrera
        en todo momento $\Rightarrow$ la barrera impermeable es una \emph{línea
        de flujo} ($\psi$).
  \item \textbf{Eje de simetría:} no hay caudal que atraviese el eje de
        simetría $\Rightarrow$ el eje de simetría es una \emph{línea de flujo} ($\psi$).
  \item \textbf{Superficie de agua libre en reposo:} es una \emph{línea
        equipotencial} ($\phi$):
  \[h=z+\frac{p}{\gamma}=\text{cte}\quad\Rightarrow\quad\phi=Kh=\text{cte}\]
  \item \textbf{Superficie libre en movimiento:} es una \emph{línea de flujo}
        ($\psi$) con $p/\gamma=0$:
  \[\phi_1=Kz_1,\quad\phi_2=Kz_2\quad\Rightarrow\quad\Delta\phi=K\Delta z\]
\end{itemize}

\begin{figure}[H]
\centering
\begin{tikzpicture}[font=\small,scale=1.1]
  % suelo
  \fill[gray!15] (-4,-3.2) rectangle (4,0);
  \fill[pattern=north east lines] (-4,-3.45) rectangle (4,-3.2);
  % agua
  \fill[agua!70] (-4,0) rectangle (0,1.2);
  \fill[agua!70] (0,0) rectangle (4,0.6);
  \draw[rojo,thick] (-4,0) -- (4,0);
  \draw[cota,rojo] (-3.5,0) -- (-3.5,1.2) node[midway,right]{$h_1$};
  \draw[cota,rojo] (0.5,0) -- (0.5,0.6) node[midway,right]{$h_2$};
  \pic at (-2,1.22) {nivelagua}; \pic at (2,0.62) {nivelagua};
  % tablestaca
  \draw[line width=2.5pt] (0,1.6) -- (0,-1);
  % líneas de flujo: semicircunferencias alrededor de la punta
  \foreach \r in {0.5,1.1,1.8,2.6}
    \draw[azulusm,thick] (-\r,0) arc (180:360:\r);
  % equipotenciales: rayos desde la punta de la tablestaca
  \begin{scope}
    \clip (-4,-3.2) rectangle (4,0);
    \foreach \a in {200,225,250,270,290,315,340}
      \draw[rojo,dashed] (0,-1) -- ++(\a:5);
  \end{scope}
  \node[rojo,below] at (-3.6,0) {$\phi_1=Kh_1$};
  \node[rojo,below] at (3.5,0) {$\phi_9=Kh_2$};
  \draw[<-] (0.05,1.4) -- (1.2,1.9) node[right]{Eje de simetría / tablestaca};
  \node[azulusm] at (-2.4,-2.9) {$\psi$}; \node[rojo] at (2.9,-2.9) {$\phi$};
\end{tikzpicture}
\caption{Red de flujo bajo una tablestaca: las superficies de agua en reposo
son equipotenciales y la tablestaca es una línea de flujo.}
\end{figure}

\subsection{Trazado de una red de flujo: procedimiento}
\begin{enumerate}
  \item Dibujar las condiciones de borde (límites del dominio).
  \item Dibujar tentativamente 3 o 4 líneas de corriente.
  \item Trazar tentativamente equipotenciales, ortogonales a las líneas de corriente.
  \item Ajustar (verificar que se generen ángulos rectos).
  \item Validar la red de flujo: trazar las líneas diagonales a los
        rectángulos o cuadrados que forman las líneas de flujo y
        equipotenciales. Estas líneas se deben interceptar también en ángulos
        rectos. Si esto no es así, volver al paso 4.
\end{enumerate}

\subsection{Medio heterogéneo}
Al pasar de un medio con $K_1$ a otro con $K_2$, por continuidad en un canal de flujo:
\[
\Delta q=K_1\frac{\Delta h_1}{l_1}a_1=K_2\frac{\Delta h_2}{l_2}a_2
\]
Como $\Delta\phi=Kh=\text{cte}\Rightarrow\Delta h_1=\Delta h_2$. Además,
\[
\operatorname{tg}(\alpha)=\frac{l_1}{a_1},\qquad\operatorname{tg}(\beta)=\frac{l_2}{a_2}
\]
\begin{formula}
$\displaystyle\frac{K_1}{K_2}=\frac{\operatorname{tg}(\alpha)}{\operatorname{tg}(\beta)}$\qquad Ley de refracción (análoga a la óptica)
\end{formula}

\begin{figure}[H]
\centering
\begin{tikzpicture}[font=\small]
  \draw (0,-1.5) -- (0,3) node[above left]{$K_1$} node[above right]{$K_2$};
  \draw[dashed] (-4,0) -- (4,0);
  \draw[azulusm,thick] (-4,2.6) -- (0,0.7) -- (4,0.2);
  \draw[azulusm,thick] (-4,1.6) -- (0,-0.3) -- (4,-0.8);
  \foreach \t in {0.3,0.55,0.8}{
    \draw[rojo,thick] ($(-4,2.6)!\t!(0,0.7)$) -- ($(-4,1.6)!\t!(0,-0.3)$);}
  \foreach \t in {0.15,0.4,0.65}{
    \draw[rojo,thick] ($(0,0.7)!\t!(4,0.2)$) -- ($(0,-0.3)!\t!(4,-0.8)$);}
  \draw[flecha] (-3.3,2.8) -- (-2.6,2.45) node[pos=0,left]{$\Delta q$};
  \draw[flecha] (2.8,0.8) -- (3.6,0.7) node[pos=0,above]{$\Delta q$};
  \draw (-0.8,-0.3) arc (180:155:0.8); \node at (-1.1,-0.5) {$\alpha$};
  \draw (0.8,-0.3) arc (0:-8:0.8); \node at (1.1,-0.55) {$\beta$};
  \node at (-2,-1.2) {\color{rojo}$\phi$ \color{azulusm}$\psi$};
\end{tikzpicture}
\caption{Refracción de las líneas de flujo en la interfaz entre dos medios.}
\end{figure}

\subsection{Medio anisotrópico}
\[
K_x\frac{\partial^2h}{\partial x^2}+K_z\frac{\partial^2h}{\partial z^2}=0
\quad\Longrightarrow\quad
\frac{\partial^2h}{(K_z/K_x)\,\partial x^2}+\frac{\partial^2h}{\partial z^2}=0
\]
¡Las familias de curvas $\psi(x,z)$ y $\phi(x,z)$ \textbf{no son ortogonales}!
Sea el cambio de variable
\[
x'=a\cdot x\ \Rightarrow\ \frac{\dd x'}{\dd x}=a,\qquad
\frac{\partial h}{\partial x}=\frac{\partial h}{\partial x'}\frac{\dd x'}{\dd x}=a\frac{\partial h}{\partial x'},\qquad
\frac{\partial^2h}{\partial x^2}=a^2\frac{\partial^2h}{\partial x'^2}
\]
\[
\Longrightarrow\quad\frac{K_x}{K_z}a^2\frac{\partial^2h}{\partial x'^2}+\frac{\partial^2h}{\partial z^2}=0
\]
Si
\begin{formula}
$\displaystyle a=\sqrt{\frac{K_z}{K_x}}\quad\Longrightarrow\quad\frac{\partial^2h}{\partial x'^2}+\frac{\partial^2h}{\partial z^2}=0$\quad (ecuación de Laplace),
\qquad $K_{eq}=\sqrt{K_x\cdot K_z}$
\end{formula}
Procedimiento: se transforma el dominio $(x,z)\to(x',z)=(a\,x,z)$, se traza
la red de flujo ortogonal en el dominio transformado usando $K_{eq}$, y luego
se vuelve al dominio real $(x,z)=\left(\frac{x'}{a},z\right)$.

\paragraph{Ejemplo:} $K_x\neq K_z$ con $a=\sqrt{K_z/K_x}=10$: un terraplén que
en el dominio real se extiende entre $x=2$ y $x=10$ se transforma en uno que
se extiende entre $x'=20$ y $x'=100$.

\begin{figure}[H]
\centering
\begin{tikzpicture}[font=\small,scale=0.9]
  \begin{scope}
    \draw[flecha] (0,0) -- (4,0) node[below]{$x$}; \draw[flecha] (0,0) -- (0,2.5) node[left]{$z$};
    \fill[gray!30,draw] (1.0,0) -- (1.6,1.2) -- (2.6,1.2) -- (3.3,0) -- cycle;
    \draw[dashed] (0,1.2) -- (1.6,1.2);
    \node[below] at (1,0) {2}; \node[below] at (3.3,0) {10};
    \fill (2.0,0.6) circle (1.5pt) node[above,font=\scriptsize]{$(x,z)$};
  \end{scope}
  \draw[flecha,azulusm,line width=2pt] (4.6,1) -- (5.6,1) node[midway,above,black]{$x'=a\cdot x$};
  \begin{scope}[xshift=6.3cm]
    \draw[flecha] (0,0) -- (6.5,0) node[below]{$x'$}; \draw[flecha] (0,0) -- (0,2.5) node[left]{$z$};
    \fill[gray!30,draw] (0.8,0) -- (1.6,1.2) -- (5.2,1.2) -- (6.0,0) -- cycle;
    \draw[dashed] (0,1.2) -- (1.6,1.2);
    \node[below] at (0.8,0) {20}; \node[below] at (6.0,0) {100};
    \fill (2.6,0.6) circle (1.5pt) node[above,font=\scriptsize]{$(x',z)$};
    \node at (3.4,1.8) {$K_{eq}=\sqrt{K_x\cdot K_z}$};
  \end{scope}
\end{tikzpicture}
\caption{Transformación de coordenadas para un medio anisotrópico.}
\end{figure}

\section{Modelos analógicos}

\subsection{Analogía gradiente térmico}
Flujo de calor a través de una barra:
\[
F=\alpha\frac{\dd T}{\dd x}\quad\Longleftrightarrow\quad v=K\frac{\dd h}{\dd x}
\]
Es más complicado, pues el calor se escapa por cualquier lado del modelo.

\subsection{Analogía fluido viscoso (Hele--Shaw)}
En fluido newtoniano, de régimen laminar, entre dos placas paralelas
separadas una distancia $b$:
\[
\bar v=\frac{\gamma b^2}{12\mu}\frac{\dd h}{\dd x}\quad\Longleftrightarrow\quad v=K\frac{\dd h}{\dd x}
\]

\begin{figure}[H]
\centering
\begin{tikzpicture}[font=\small]
  \fill[gray!40] (0,1.2) rectangle (6,1.4); \fill[gray!40] (0,-1.4) rectangle (6,-1.2);
  \fill[left color=azulusm!80,right color=aguaclara] (0,-1.2) rectangle (6,1.2);
  \draw[dashed] (0,0) -- (6,0);
  \draw[cota] (1,-1.2) -- (1,1.2) node[midway,left]{$b$};
  \draw[thick] (3,-1.2) -- (3,1.2);
  \draw[thick,domain=-1.2:1.2,samples=30] plot ({3+0.8*(1-(\x/1.2)^2)},\x);
  \foreach \y in {-1.0,-0.8,...,1.0} \draw[->] (3,\y) -- ({3+0.8*(1-(\y/1.2)^2)},\y);
  \node at (4.2,0.9) {$v(y)$};
  \draw[flecha] (3,0) -- (4.5,0) node[below]{$x$};
  \draw[flecha] (3,0) -- (3,1.8) node[left]{$y$};
\end{tikzpicture}
\caption{Modelo de Hele--Shaw: flujo laminar entre placas paralelas.}
\end{figure}

\subsection{Analogía eléctrica}
Ley de Ohm:
\[
V=Ri\ \Rightarrow\ i=\frac1R V\quad\Longleftrightarrow\quad v=K\frac{\Delta h}{L}
\]
Donde las analogías son:
\begin{itemize}
  \item $V$ = gradiente hidráulico $\left(\frac{\Delta h}{L}\right)$.
  \item $i$ = velocidad del flujo ($v$).
  \item $\frac1R$ = permeabilidad ($K$).
\end{itemize}
Por lo tanto, son válidas las reglas de los circuitos en serie y paralelo.
Es más precisa que la analogía térmica.

\begin{figure}[H]
\centering
\begin{tikzpicture}[font=\small]
  \draw (0,0) rectangle (4,3.5);
  \fill[white] (-0.3,1.4) rectangle (0.3,2.1);
  \draw[thick] (-0.35,1.9) -- (0.35,1.9); \draw (-0.2,1.75) -- (0.2,1.75);
  \draw[thick] (-0.35,1.6) -- (0.35,1.6); \draw (-0.2,1.45) -- (0.2,1.45);
  \node[left] at (-0.4,1.7) {10};
  \fill[white] (3.2,0.4) rectangle (4.8,3.1);
  \draw (3.2,0.4) rectangle (4.8,3.1);
  \fill[gray!60] (3.4,2.7) rectangle (4.6,2.9); \fill[gray!60] (3.4,0.6) rectangle (4.6,0.8);
  \fill[aguaclara] (3.4,0.8) rectangle (4.6,2.7);
  \foreach \y/\l in {2.35/2,1.95/4,1.55/6,1.15/8} {\draw[rojo] (3.4,\y) -- (4.6,\y); \node[right,font=\scriptsize] at (4.6,\y) {\l};}
  \foreach \x in {3.7,4.0,4.3} \draw[azulusm] (\x,0.8) -- (\x,2.7);
  \node[font=\scriptsize] at (3.6,2.55) {$K$};
  \node[above right,font=\scriptsize] at (4,3.1) {$V=0$};
  \node[below right,font=\scriptsize] at (4,0.4) {$V=10$};
  \node[right] at (4.9,1.8) {$R$};
\end{tikzpicture}
\caption{Analogía eléctrica: papel conductor con equipotenciales medidas.}
\end{figure}

\section{Modelos numéricos}
Resolución de la ecuación de Laplace mediante:
\begin{itemize}
  \item Elementos finitos.
  \item Diferencias finitas.
\end{itemize}
Algunos programas: \textbf{MODFLOW} (USGS; software libre, diferencias
finitas), \textbf{Visual MODFLOW} y \textbf{FEFLOW} (DHI; elementos finitos).

% =====================================================================
%  CLASE 05 -- Acuíferos costeros, hidráulica de pozos en régimen
%              impermanente
% =====================================================================
\clase{05}{Acuíferos costeros e hidráulica de pozos en régimen impermanente}

\section{Acuíferos costeros: intrusión salina}

La intrusión salina se analiza como un \textbf{equilibrio estático de
columnas de agua de diferente densidad} (aproximación de Ghyben--Herzberg):
\begin{itemize}
  \item El flujo de agua dulce es perfectamente horizontal.
  \item No existe flujo de agua salada.
  \item La interfaz es un plano, no existiendo zona de mezcla.
\end{itemize}

\begin{figure}[H]
\centering
\begin{tikzpicture}[font=\small,scale=1.1]
  % agua subterránea dulce
  \fill[aguaclara] (0,0) -- (0.8,0) .. controls (2.5,0.5) and (4.5,1.6) .. (5.5,2.6)
     .. controls (4,2.85) and (2,3.1) .. (0,3.2) -- cycle;
  % agua subterránea salada
  \fill[azulusm!75] (0.8,0) .. controls (2.5,0.5) and (4.5,1.6) .. (5.5,2.6)
     .. controls (6.5,2.0) and (7.5,1.1) .. (9,0.8) -- (9,0) -- cycle;
  % mar
  \fill[azulusm!45] (5.5,2.6) .. controls (6.5,2.0) and (7.5,1.1) .. (9,0.8) -- (9,2.6) -- cycle;
  % zona no saturada
  \fill[gray!35] (0,3.2) .. controls (2,3.1) and (4,2.85) .. (5.5,2.6)
     .. controls (4,3.2) and (2,3.7) .. (0,3.8) -- cycle;
  \draw (0,3.8) .. controls (2,3.7) and (4,3.2) .. (5.5,2.6) .. controls (6.5,2.0) and (7.5,1.1) .. (9,0.8);
  \draw[thick,azulusm] (0,3.2) .. controls (2,3.1) and (4,2.85) .. (5.5,2.6);
  \draw (5.5,2.6) -- (9,2.6);
  \fill[pattern=north east lines] (0,-0.25) rectangle (9,0);
  \pic at (0.4,3.18) {nivelagua}; \pic at (8.6,2.62) {nivelagua};
  \draw[dashed] (0,2.6) -- (5.5,2.6);
  \draw[cota] (3.3,2.6) -- (3.3,2.93) node[midway,left]{$h_d$};
  \draw[cota] (3.3,1.25) -- (3.3,2.6) node[midway,left]{$h_s$};
  \draw[dashed] (3.3,1.25) -- (8.4,1.25);
  \draw[cota] (8.4,1.25) -- (8.4,2.6) node[midway,left]{$h_s$};
  \node at (0.5,2.2) {$\gamma_d$}; \node[align=center] at (1.2,1.4) {Agua\\dulce};
  \node at (8.7,1.0) {$\gamma_s$}; \node[align=center,white] at (6.0,0.6) {Agua\\salada};
  \node at (8.6,2.2) {Mar};
  \draw[<-] (5.5,2.65) -- (5.1,3.5) node[above left,azulusm]{Línea de costa};
  \draw[<-] (4.3,1.6) -- (7.6,3.4) node[above,azulusm]{Intrusión salina};
\end{tikzpicture}
\caption{Equilibrio de agua dulce y salada en un acuífero costero.}
\end{figure}

Las presiones sobre la interfaz son:
\[
P_d=\gamma_d\,(h_d+h_s),\qquad P_s=\gamma_s\,h_s
\]
Considerando $P_s=P_d$:
\begin{formula}
$\displaystyle h_s=h_d\left(\frac{\gamma_d}{\gamma_s-\gamma_d}\right)$
\end{formula}
Considerando que $\gamma_d=1000$ [kgf/m$^3$] y $\gamma_s=1025$ [kgf/m$^3$]:
\begin{formula}
$h_s=40\,h_d$
\end{formula}
La manera segura de evitar que un pozo entregue agua salada es que su
profundidad no pase el nivel del mar: al bombear, el descenso del nivel
de agua dulce provoca un ascenso de la interfaz 40 veces mayor
(\emph{upconing}).

\begin{nota}[title={Ejemplo: intrusión salina en una isla}]
Una isla de forma aproximadamente circular de 10 [km] de diámetro está
geológicamente constituida por suelos con $K=1\times10^{-4}$ [m/s]. La isla
recibe una recarga por infiltración de lluvias de aproximadamente 4 [mm/día],
formando un acuífero de agua dulce rodeado de agua salada de densidad
relativa $\gamma_{sr}=1.025$.
\begin{enumerate}[label=\alph*)]
  \item Estimar la cota sobre el nivel del mar a la que se encontrará la napa
        subterránea en su punto más alto en la isla.
  \item ¿A qué profundidad se hallará agua salada a 1 [km] de distancia de la costa?
\end{enumerate}
\end{nota}

\section{Ecuaciones generales (recordatorio)}
Conservación de la masa, flujo saturado:
\begin{formula}
$\displaystyle
\frac{\partial}{\partial x}\left(K_x\frac{\partial h}{\partial x}\right)+
\frac{\partial}{\partial y}\left(K_y\frac{\partial h}{\partial y}\right)+
\frac{\partial}{\partial z}\left(K_z\frac{\partial h}{\partial z}\right)=S_s\frac{\partial h}{\partial t}$
\end{formula}
Flujo transiente en medio homogéneo e isotrópico:
\[
\frac{\partial^2h}{\partial x^2}+\frac{\partial^2h}{\partial y^2}+\frac{\partial^2h}{\partial z^2}=\frac{S_s}{K}\frac{\partial h}{\partial t}
\]

\section{Acuíferos confinados: régimen transiente}

\begin{itemize}
  \item En un acuífero confinado el agua se obtiene del almacenamiento
        específico del acuífero, $S_s=S/b$, con $b$ el espesor del acuífero.
  \item En un acuífero confinado, $S$ es generalmente muy pequeño (0.005 o
        menos), entonces el bombeo afecta un área relativamente extensa del acuífero.
  \item Acuífero horizontal, confinado en la base y techo. Sin fuente de recarga.
  \item Los coeficientes $K$, $S$ son constantes e independientes del tiempo
        durante el bombeo.
  \item Caudal de bombeo a tasa constante, sin llegar a régimen permanente.
  \item El agua de almacenamiento se moviliza hacia el pozo de bombeo
        inmediatamente después de producirse las depresiones de la napa (el
        agua se libera instantánea e inmediatamente a la extracción).
  \item En el caso del régimen transiente (no se alcanza la estabilización del
        cono), cuando se indica un $Q$ de extracción se debe hacer referencia
        al tiempo desde que se inicia el bombeo.
\end{itemize}
\textcolor{rojo}{Estamos resolviendo el sistema para una distancia radial $r$
del pozo de extracción, de modo que ahora la altura piezométrica variará en
el tiempo y en el espacio (no se alcanza el equilibrio).}

% Macro: pozo en régimen transiente
\newcommand{\pozotransiente}[1]{% #1 = 1 confinado, 0 libre
\begin{tikzpicture}[font=\small,scale=0.9]
  \draw[thick] (-4,0) -- (4,0);
  \foreach \x in {-3.9,-3.6,...,3.9} \draw (\x,0) -- (\x-0.2,-0.25);
  \draw[thick] (-0.4,0.05) -- (-0.4,4.4) (0.4,0.05) -- (0.4,4.4);
  \draw[thick] (-4,4.4) -- (-0.4,4.4) (0.4,4.4) -- (4,4.4);
  \fill[azulusm] (-0.1,4.2) rectangle (0.1,4.9);
  \fill[azulusm] (-0.25,4.9) -- (0.25,4.9) -- (0,5.25) -- cycle; \node at (0.55,5.1) {$Q$};
  \draw[cyan!70!blue,thick] (-4,3.6) -- (-0.4,3.6) (0.4,3.6) -- (4,3.6);
  \node[left] at (-4,3.6) {$t=0$};
  \draw[dashed,cyan!70!blue,thick] (-4,3.1) .. controls (-2,2.9) and (-1,2.6) .. (-0.4,2.3);
  \draw[dashed,cyan!70!blue,thick] (0.4,2.3) .. controls (1,2.6) and (2,2.9) .. (4,3.1);
  \draw[cyan!70!blue,thick] (-0.4,2.3) -- (0.4,2.3);
  \node[left] at (-4,3.1) {$t=t$};
  \ifnum#1=1
    \draw[thick] (-4,2.0) -- (-0.4,2.0) (0.4,2.0) -- (4,2.0);
    \foreach \x in {-3.9,-3.6,...,-0.6,0.6,0.9,...,3.9} \draw (\x,2.0) -- (\x+0.15,2.25);
    \draw[cota] (-3.3,0) -- (-3.3,2.0) node[midway,left]{$m$};
    \def\ytop{1.8}
  \else
    \def\ytop{2.0}
  \fi
  \foreach \y in {0.3,0.6,...,\ytop}{\draw[->,azulusm] (-1.3,\y) -- (-0.5,\y);\draw[->,azulusm] (1.3,\y) -- (0.5,\y);}
  \node[right] at (1.3,1.0) {$Q_i$};
  \draw[cota] (-2.3,0) -- (-2.3,3.6) node[midway,right]{$H$};
  \draw[cota] (2.8,0) -- (2.8,2.95) node[midway,right]{$h(r,t)$};
  \draw[->] (0,-0.5) -- (2.8,-0.5) node[midway,below]{$r$};
\end{tikzpicture}}

\begin{figure}[H]
\centering
\pozotransiente{1}
\caption{Pozo en acuífero confinado en régimen transiente.}
\end{figure}

\subsection{Solución de Theis (1935)}
Hipótesis: acuífero horizontal, confinado entre formaciones impermeables, de
extensión horizontal infinita, espesor constante, homogéneo e isotrópico,
pozo de radio infinitesimalmente pequeño. Se asume además que $Q_i=Q_s$ (no
hay acumulación en el pozo), lo que se cumple si el radio del pozo es pequeño.

En coordenadas cilíndricas:
\begin{formula}
$\displaystyle\frac{\partial^2h}{\partial r^2}+\frac1r\frac{\partial h}{\partial r}=\frac ST\frac{\partial h}{\partial t}$
\end{formula}
La solución de Theis es:
\begin{formula}
$\displaystyle\Delta h(r,t)=H-h(r,t)=\frac{Q}{4\pi T}W(u),\qquad
W(u)=\int_u^\infty\frac{e^{-u}}{u}\,\dd u,\qquad u=\frac{r^2\cdot S}{4\cdot T\cdot t}$
\end{formula}
$W(u)$ es conocida como la \textbf{función del pozo} (\emph{well function}):
\[
W(u)=\int_u^\infty\frac{e^{-u}}{u}\,\dd u=-0.577216-\ln u+u-\frac{u^2}{2\cdot2!}+\frac{u^3}{3\cdot3!}-\dots
\]

\subsection{Aproximación de Cooper--Jacob (1946)}
Si $r$ es pequeño ($u<0.01$) y $t\to\infty$, basta con los dos primeros
términos de la serie:
\begin{formula}
$\displaystyle H-h(r,t)\approx\frac{Q}{4\pi T}\ln\left(\frac{2.24\cdot T\cdot t}{S\cdot r^2}\right)$\qquad\textbf{Cooper--Jacob (1946)}
\end{formula}
\begin{table}[H]
\centering
\begin{tabular}{cc}
\toprule
$u$ & Error\\
\midrule
0.20 & 20\%\\
0.14 & 10\%\\
0.09 & 5\%\\
\bottomrule
\end{tabular}
\caption{Error de la aproximación de Cooper--Jacob.}
\end{table}
De este modo, conociendo las características hidrogeológicas ($S$, $T$), se
puede estimar el descenso en cualquier parte del acuífero, una vez iniciado el bombeo.

\subsection{Función del pozo, $W(u)$}
\begin{figure}[H]
\centering
\begin{tikzpicture}
\begin{loglogaxis}[width=11cm,height=7cm,xmin=0.1,xmax=1e4,ymin=0.01,ymax=10,
  xlabel={$1/u$},ylabel={$W(u)$},grid=both,minor grid style={gray!20},major grid style={gray!50}]
\addplot[thick,cyan!70!blue] coordinates {(0.3162,0.01063) (0.3981,0.02453) (0.5012,0.04922) (0.631,0.08824) (0.7943,0.1444) (1,0.2194) (1.259,0.3138) (1.585,0.4272) (1.995,0.5583) (2.512,0.7056) (3.162,0.867) (3.981,1.041) (5.012,1.225) (6.31,1.417) (7.943,1.617) (10,1.823) (12.59,2.034) (15.85,2.248) (19.95,2.466) (25.12,2.686) (31.62,2.908) (39.81,3.132) (50.12,3.357) (63.1,3.583) (79.43,3.81) (100,4.038) (125.9,4.266) (158.5,4.495) (199.5,4.724) (251.2,4.953) (316.2,5.182) (398.1,5.412) (501.2,5.642) (631,5.872) (794.3,6.102) (1000,6.332) (1259,6.562) (1585,6.792) (1995,7.022) (2512,7.252) (3162,7.482) (3981,7.712) (5012,7.943) (6310,8.173) (7943,8.403) (1e+04,8.633)};
\end{loglogaxis}
\end{tikzpicture}
\caption{Curva tipo de Theis: $W(u)$ versus $1/u$.}
\end{figure}

\begin{table}[H]
\centering\footnotesize
\begin{tabular}{c|ccccc}
\toprule
$N$ & $N\times10^{-14}$ & $N\times10^{-13}$ & $N\times10^{-12}$ & $N\times10^{-11}$ & $N\times10^{-10}$\\
\midrule
1 & 31.66 & 29.36 & 27.05 & 24.75 & 22.45\\
2 & 30.97 & 28.66 & 26.36 & 24.06 & 21.76\\
3 & 30.56 & 28.26 & 25.96 & 23.65 & 21.35\\
4 & 30.27 & 27.97 & 25.67 & 23.36 & 21.06\\
5 & 30.05 & 27.75 & 25.44 & 23.14 & 20.84\\
6 & 29.87 & 27.56 & 25.26 & 22.96 & 20.66\\
7 & 29.71 & 27.41 & 25.11 & 22.81 & 20.50\\
8 & 29.58 & 27.28 & 24.97 & 22.67 & 20.37\\
9 & 29.46 & 27.16 & 24.86 & 22.55 & 20.25\\
\bottomrule
\end{tabular}\par\medskip
\begin{tabular}{c|ccccc}
\toprule
$N$ & $N\times10^{-9}$ & $N\times10^{-8}$ & $N\times10^{-7}$ & $N\times10^{-6}$ & $N\times10^{-5}$\\
\midrule
1 & 20.15 & 17.84 & 15.54 & 13.24 & 10.94\\
2 & 19.45 & 17.15 & 14.85 & 12.55 & 10.24\\
3 & 19.05 & 16.74 & 14.44 & 12.14 & 9.84\\
4 & 18.76 & 16.46 & 14.15 & 11.85 & 9.55\\
5 & 18.54 & 16.23 & 13.93 & 11.63 & 9.33\\
6 & 18.35 & 16.05 & 13.75 & 11.45 & 9.14\\
7 & 18.20 & 15.90 & 13.59 & 11.29 & 8.99\\
8 & 18.07 & 15.76 & 13.46 & 11.16 & 8.86\\
9 & 17.95 & 15.65 & 13.34 & 11.04 & 8.74\\
\bottomrule
\end{tabular}\par\medskip
\begin{tabular}{c|ccccc}
\toprule
$N$ & $N\times10^{-4}$ & $N\times10^{-3}$ & $N\times10^{-2}$ & $N\times10^{-1}$ & $N\times10^{0}$\\
\midrule
1 & 8.63 & 6.33 & 4.04 & 1.82 & 0.22\\
2 & 7.94 & 5.64 & 3.35 & 1.22 & $4.9\times10^{-2}$\\
3 & 7.53 & 5.23 & 2.96 & 0.91 & $1.3\times10^{-2}$\\
4 & 7.25 & 4.95 & 2.68 & 0.70 & $3.8\times10^{-3}$\\
5 & 7.02 & 4.73 & 2.47 & 0.56 & $1.1\times10^{-3}$\\
6 & 6.84 & 4.54 & 2.30 & 0.45 & $3.6\times10^{-4}$\\
7 & 6.69 & 4.39 & 2.15 & 0.37 & $1.2\times10^{-4}$\\
8 & 6.55 & 4.26 & 2.03 & 0.31 & $3.8\times10^{-5}$\\
9 & 6.44 & 4.14 & 1.92 & 0.26 & $1.2\times10^{-5}$\\
\bottomrule
\end{tabular}\par\medskip
\caption{Valores de la función del pozo $W(u)$ para $u=N\times10^{k}$.}
\end{table}

\subsection{Radio de influencia dependiente del tiempo}
En equilibrio: $\displaystyle\Delta h=\frac{Q}{2\pi T}\ln\left(\frac Rr\right)$.
A partir de Cooper--Jacob se define:
\begin{formula}
$\displaystyle R(t)=\sqrt{\frac{2.24\cdot T\cdot t}{S}}
\quad\Longrightarrow\quad
\Delta h=\frac{Q}{2\pi T}\ln\left(\frac{R(t)}{r}\right)$\qquad (desequilibrio)
\end{formula}

\section{Acuíferos libres: régimen transiente}
Ecuación gobernante:
\[
2\frac hr\frac{\partial h}{\partial r}+\frac{\partial^2h}{\partial r^2}=\frac{S}{K\cdot h}\frac{\partial h}{\partial t}
\]
Siguiendo el resultado del caso confinado, se establece la analogía
(en equilibrio $H^2-h^2=\frac{Q}{\pi K}\ln\frac Rr$):
\begin{formula}
$\displaystyle R(t)=\sqrt{\frac{2.24\cdot T\cdot t}{S}}\quad\Longrightarrow\quad
H^2-h(r,t)^2=\frac{Q}{\pi K}\ln\left(\frac{R(t)}{r}\right)$\qquad (desequilibrio)
\end{formula}

\begin{figure}[H]
\centering
\pozotransiente{0}
\caption{Pozo en acuífero libre en régimen transiente.}
\end{figure}

\section{Interferencia de pozos: régimen impermanente}
\paragraph{Depresión producto de múltiples pozos.} Similar a régimen
permanente $\Rightarrow$ solución: superposición.
\begin{formula}
$\displaystyle\Delta h_T=\sum_{i=1}^n\Delta h_i(r_i)$
\end{formula}
Siempre que $R(t)>r_i$. Esta restricción define el concepto de
\textbf{tiempo de influencia} ($t_i$):
\[
R(t)=\sqrt{\frac{2.24\cdot T\cdot t_i}{S}}=r_i\quad\Longrightarrow\quad
\boxed{t_i=\frac{r_i^2\cdot S}{2.24\cdot T}}
\]

Usando la solución de Jacob (acuíferos confinados):
\begin{itemize}
  \item Si todos los pozos afectan al punto ``A'':
  \begin{formula}
  $\displaystyle\Delta h_A=\sum_{i=1}^n\frac{Q_i}{4\cdot\pi\cdot T}\ln\left(\frac{2.24\cdot T\cdot t}{r_{i-A}^2\cdot S}\right),
  \qquad r_{i-A}\leq\sqrt{\frac{2.24\cdot T\cdot t}{S}}$
  \end{formula}
  \item Si todos los pozos extraen el mismo caudal:
  \begin{formula}
  $\displaystyle\Delta h_A=Q\sum_{i=1}^n\frac{1}{4\cdot\pi\cdot T}\ln\left(\frac{2.24\cdot T\cdot t}{r_{i-A}^2\cdot S}\right),
  \qquad r_{i-A}\leq\sqrt{\frac{2.24\cdot T\cdot t}{S}}$
  \end{formula}
  \item ¡Para acuíferos libres, linealizar primero!
\end{itemize}

\section{Método de las imágenes en régimen impermanente}
Superposición de soluciones para obtener la respuesta a condiciones de borde
específicas: recargas lineales (masas de agua con nivel constante) y barreras
impermeables. Es similar a régimen permanente, pero se debe tener en cuenta
el tiempo $t$ en el que se toma la medición y los tiempos de influencia de
los pozos ($t_r$: pozo real, $t_i$: pozo imagen).

\paragraph{Recarga lineal}
\begin{itemize}
  \item Si $t_r<t<t_i$:\quad $\displaystyle\Delta h_T=\Delta h_{real}=\frac{Q}{2\pi T}\ln\left(\frac{R(t)}{r_r}\right)$
  \item Si $t>t_i$:\quad $\Delta h_T=\Delta h_{real}-\Delta h_{imag}$
  \begin{formula}
  $\displaystyle\Delta h_T=\frac{Q}{2\pi T}\ln\left(\frac{r_i}{r_r}\right)$
  \end{formula}
\end{itemize}

\paragraph{Barrera impermeable}
\begin{itemize}
  \item Si $t_r<t<t_i$:\quad $\displaystyle\Delta h_T=\Delta h_{real}=\frac{Q}{2\pi T}\ln\left(\frac{R(t)}{r_r}\right)$
  \item Si $t>t_i$:\quad $\Delta h_T=\Delta h_{real}+\Delta h_{imag}$
  \begin{formula}
  $\displaystyle\Delta h_T=\frac{Q}{2\pi T}\ln\left(\frac{R(t)^2}{r_r\cdot r_i}\right)$
  \end{formula}
\end{itemize}

\section{Bombeo variable}
Usando la solución de Jacob (acuíferos confinados), un bombeo intermitente
se descompone como la superposición de pozos que se encienden ($+Q$) y se
apagan ($-Q$) en los instantes $t_i$.

\begin{figure}[H]
\centering
\begin{tikzpicture}[font=\small]
  \begin{scope}[yshift=3.2cm]
    \draw[flecha] (0,0) -- (6.5,0) node[below]{$t$};
    \draw[flecha] (0,0) -- (0,2.2) node[left]{$Q$};
    \foreach \a in {0,2,4} \fill[azulusm!80] (\a,0) rectangle (\a+1,1.5);
    \node[left] at (0,1.5) {$Q_1$};
    \foreach \x/\n in {1/1,2/2,3/3,4/4,5/5} \node[below] at (\x,0) {$t_\n$};
  \end{scope}
  \draw[flecha] (0,0) -- (6.5,0) node[below]{$t$};
  \draw[flecha] (0,-1.8) -- (0,2.4) node[left]{$Q$};
  \draw[azulusm,very thick] (0,0.5) -- (6,0.5); \node[left] at (0,0.5) {$Q_1$};
  \draw[azulusm,very thick] (2,1.1) -- (6,1.1); \node[left] at (2,1.1) {$Q_1$};
  \draw[azulusm,very thick] (4,1.7) -- (6,1.7); \node[left] at (4,1.7) {$Q_1$};
  \draw[gray,very thick] (1,-0.5) -- (6,-0.5); \node[left] at (1,-0.5) {$-Q_1$};
  \draw[gray,very thick] (3,-1.0) -- (6,-1.0); \node[left] at (3,-1.0) {$-Q_1$};
  \draw[gray,very thick] (5,-1.5) -- (6,-1.5); \node[left] at (5,-1.5) {$-Q_1$};
  \foreach \x in {1,2,3,4,5} \draw[dashed,gray] (\x,-1.6) -- (\x,3.2);
\end{tikzpicture}
\caption{Bombeo intermitente descompuesto en pozos ficticios que se encienden y apagan.}
\end{figure}

Por lo tanto,
\[
\Delta h=\sum\Delta h_i=\frac{1}{4\cdot\pi\cdot T}\sum_{i=1}^n\pm Q\cdot\ln\left(\frac{2.24\cdot T}{r^2\cdot S}(t-t_i)\right),
\qquad\text{para } t\geq t_i
\]
Generalizando:
\begin{formula}
$\displaystyle\Delta h=\sum\Delta h_i=\frac{1}{4\cdot\pi\cdot T}\sum_{i=1}^n\pm Q_i\cdot\ln\left(\frac{2.24\cdot T}{r^2\cdot S}(t-t_i)\right),
\qquad t\geq t_i$
\end{formula}

% =====================================================================
%  CLASE 06 -- Pruebas de bombeo
% =====================================================================
\clase{06}{Pruebas de bombeo}

\section{Resumen}
\begin{table}[H]
\centering
\renewcommand{\arraystretch}{2.2}
\begin{tabular}{l>{$\displaystyle}c<{$}}
\toprule
\textbf{Caso} & \textbf{Ecuación}\\
\midrule
Penetración total, acuífero confinado & \Delta h=H-h=\frac{Q}{2\pi T}\ln\left(\frac Rr\right)\\
Penetración total, acuífero libre & H^2-h^2=\frac{Q}{\pi K}\ln\left(\frac Rr\right)\\
Acuífero artesiano, impermanente & \Delta h=H-h=\frac{Q}{2\pi T}\ln\left(\frac{R(t)}{r}\right)\\
Acuífero libre, impermanente & H^2-h^2=\frac{Q}{\pi K}\ln\left(\frac{R(t)}{r}\right)\\
Radio de influencia transiente & R(t)=\sqrt{\frac{2.24\cdot T\cdot t}{S}}\\
\bottomrule
\end{tabular}
\end{table}

\section{Pruebas de bombeo}
\begin{itemize}
  \item Para determinar parámetros del acuífero ($K$, $S$).
  \item Extracción de agua y control de niveles en el pozo.
  \item En condiciones de equilibrio y desequilibrio.
  \item Cálculo del comportamiento a largo plazo de un pozo.
  \item Cálculo del radio de protección de un pozo (para derechos de agua, por ejemplo).
  \item Ideal contar con pozos de observación.
  \item Generalmente las pruebas se realizan en el mismo pozo, por lo que los
        resultados presentan incertidumbre.
\end{itemize}

\begin{figure}[H]
\centering
\begin{tikzpicture}[font=\scriptsize,scale=1.0]
  \fill[suelo!60] (0,0) rectangle (9,0.4); \node at (1.2,0.2) {Basamento};
  \fill[aguaclara] (0,0.4) rectangle (9,2.6);
  \fill[sueloclaro] (0,2.6) -- (9,2.6) -- (9,3.3) -- (0,3.3) -- cycle;
  \fill[aguaclara] (0,2.6) -- (0,2.4) .. controls (2,2.3) and (3,1.3) .. (3.4,1.2) -- (3.8,1.2) .. controls (4.2,1.3) and (6,2.3) .. (9,2.4) -- (9,2.6) -- cycle;
  \fill[sueloclaro] (0,2.4) .. controls (2,2.3) and (3,1.3) .. (3.4,1.2) -- (3.8,1.2) .. controls (4.2,1.3) and (6,2.3) .. (9,2.4) -- (9,2.6) -- (0,2.6) -- cycle;
  \draw[dashed,azulusm,thick] (0,2.6) -- (9,2.6);
  \draw[azulusm,thick] (0,2.4) .. controls (2,2.3) and (3,1.3) .. (3.4,1.2) -- (3.8,1.2) .. controls (4.2,1.3) and (6,2.3) .. (9,2.4);
  \draw (0,3.3) -- (9,3.3);
  \fill[white,draw] (3.4,0.4) rectangle (3.8,3.9); \fill[agua] (3.4,0.4) rectangle (3.8,1.2);
  \fill[white,draw] (6.4,0.8) rectangle (6.6,3.9); \fill[agua] (6.4,0.8) rectangle (6.6,2.2);
  \draw[flecha,azulusm] (3.6,1.4) -- (3.6,4.3) -- (2.8,4.3) node[left,black]{Caudal de bombeo};
  \node[above] at (3.6,4.4) {Pozo de bombeo};
  \node[above,align=center] at (6.5,3.9) {Pozo de\\observación};
  \draw[cota] (7.0,2.2) -- (7.0,2.6) node[midway,right]{Descenso};
  \node[right,align=left] at (7.3,3.0) {Nivel estático original};
  \draw[<-] (2.2,1.9) -- (1.2,1.3) node[below,align=center]{Nivel de agua luego\\del bombeo};
  \draw[<-] (2.8,1.6) -- (1.0,3.7) node[above]{Cono de depresión};
  \node at (5,0.9) {Acuífero};
\end{tikzpicture}
\caption{Esquema de una prueba de bombeo con pozo de observación.}
\end{figure}

\subsection{Solicitud de derecho subterráneo (DGA)}
Según la \emph{Guía para la presentación de Solicitudes de Derechos de
Aprovechamiento de Aguas Subterráneas} (DGA, versión 25.05.2022), los
antecedentes técnicos mínimos son:
\begin{itemize}
  \item \textbf{Pozos profundos y punteras:} perfil estratigráfico;
        habilitación; \textbf{prueba de bombeo de gasto constante} para el
        caudal solicitado, con una duración de 24 horas como mínimo y con un
        tiempo de estabilización de niveles de 180 minutos (sin perjuicio de
        otros antecedentes, como una prueba de gasto variable, si la hubiere);
        los datos de recuperación completa de la obra, o hasta que se haya
        estabilizado la recuperación.
  \item \textbf{Norias y drenes:} prueba de bombeo de gasto constante para el
        caudal solicitado, con estabilización de niveles de por lo menos 180
        minutos; para drenes podrían requerirse registros de caudales medidos
        mediante aforos; para norias, los datos de recuperación completa de la obra.
  \item \textbf{Pozos zanjas:} prueba de bombeo de gasto constante para el
        caudal solicitado durante 24 horas, con estabilización de niveles de
        por lo menos 180 minutos; perfil estratigráfico y dimensiones de la obra.
\end{itemize}

\section{Prueba de gasto variable}
\begin{itemize}
  \item Para determinar el caudal máximo a extraer del pozo y la ubicación de la bomba.
  \item Se aumenta paulatinamente el caudal bombeado.
  \item Para cada caudal, se deja estabilizar el sistema.
  \item Si el acuífero es artesiano se obtiene una recta (uniendo los puntos
        de estabilización de cada escalón); si es libre, se puede linealizar el problema.
\end{itemize}

\begin{figure}[H]
\centering
\begin{tikzpicture}
\begin{axis}[width=10cm,height=6cm,xmin=0,xmax=18,ymin=0,ymax=25,
  xlabel={Tiempo [h]},ylabel={Descenso [m]},grid=major,grid style={dashed,gray!40}]
  \addplot[thick,black] coordinates {(0,0)(0.1,3)(0.5,4)(3,4.5)};
  \addplot[thick,red] coordinates {(3,4.5)(3.1,5.8)(3.5,6.2)(6,6.5)};
  \addplot[thick,green!60!black] coordinates {(6,6.5)(6.1,8.5)(6.5,9)(9,9.3)};
  \addplot[thick,blue] coordinates {(9,9.3)(9.1,12)(9.5,12.7)(12,13)};
  \addplot[thick,magenta] coordinates {(12,13)(12.1,16)(12.5,16.6)(15,17)};
  \addplot[thick,green!40!black] coordinates {(15,17)(15.1,20.5)(15.5,21.5)(18,22)};
  \addplot[very thick,dashed,azulusm] coordinates {(3,4.5)(6,6.5)(9,9.3)(12,13)(15,17)(18,22)};
\end{axis}
\end{tikzpicture}
\caption{Prueba de gasto variable (escalonada): descenso versus tiempo.}
\end{figure}

\section{Prueba de gasto constante}
\begin{itemize}
  \item Para determinar los parámetros elásticos del acuífero ($K$, $S$).
  \item En Chile se realizan por 24 horas, con 3 horas de estabilización.
  \item Se realiza para el 80\% del $Q$ máximo de la prueba de gasto variable.
\end{itemize}

\paragraph{Prueba de recuperación}
\begin{itemize}
  \item En un tiempo $t_f$ se deja de bombear, para que los niveles se recuperen.
  \item Permite la determinación de $K$.
\end{itemize}

\begin{figure}[H]
\centering
\begin{subfigure}{0.48\textwidth}\centering
\begin{tikzpicture}
\begin{axis}[width=\linewidth,height=5cm,xmin=0,xmax=250,ymin=0,ymax=1.2,y dir=reverse,
  axis x line*=top,xlabel={tiempo [min]},ylabel={$\Delta h$ [m]},grid=major,grid style={dashed,gray!40},font=\scriptsize]
  \addplot[only marks,mark=*,orange,mark size=1.5pt,domain=1:210,samples=40] {1.0*(1-exp(-x/15))};
  \addplot[only marks,mark=*,orange,mark size=1.5pt] coordinates {(0,0)};
\end{axis}
\end{tikzpicture}
\caption{Descenso en prueba de gasto constante.}
\end{subfigure}\hfill
\begin{subfigure}{0.48\textwidth}\centering
\begin{tikzpicture}[font=\small]
  \draw[flecha] (0,0) -- (5,0) node[right]{$t$};
  \draw[flecha] (0,0) -- (0,-2.6) node[left]{$\Delta h$};
  \draw[orange,very thick,dashed] (0,0) .. controls (0.5,-1.5) and (1.5,-1.8) .. (3,-1.95);
  \draw[orange,very thick] (3,-1.95) .. controls (3.3,-1.2) and (3.6,-0.8) .. (4.9,-0.3);
  \draw[dashed] (3,0.3) node[above]{$t_f$} -- (3,-2.2);
\end{tikzpicture}
\caption{Prueba de recuperación a partir de $t_f$.}
\end{subfigure}
\caption{Prueba de gasto constante y prueba de recuperación.}
\end{figure}

\section{Interpretación de la prueba de gasto constante}

\subsection{Método de la curva tipo (Theis)}
Solución gráfica propuesta por Theis. Como $u$ es sensible para valores de
$t$ pequeños, se definen dos soluciones aproximadas distintas de $W(u)$ con
más detalle, por medio de series de Taylor:
\begin{itemize}
  \item Para $u<1$:
  \begin{formula}
  $W(u)=u-0.2499u^2+0.0552u^3-0.0098u^4+0.00108u^5-\ln(u)-0.5772$
  \end{formula}
  \item Para $u\geq1$:
  \begin{formula}
  $\displaystyle W(u)=\frac{mH_1}{nH_1\cdot u\cdot e^u}$
  \end{formula}
  donde
  \[
  mH_1=u^4+8.573u^3+18.059u^2+8.635u+0.268,
  \]
  \[
  nH_1=u^4+9.573u^3+25.633u^2+21.097u+3.958
  \]
\end{itemize}

\paragraph{Procedimiento}
\begin{enumerate}
  \item Graficar $W(u)$ en función de $1/u$ (ejes logarítmicos).
  \item Graficar la depresión en el pozo en función del tiempo (ejes logarítmicos, misma escala).
  \item Mover los gráficos en dirección horizontal y vertical, hasta que los
        datos teóricos se ubiquen sobre los datos de terreno. ¡Cuidado de no
        girar los gráficos!
  \item Una vez conseguido el mejor ajuste, se selecciona un punto en el
        gráfico y se determina el valor de los 4 ejes: $W(u)$, $1/u$,
        $\Delta h^*$ y $t^*$.
  \item Con estos valores se determinan los parámetros del acuífero:
\end{enumerate}
\begin{formula}
$\displaystyle T=\frac{Q}{4\cdot\pi\cdot\Delta h^*}\cdot W(u)\qquad\qquad
S=\frac{4\cdot T\cdot u\cdot t^*}{r^2}$
\end{formula}
¡Imprescindible para tiempos de bombeo pequeños!

\begin{figure}[H]
\centering
\begin{tikzpicture}
\begin{loglogaxis}[width=10cm,height=7cm,xmin=0.1,xmax=1e4,ymin=0.01,ymax=10,
  xlabel={$1/u$},ylabel={$W(u)$},grid=both,minor grid style={gray!15},major grid style={gray!40},
  clip=false]
\addplot[thick,orange] coordinates {(0.3162,0.01063) (0.3981,0.02453) (0.5012,0.04922) (0.631,0.08824) (0.7943,0.1444) (1,0.2194) (1.259,0.3138) (1.585,0.4272) (1.995,0.5583) (2.512,0.7056) (3.162,0.867) (3.981,1.041) (5.012,1.225) (6.31,1.417) (7.943,1.617) (10,1.823) (12.59,2.034) (15.85,2.248) (19.95,2.466) (25.12,2.686) (31.62,2.908) (39.81,3.132) (50.12,3.357) (63.1,3.583) (79.43,3.81) (100,4.038) (125.9,4.266) (158.5,4.495) (199.5,4.724) (251.2,4.953) (316.2,5.182) (398.1,5.412) (501.2,5.642) (631,5.872) (794.3,6.102) (1000,6.332) (1259,6.562) (1585,6.792) (1995,7.022) (2512,7.252) (3162,7.482) (3981,7.712) (5012,7.943) (6310,8.173) (7943,8.403) (1e+04,8.633)};
\addplot[only marks,mark=o,azulusm,mark size=1.3pt] coordinates {(5,1.2)(8,1.6)(12,2.0)(20,2.45)(30,2.9)(50,3.35)(80,3.8)(120,4.2)(200,4.7)(300,5.1)(500,5.6)(800,6.1)(1200,6.5)(2000,7.0)};
\draw[dashed,rojo,thick] (axis cs:0.1,4.7) -- (axis cs:200,4.7) -- (axis cs:200,0.01);
\fill[rojo] (axis cs:200,4.7) circle (2pt);
\node[rojo,above] at (axis cs:200,5.2) {Punto elegido};
\node[rojo,left] at (axis cs:0.1,4.7) {$W(u)$, $\Delta h^*$};
\node[rojo,below] at (axis cs:200,0.01) {$1/u$, $t^*$};
\end{loglogaxis}
\end{tikzpicture}
\caption{Método de la curva tipo: superposición de las mediciones ($\Delta h$
vs.\ $t$, puntos) sobre la curva de Theis ($W(u)$ vs.\ $1/u$) y elección de un
punto de ajuste.}
\end{figure}

\subsection{Método de Cooper--Jacob (1946)}
\begin{itemize}
  \item Graficar los datos de terreno en papel semilog ($\Delta h$ vs.\ $\ln t$).
  \item Las constantes del acuífero se obtienen a partir del ajuste lineal.
  \item Se deben cumplir las hipótesis de Jacob ($u$ pequeño).
\end{itemize}
\begin{formula}
$\displaystyle T=\frac{Q}{4\cdot\pi\cdot a}\qquad a$: pendiente de la recta\\[6pt]
$\displaystyle S=\frac{2.24\cdot T\cdot t_0}{r^2}\qquad t_0$: punto artificial donde $\Delta h=0$
\end{formula}

\begin{figure}[H]
\centering
\begin{tikzpicture}
\begin{axis}[width=9cm,height=5.5cm,xmin=0,xmax=6,ymin=0,ymax=1.2,
  xlabel={$\ln(t)$},ylabel={Desnivel [m]},grid=major,grid style={dashed,gray!40}]
  \addplot[only marks,mark=*,orange,mark size=1.5pt] coordinates {(0,0.45)(0.7,0.5)(1.1,0.52)(1.4,0.55)(1.8,0.6)(2.1,0.65)(2.4,0.7)(2.7,0.74)(3.0,0.8)(3.3,0.86)(3.6,0.92)(3.9,0.97)(4.2,1.0)(4.5,1.03)(4.8,1.06)(5.0,1.08)(5.2,1.1)};
  \addplot[very thick,black,domain=0.2:5.3] {0.78+0.06*x};
\end{axis}
\end{tikzpicture}
\caption{Método de Cooper--Jacob: ajuste lineal en papel semilogarítmico.}
\end{figure}

\subsection{Método de Papadopulos--Cooper}
\begin{itemize}
  \item Útil para pozos de gran diámetro.
  \item No es posible despreciar el almacenamiento de agua dentro del pozo.
\end{itemize}
En régimen transiente se cumple $Q\neq Q_B$ con $Q_B>Q$. Considerando el
volumen de control definido por el pozo:
\[
Q_i-Q_s=\frac{\partial V}{\partial t}\;\Rightarrow\;Q_i-Q_s=\pi r_0^2\frac{\partial h_0}{\partial t}
\]
Por otro lado, considerando
\[
\frac{\partial^2h}{\partial r^2}+\frac1r\frac{\partial h}{\partial r}=\frac ST\frac{\partial h}{\partial t}
\]
con $r_0$: radio efectivo del pozo (entubación) y $r$: radio de habilitación
del pozo (incluye el filtro de grava). Se obtiene:
\begin{formula}
$\displaystyle\Delta h=\frac{Q}{4\cdot\pi\cdot T}\cdot W(u,\alpha),\qquad
u=\frac{r^2\cdot S}{4\cdot T\cdot t},\qquad\alpha=S\cdot\left(\frac{r_0}{r}\right)^2$
\end{formula}
¡Similar a Theis, se puede utilizar el método de la curva tipo (una curva
por cada valor de $\alpha$)!

\subsection{Método de Hantush--Jacob (1955)}
\begin{itemize}
  \item Acuífero semi-confinado, de pequeño diámetro.
  \item Desprecia el almacenamiento de agua dentro del pozo.
\end{itemize}
Resolvió la siguiente ecuación diferencial (con $s$ el descenso, $K'$ y $m'$
la conductividad y espesor de la capa semiconfinante):
\[
\frac{\partial^2s}{\partial r^2}+\frac1r\frac{\partial s}{\partial r}-\frac{K'}{Tm'}s=\frac ST\frac{\partial s}{\partial t}
\]
obteniendo:
\begin{formula}
$\displaystyle\Delta h=\frac{Q}{4\cdot\pi\cdot T}\cdot W\!\left(u,\frac rB\right),\qquad
u=\frac{r^2\cdot S}{4\cdot T\cdot t},\qquad B=\sqrt{\frac{T\cdot m'}{K'}}$
\end{formula}
¡Similar a Theis, se puede utilizar el método de la curva tipo (familia de
curvas para distintos $r/B$)!

\begin{figure}[H]
\centering
\begin{tikzpicture}[font=\small,scale=0.9]
  \draw[thick] (-4,0) -- (4,0);
  \foreach \x in {-3.9,-3.6,...,3.9} \draw (\x,0) -- (\x-0.2,-0.25);
  \fill[pattern=north east lines] (-4,2.0) rectangle (-0.4,2.4);
  \fill[pattern=north east lines] (0.4,2.0) rectangle (4,2.4);
  \draw (-4,2.0) -- (-0.4,2.0) (0.4,2.0) -- (4,2.0) (-4,2.4) -- (-0.4,2.4) (0.4,2.4) -- (4,2.4);
  \foreach \x in {-3.5,-3.0,...,-1.0,1.0,1.5,...,3.5} \draw[->,rojo] (\x,2.9) -- (\x,2.05);
  \draw[thick] (-0.4,0) -- (-0.4,4.2) (0.4,0) -- (0.4,4.2);
  \draw[thick] (-4,4.2) -- (-0.4,4.2) (0.4,4.2) -- (4,4.2);
  \draw[dashed,cyan!70!blue,thick] (-4,3.7) -- (-0.4,3.7) (0.4,3.7) -- (4,3.7) node[right]{$t=0$};
  \draw[cyan!70!blue,thick] (-4,3.4) .. controls (-2,3.3) and (-1,2.8) .. (-0.4,2.5) -- (0.4,2.5) .. controls (1,2.8) and (2,3.3) .. (4,3.4) node[right]{$t=t$};
  \fill[azulusm] (-0.1,4.1) rectangle (0.1,4.7); \fill[azulusm] (-0.25,4.7) -- (0.25,4.7) -- (0,5.0) -- cycle; \node at (0.5,4.9) {$Q_s$};
  \draw[cota] (-3.3,0) -- (-3.3,2.0) node[midway,left]{$m$};
  \draw[cota] (-4.3,2.0) -- (-4.3,2.4) node[midway,left]{$m'$};
  \node at (-2.2,1.0) {$K,S$}; \node at (-1.5,2.2) [fill=white,inner sep=1pt]{$K'$};
  \node[azulusm,font=\scriptsize] at (2.3,1.7) {Capa confinante};
  \draw[->,azulusm] (1.5,1.0) -- (0.5,1.0) node[pos=0,right]{$Q_i$};
\end{tikzpicture}
\caption{Pozo en acuífero semiconfinado (Hantush--Jacob): filtración vertical a través del acuitardo.}
\end{figure}

\subsection{Acuíferos libres}
Para acuíferos libres se puede usar la linealización del problema:
\begin{itemize}
  \item El método gráfico original no es válido; hay modificaciones (se
        considera el drenaje diferido).
  \item Sí es válido el método de Jacob.
  \item La prueba de recuperación es válida.
\end{itemize}
\begin{formula}
$\displaystyle\Delta h'=\Delta h-\frac{(\Delta h)^2}{2H},\qquad
\Delta h'=\frac{Q}{4\cdot\pi\cdot K\cdot H}\ln\left(\frac{2.24\cdot T\cdot t}{S\cdot r^2}\right)$
\end{formula}
Tener cuidado si existen recargas o barreras impermeables en las cercanías del pozo.

\section{Interpretación de la prueba de recuperación}
\begin{itemize}
  \item Puede modelarse por superposición: a partir de $t_f$ se agrega un
        pozo imagen ficticio que inyecta $-Q$ en el mismo lugar.
  \item Debe respetar el supuesto de Jacob.
  \item Para la zona de recuperación se grafica $\Delta h$ vs.\
        $\ln\frac{t}{t-t_f}$; la parte inicial es una recta.
  \item ¡Similar a Theis, se puede utilizar el método de la curva tipo!
\end{itemize}

\begin{figure}[H]
\centering
\begin{tikzpicture}[font=\small]
  \draw[flecha] (0,0) -- (6,0) node[right]{$t$};
  \draw[flecha] (0,1.2) -- (0,-2.6) node[left]{$\Delta h$};
  \node[left] at (0,1.0) {$Q$};
  \fill[gray!20,draw] (0,0) rectangle (5.8,0.4); \node[left] at (0,0.4) {$+Q$};
  \fill[gray!20,draw] (3.5,0) rectangle (5.8,-0.4); \node[right] at (5.8,-0.4) {$-Q$};
  \draw[orange,very thick,dashed] (0,0) .. controls (0.5,-1.5) and (1.5,-1.9) .. (3.5,-2.1);
  \draw[orange,very thick] (3.5,-2.1) .. controls (3.8,-1.3) and (4.2,-0.9) .. (5.8,-0.5);
  \draw[dashed] (3.5,0.6) node[above]{$t_f$} -- (3.5,-2.3);
\end{tikzpicture}
\caption{La recuperación se modela superponiendo un caudal $-Q$ desde $t_f$.}
\end{figure}

\section{Pruebas de bombeo con condiciones de borde}

\subsection{Barrera impermeable}
Se utiliza el método de las imágenes. En el gráfico $\Delta h$ vs.\ $\ln t$
la pendiente se \textbf{duplica} cuando el pozo imagen comienza a afectar.
\begin{itemize}
  \item Estimar los parámetros del acuífero con la parte no perturbada del gráfico.
  \item Es posible estimar la distancia hasta la barrera.
\end{itemize}

\subsection{Recarga lineal}
Se utiliza el método de las imágenes. En el gráfico $\Delta h$ vs.\ $\ln t$
la curva se \textbf{estabiliza} (pendiente nula) cuando el pozo imagen
comienza a afectar.
\begin{itemize}
  \item Estimar los parámetros del acuífero con la parte no perturbada del gráfico.
  \item Es posible estimar la distancia hasta la fuente de recarga.
\end{itemize}

\begin{figure}[H]
\centering
\begin{subfigure}{0.48\textwidth}\centering
\begin{tikzpicture}[font=\small]
  \draw[flecha] (0,0) -- (5,0) node[right]{$\ln(t)$};
  \draw[flecha] (0,0) -- (0,-3.5) node[left]{$\Delta h$};
  \draw[azulusm,thick] (0,0) -- (2.8,-1.4) -- (4.2,-3.2);
  \draw[azulusm,dashed] (2.8,-1.4) -- (4.8,-2.4);
  \draw[cota] (1.2,0) -- (1.2,-0.6) node[midway,right]{$\Delta h_1$};
  \node[above] at (1.2,0) {$\ln(t_1)$};
  \draw[dashed] (3.9,0) node[above]{$\ln(t_2)$} -- (3.9,-2.9);
  \draw[cota] (3.9,-1.95) -- (3.9,-2.8) node[midway,right]{$\Delta h_2$};
\end{tikzpicture}
\caption{Barrera impermeable.}
\end{subfigure}\hfill
\begin{subfigure}{0.48\textwidth}\centering
\begin{tikzpicture}[font=\small]
  \draw[flecha] (0,0) -- (5,0) node[right]{$\ln(t)$};
  \draw[flecha] (0,0) -- (0,-3.5) node[left]{$\Delta h$};
  \draw[azulusm,thick] (0,0) -- (2.6,-1.3) -- (4.8,-1.3);
  \draw[azulusm,dashed] (2.6,-1.3) -- (4.8,-2.4);
  \draw[cota] (1.2,0) -- (1.2,-0.6) node[midway,right]{$\Delta h_1$};
  \node[above] at (1.2,0) {$\ln(t_1)$};
  \draw[dashed] (4.2,0) node[above]{$\ln(t_2)$} -- (4.2,-2.2);
  \draw[cota] (4.2,-1.3) -- (4.2,-2.1) node[midway,right]{$\Delta h_2$};
\end{tikzpicture}
\caption{Recarga lineal.}
\end{subfigure}
\caption{Efecto de los bordes sobre la curva descenso--tiempo.}
\end{figure}

Si se tiene un pozo de observación a distancia $r_1$ del pozo de bombeo, se
eligen $t_1$ y $t_2$ tales que $\Delta h_1=\Delta h_2$ (desviación de la
recta igual a la depresión de la recta original):
\begin{formula}
$\displaystyle r_2=r_1\sqrt{\frac{t_2}{t_1}}\qquad\qquad d=\frac{r_1+r_2}{2}$
\end{formula}
donde $r_2$ es la distancia del pozo imagen al pozo de observación y $d$ la
distancia estimada al borde.

% =====================================================================
%  CLASE 07 -- Drenaje subterráneo: régimen permanente
% =====================================================================
\clase{07}{Drenaje subterráneo: régimen permanente}

\section{Motivación}
Suelos agrícolas anegados, excavaciones inundadas por la napa o canchas
deportivas saturadas: ¿cómo podemos solucionar este problema?

\section{Drenaje subterráneo}
Los drenes son utilizados con distintos fines:
\begin{itemize}
  \item Disminuir los niveles de agua para evitar la saturación de las raíces
        de una plantación.
  \item Evitar la saturación superficial del suelo de complejos deportivos.
  \item Captación de aguas subterráneas con fines de almacenamiento de agua
        potable y/o superficial.
\end{itemize}
Tipos de obras:
\begin{itemize}
  \item Pozos.
  \item Drenes:
  \begin{itemize}
    \item Drenes verticales o punteras.
    \item Drenes horizontales: naturales, interceptores y paralelos.
  \end{itemize}
\end{itemize}

\section{Drenes horizontales}
\begin{itemize}
  \item Conocidos como drenes o galerías de infiltración.
  \item No son tan populares y usados como los pozos.
  \item Usados en napas superficiales de poco espesor.
  \item Permiten eliminar el consumo energético de bombas.
\end{itemize}

\subsection{Drenes interceptores}
\begin{itemize}
  \item Evacúan sólo una parte del flujo.
  \item $h_0$ es la altura en el canal de evacuación.
  \item Dimensionar $h_0$ tal que el problema se solucione.
  \item $h_0$ puede tomarse como la altura normal en el canal (dren interceptor).
\end{itemize}

\begin{figure}[H]
\centering
\begin{tikzpicture}[font=\small]
  \fill[pattern=crosshatch,pattern color=gray!60] (0,-0.25) rectangle (9,0);
  \draw (0,0) -- (9,0);
  % canal
  \draw[thick] (3,2.4) -- (3,0) -- (4.6,0) -- (4.6,2.4);
  \fill[agua!50] (3,0) rectangle (4.6,0.9);
  \pic at (4.2,0.92) {nivelagua};
  % napa aguas arriba
  \draw[azulusm,thick] (4.6,0.9) .. controls (5.3,1.0) and (5.6,2.0) .. (6.5,2.1) -- (9,2.15);
  \draw[azulusm,thick] (0,0.8) -- (3,0.9);
  \draw[cota] (3.3,0) -- (3.3,0.9) node[midway,right]{$h_0$};
  \draw[cota] (5.6,0) -- (5.6,1.75) node[midway,right]{$h$};
  \draw[cota] (8.2,0) -- (8.2,2.15) node[midway,right]{$H$};
  \draw[cota] (3,2.6) -- (4.6,2.6) node[midway,above]{$b$};
  \draw[->] (4.6,-0.5) -- (5.6,-0.5) node[midway,below]{$x$};
  \draw[->] (4.6,-1.1) -- (8.2,-1.1) node[midway,below]{$x_1$};
  \draw[flecha,azulusm] (7.6,1.2) -- (6.8,1.2) node[pos=0,above]{$q$};
  \draw[flecha,azulusm] (2.2,0.4) -- (1.2,0.4) node[pos=0,above]{$q'$};
\end{tikzpicture}
\caption{Dren interceptor (corte transversal): el canal capta el caudal $q$
y deja pasar $q'$.}
\end{figure}

Aplicando Dupuit al flujo hacia el dren:
\begin{formula}
$\displaystyle h^2-h_0^2=\frac{2\cdot q}{K}\cdot(x-x_0)$\\[4pt]
Si para $x=x_1$, $h=H$ y $x_0=0$:\qquad $\displaystyle H^2-h_0^2=\frac{2\cdot q}{K}\cdot x_1$
\end{formula}
Lo que se lleva el canal, con $L$ el largo del dren:
\begin{formula}
$Q_z=(q-q')\cdot L$
\end{formula}

\subsection{Drenes paralelos}
\paragraph{Clasificación según ubicación en el acuífero}
\begin{itemize}
  \item \textbf{Zanjas abiertas:}
  \begin{itemize}
    \item Líneas de flujo casi horizontales, excepto cerca de la captación.
    \item Equipotenciales planas cercanas a la vertical.
  \end{itemize}
  \item \textbf{Tuberías:}
  \begin{itemize}
    \item Líneas de flujo radiales dirigidas a la captación.
    \item Equipotenciales son superficies semicilíndricas concéntricas con la captación.
  \end{itemize}
\end{itemize}

\begin{figure}[H]
\centering
\begin{subfigure}{0.48\textwidth}\centering
\begin{tikzpicture}[font=\small,scale=0.8]
  \fill[aguaclara] (0,0) rectangle (6,3);
  \draw[thick] (0,3.3) -- (6,3.3);
  \draw[thick] (2.9,0) -- (2.9,3.3) (3.1,0) -- (3.1,3.3);
  \draw[azulusm,thick] (0,2.9) .. controls (2,2.8) and (2.7,2.5) .. (2.9,2.2);
  \draw[azulusm,thick] (6,2.9) .. controls (4,2.8) and (3.3,2.5) .. (3.1,2.2);
  \foreach \y in {0.5,1.2,1.9} {\draw[->,azulusm] (0.3,\y) -- (2.8,\y); \draw[->,azulusm] (5.7,\y) -- (3.2,\y);}
  \foreach \x in {0.8,1.6,2.3,3.7,4.4,5.2} \draw[dashed] (\x,0) -- (\x,2.8);
  \fill[pattern=north east lines] (0,-0.2) rectangle (6,0);
\end{tikzpicture}
\caption{Zanjas abiertas.}
\end{subfigure}\hfill
\begin{subfigure}{0.48\textwidth}\centering
\begin{tikzpicture}[font=\small,scale=0.8]
  \fill[aguaclara] (0,0) rectangle (6,3);
  \draw[thick] (0,3.3) -- (6,3.3);
  \draw[azulusm,thick] (0,2.9) .. controls (2,2.8) and (2.6,2.5) .. (3,2.3) .. controls (3.4,2.5) and (4,2.8) .. (6,2.9);
  \fill[azulusm] (3,2.3) circle (0.18);
  \foreach \a in {0,20,...,180} \draw[->,azulusm] (3,2.3) ++(-\a:2.0) -- ++(180-\a:1.7);
  \foreach \r in {0.8,1.4,2.0} \draw[dashed] (3,2.3) ++(0:\r) arc (0:-180:\r);
  \fill[pattern=north east lines] (0,-0.2) rectangle (6,0);
\end{tikzpicture}
\caption{Tuberías.}
\end{subfigure}
\caption{Redes de flujo hacia drenes paralelos.}
\end{figure}

\paragraph{Régimen de escurrimiento}
\begin{itemize}
  \item Condición de equilibrio (régimen permanente).
  \item Condición de desequilibrio (régimen impermanente).
\end{itemize}

% Macro: drenes paralelos abiertos (#1 = separación al fondo impermeable d)
\newcommand{\drenesabiertos}[1]{%
\begin{tikzpicture}[font=\small,scale=0.9]
  \def\d{#1}
  \fill[pattern=crosshatch,pattern color=gray!60] (-0.5,-\d-0.25) rectangle (8.5,-\d);
  \draw (-0.5,-\d) -- (8.5,-\d);
  % zanjas
  \draw[thick] (-0.5,3) -- (0.6,3) -- (0.6,0) -- (1.0,0) -- (1.0,3) -- (7.0,3) -- (7.0,0) -- (7.4,0) -- (7.4,3) -- (8.5,3);
  \fill[agua!60] (0.6,0) rectangle (1.0,1.0); \fill[agua!60] (7.0,0) rectangle (7.4,1.0);
  % volumen de control
  \fill[rojo,opacity=0.15,draw=rojo] (1.0,0) rectangle (3.3,3);
  \node[rojo] at (1.6,2.6) {\large$V_c$};
  % napa
  \draw[azulusm,thick] (1.0,1.0) .. controls (1.3,2.0) and (2.5,2.3) .. (4.0,2.3) .. controls (5.5,2.3) and (6.7,2.0) .. (7.0,1.0);
  \pic at (0.8,1.02) {nivelagua}; \pic at (7.2,1.02) {nivelagua};
  % lluvia
  \foreach \x in {-0.3,0.1,1.6,2.1,2.6,3.1,3.6,4.1,4.6,5.1,5.6,6.1,7.8,8.2}
    \draw[->,azulusm,thick] (\x,4.0) -- (\x,3.1);
  \node at (0.0,4.3) {$f$}; \node at (3.8,4.3) {$f$}; \node at (8.0,4.3) {$f$};
  \draw[cota] (0.2,0) -- (0.2,1.0) node[midway,left]{$h_0$};
  \draw[cota] (7.8,0) -- (7.8,1.0) node[midway,right]{$h_0$};
  \draw[cota] (4.0,0) -- (4.0,2.3) node[midway,right]{$H$};
  \draw[cota] (1.3,0) -- (1.3,1.45) node[midway,right]{$h$};
  \draw[->] (1.0,0.15) -- (2.0,0.15) node[midway,above]{$x$};
  \draw[flecha,azulusm] (2.4,1.1) -- (1.7,1.1) node[pos=0,right]{$q$};
  \draw[flecha,azulusm] (5.8,1.1) -- (6.5,1.1) node[pos=0,left]{$q$};
  \draw[cota] (0.8,-\d-0.6) -- (7.2,-\d-0.6) node[midway,below]{$D$};
  \ifdim\d pt>0pt \draw[cota] (7.6,-\d) -- (7.6,0) node[midway,right]{$d$};\fi
\end{tikzpicture}}

\section{Drenes abiertos: drenes paralelos en condición de equilibrio}

\subsection{Drenes que abarcan toda la napa: ecuación de Donnan}
\begin{itemize}
  \item Acuífero homogéneo e isotrópico.
  \item Drenes abarcan toda la napa (llegan al fondo impermeable).
  \item Espaciados equidistantemente.
  \item Alimentación constante por unidad de superficie ($f$).
  \item Se desprecia la curvatura del flujo cerca de la zanja.
\end{itemize}

\begin{figure}[H]
\centering
\drenesabiertos{0}
\caption{Drenes paralelos abiertos que abarcan toda la napa.}
\end{figure}

\begin{nota}[title={Deducción}]
Balance en el volumen de control $V_c$ entre el dren ($x=0$) y la sección
$x$, y Dupuit: el caudal que sale por $x=0$ es $q=f\left(\frac D2-x\right)$ y
\[
q=K\,h\frac{\dd h}{\dd x}\ \Rightarrow\ f\left(\frac D2-x\right)\dd x=K\,h\,\dd h
\ \Rightarrow\ \int_0^{D/2}f\left(\frac D2-x\right)\dd x=\int_{h_0}^{H}K\,h\,\dd h
\]
\end{nota}
\begin{formula}
$\displaystyle H^2-h_0^2=\frac{f\cdot D^2}{4\cdot K}$\qquad\textbf{Ecuación de Donnan}
\end{formula}

\subsection{Drenes que no abarcan toda la napa: ecuación de Hooghoudt}
\begin{itemize}
  \item Acuífero homogéneo e isotrópico.
  \item Drenes \textbf{no} abarcan toda la napa.
  \item Espaciados equidistantemente.
  \item Alimentación constante por unidad de superficie.
  \item La distancia al fondo impermeable no debe ser muy grande ($d<0.6$ m).
  \item Se desprecia la curvatura del flujo cerca de la zanja.
\end{itemize}

\begin{figure}[H]
\centering
\drenesabiertos{0.8}
\caption{Drenes paralelos abiertos a una distancia $d$ del estrato impermeable.}
\end{figure}

\begin{formula}
$\displaystyle(H-h_0)(H+h_0+2d)=\frac{f\cdot D^2}{4\cdot K}$\qquad\textbf{Ecuación de Hooghoudt}
\end{formula}

\paragraph{Corrección de Hooghoudt.} Para $d>0.6$ m (flujo radial hacia el
dren) se utiliza una profundidad equivalente $d'$, que se obtiene de un
ábaco en función de la profundidad del estrato $d$ y del espaciamiento $D$
($d'<d$, y $d'$ tiende a un valor límite al crecer $d$):
\begin{formula}
$\displaystyle(H-h_0)(H+h_0+2d')=\frac{f\cdot D^2}{4\cdot K}$
\end{formula}

\section{Drenes cerrados: drenes paralelos en condición de equilibrio}

% Macro: drenes cerrados (tuberías)
\newcommand{\drenescerrados}[1]{% #1: 1 = fondo a distancia d, 0 = infinito
\begin{tikzpicture}[font=\small,scale=0.9]
  \foreach \x in {0,0.4,...,8} \draw[->,azulusm,thick] (\x,3.4) -- (\x,2.6);
  \node at (0.6,3.7) {$f$}; \node at (4,3.7) {$f$}; \node at (7.4,3.7) {$f$};
  \draw[thick] (-0.2,2.5) -- (8.2,2.5);
  \draw[dashed] (-0.5,1.0) -- (8.5,1.0);
  \foreach \x in {0.4,4.0,7.6} \draw[thick,fill=white] (\x,1.0) circle (0.35);
  \draw[azulusm,thick] (0.75,1.1) .. controls (1.2,2.1) and (3.2,2.1) .. (3.65,1.1);
  \draw[azulusm,thick] (4.35,1.1) .. controls (4.8,2.1) and (6.8,2.1) .. (7.25,1.1);
  \draw[cota] (2.2,1.0) -- (2.2,1.85) node[midway,right]{$H$};
  \draw[cota] (0.05,1.6) -- (0.75,1.6) node[midway,above]{$2r_0$};
  \draw[cota] (0.4,0.3) -- (4.0,0.3) node[midway,below]{$D$};
  \draw[cota] (4.0,0.3) -- (7.6,0.3) node[midway,below]{$D$};
  \ifnum#1=1
    \fill[pattern=crosshatch,pattern color=gray!60] (-0.5,-0.9) rectangle (8.5,-0.65);
    \draw (-0.5,-0.65) -- (8.5,-0.65);
    \draw[cota] (8.2,-0.65) -- (8.2,1.0) node[midway,right]{$d$};
    \fill[rojo,opacity=0.15,draw=rojo] (4.0,-0.65) rectangle (5.8,2.5);
    \node[rojo] at (4.9,2.2) {$V_c$};
  \else
    \fill[rojo,opacity=0.15,draw=rojo] (2.2,2.5) -- (2.2,1.0) arc (180:360:1.8) -- (5.8,2.5) -- cycle;
    \node[rojo] at (4.0,2.2) {$V_c$};
    \foreach \a in {200,230,...,340} \draw[->,azulusm] (4,1.0) ++(\a:1.7) -- ++(\a+180:1.2);
    \node at (7.5,-0.7) {$d\approx\infty$};
  \fi
\end{tikzpicture}}

\subsection{Corrección de Dagan (1964)}
\begin{itemize}
  \item Para distancias ($d$) al estrato impermeable pequeñas.
  \item Incorpora la curvatura cerca del dren.
  \item Se aplica un factor de corrección a la fórmula de Donnan.
  \item También corrige la fórmula de Donnan ($d=h_0$).
\end{itemize}

\begin{figure}[H]
\centering
\drenescerrados{1}
\caption{Drenes paralelos cerrados (tuberías) a una distancia $d$ del estrato impermeable.}
\end{figure}

\begin{formula}
$\displaystyle(H+d)^2-d^2=\frac{f\cdot D^2}{4\cdot K}\left(1+\alpha\frac{4\cdot d}{D}\right)$\\[8pt]
$\displaystyle\alpha=-\frac1\pi\ln\left[2\left(\cosh\left(\frac{\pi\cdot r_0}{d}\right)-1\right)\right]$
\end{formula}

\begin{table}[H]
\centering
\begin{tabular}{cc}
\toprule
$r_0/d$ & $\alpha$\\
\midrule
0.01 & 2.16\\
0.02 & 1.75\\
0.04 & 1.31\\
0.06 & 1.05\\
0.08 & 0.86\\
0.10 & 0.75\\
0.20 & 0.29\\
0.30 & 0.01\\
\bottomrule
\end{tabular}
\caption{Factor $\alpha$ de Dagan. Para $r_0/d\geq0.3\Rightarrow\alpha\to0$.}
\end{table}

\subsection{Drenes superficiales (acuífero de espesor infinito)}
\begin{itemize}
  \item Acuífero homogéneo e isotrópico.
  \item Acuífero de espesor muy grande o indefinido (infinito).
  \item Espaciados equidistantemente.
  \item Alimentación constante por unidad de superficie.
  \item Escurrimiento semi-cilíndrico.
\end{itemize}

\begin{figure}[H]
\centering
\drenescerrados{0}
\caption{Drenes superficiales en acuífero de espesor infinito: escurrimiento semicilíndrico.}
\end{figure}

\begin{formula}
$\displaystyle H=\frac{D\cdot f}{K\cdot\pi}\ln\left(\frac{D}{5.44\cdot r_0}\right)$
\end{formula}

\subsection{Acuíferos de espesor finito: fórmula de Kirkham (1958)}
Incorpora la curvatura cerca del dren:
\begin{formula}
$\displaystyle H=\frac{D\cdot f}{K}\cdot F\left(\frac{2r_0}{D},\frac dD\right)$\\[8pt]
$\displaystyle F(x,y)=\frac1\pi\left[\ln\left(\frac{2}{\pi x}\right)+\sum_{m=1}^\infty\frac1m\big(\cos(m\pi x)-\cos(m\pi)\big)\big(\coth(2m\pi y)-1\big)\right]$
\end{formula}

\begin{table}[H]
\centering
\begin{tabular}{c|ccccc}
\toprule
& \multicolumn{5}{c}{$2r_0/D$}\\
$d/D$ & 0.0025 & 0.005 & 0.01 & 0.02 & 0.04\\
\midrule
0.01 & 12.79 & 12.57 & 12.33 & 12.03 & 11.52\\
0.02 & 6.761 & 6.541 & 6.318 & 6.077 & 5.771\\
0.04 & 3.864 & 3.643 & 3.421 & 3.195 & 2.954\\
0.08 & 2.522 & 2.301 & 2.080 & 1.858 & 1.633\\
0.16 & 1.961 & 1.741 & 1.520 & 1.299 & 1.077\\
0.32 & 1.787 & 1.566 & 1.345 & 1.125 & 0.904\\
0.64 & 1.764 & 1.543 & 1.323 & 1.102 & 0.811\\
1.00 & 1.763 & 1.543 & 1.322 & 1.101 & 0.811\\
$\infty$ & 1.763 & 1.543 & 1.322 & 1.101 & 0.811\\
\bottomrule
\end{tabular}
\caption{Valores de $F(2r_0/D,\,d/D)$ de Kirkham.}
\end{table}

\section{Régimen permanente: consideraciones de diseño}
La infiltración siempre dependerá de:
\begin{itemize}
  \item Nivel de saturación del suelo.
  \item Intensidad de precipitación.
  \item Propiedades del suelo.
  \item Recordar métodos de infiltración de hidrología para valores precisos.
\end{itemize}
\begin{table}[H]
\centering
\begin{tabular}{lc}
\toprule
\textbf{Tipo de suelo} & $f$ [mm/h]\\
\midrule
Arena & $>30$\\
Arena limosa & 20 -- 30\\
Limo & 10 -- 20\\
Arcilla limosa & 5 -- 10\\
Arcilla & 1 -- 5\\
\bottomrule
\end{tabular}
\caption{Ejemplos de valores de infiltración.}
\end{table}
El diseño en régimen de equilibrio o permanente se utiliza cuando no es
permisible tener un nivel de saturación prolongado en el tiempo. Ejemplo:
plantas con raíces que no permiten estar saturadas por un tiempo prolongado,
campos deportivos, etc.

% =====================================================================
%  CLASE 08 -- Drenaje subterráneo: régimen impermanente
% =====================================================================
\clase{08}{Drenaje subterráneo: régimen impermanente}

\section{Régimen impermanente}
\paragraph{Hipótesis}
\begin{itemize}
  \item Lluvia o sobrecarga esporádica.
  \item Acuífero homogéneo e isotrópico.
\end{itemize}
\paragraph{Condición de diseño}
\begin{itemize}
  \item La saturación no debe llegar a las raíces del cultivo.
\end{itemize}

\begin{figure}[H]
\centering
\begin{tikzpicture}[font=\small,scale=1.0]
  \fill[pattern=crosshatch,pattern color=gray!60] (-0.5,-0.25) rectangle (8.5,0);
  \draw (-0.5,0) -- (8.5,0);
  \draw[thick] (-0.5,3.2) -- (0.8,3.2) -- (0.8,0) (1.3,0) -- (1.3,3.2) -- (6.7,3.2) -- (6.7,0) (7.2,0) -- (7.2,3.2) -- (8.5,3.2);
  \draw[dashed] (0.8,1.0) -- (7.2,1.0);
  \pic at (1.05,1.02) {nivelagua}; \pic at (6.95,1.02) {nivelagua};
  \draw[azulusm,thick] (1.3,1.0) .. controls (1.8,3.1) and (6.2,3.1) .. (6.7,1.0);
  \draw[azulusm,thick,dashed] (1.3,1.0) .. controls (1.8,2.3) and (6.2,2.3) .. (6.7,1.0);
  \draw[cota] (0.2,0) -- (0.2,3.2) node[midway,left]{$H_0$};
  \draw[cota] (7.6,0) -- (7.6,1.0) node[midway,right]{$h_0$};
  \draw (2.6,0) -- (2.6,2.45) (3.0,0) -- (3.0,2.6);
  \node at (2.4,2.7) {$h$}; \node at (3.4,2.9) {$h+\frac{\partial h}{\partial x}\partial x$};
  \draw[flecha,azulusm] (2.5,0.6) -- (1.8,0.6) node[pos=0,above right]{$Q_2$};
  \draw[flecha,azulusm] (3.6,0.6) -- (3.1,0.6) node[pos=0,above]{$Q_1$};
  \node[below] at (2.4,0) {$x$}; \node[below] at (3.5,0) {$x+\partial x$};
  \draw[cota] (1.3,-0.8) -- (4.0,-0.8) node[midway,below]{$D/2$};
  \draw[dashed] (4,0) -- (4,3.2);
  \draw[<-] (2.0,2.4) -- (1.2,4.0) node[above]{Nivel inicial};
  \draw[<-] (5.4,2.0) -- (6.8,4.0) node[above]{Nivel instante $t$};
  \node[above] at (4,3.4) {Parábola de 4° grado};
\end{tikzpicture}
\caption{Descenso de la napa entre drenes paralelos después de una recarga esporádica.}
\end{figure}

Usando la aproximación de Dupuit--Forchheimer ($\partial h/\partial s\approx\partial h/\partial r$),
la ecuación de continuidad es:
\begin{formula}
$\displaystyle h\frac{\partial^2h}{\partial x^2}+\left(\frac{\partial h}{\partial x}\right)^2=\frac SK\frac{\partial h}{\partial t}$
\qquad\textbf{Ecuación de continuidad}
\end{formula}
\paragraph{Condiciones de borde}
\begin{itemize}
  \item Suelo saturado: \quad $t=0\ \Rightarrow\ x=\frac D2;\ h=H_0$.
  \item Napa vuelve a su estado natural: \quad $t=\infty\ \Rightarrow\ x=\frac D2;\ h=h_0$.
\end{itemize}

% Macro drenes cerrados con y, y0, H0, d
\newcommand{\drenesimperm}{%
\begin{tikzpicture}[font=\small,scale=0.85]
  \fill[pattern=crosshatch,pattern color=gray!60] (-0.5,-0.25) rectangle (8.5,0);
  \draw (-0.5,0) -- (8.5,0);
  \draw[thick] (-0.5,3.8) -- (8.5,3.8);
  \draw[dashed] (-0.5,1.8) -- (8.5,1.8);
  \foreach \x in {0.4,4.0,7.6} \draw[thick,fill=white] (\x,1.8) circle (0.35);
  \draw[azulusm,thick] (0.75,1.9) .. controls (1.2,3.0) and (3.2,3.0) .. (3.65,1.9);
  \draw[azulusm,thick] (4.35,1.9) .. controls (4.8,3.0) and (6.8,3.0) .. (7.25,1.9);
  \draw[cota] (-0.2,0) -- (-0.2,3.8) node[midway,left]{$H_0$};
  \draw[cota] (0.2,1.8) -- (0.2,3.8) node[midway,right]{$y_0$};
  \draw[cota] (2.2,0) -- (2.2,2.72) node[pos=0.4,right]{$h$};
  \draw[cota] (5.8,2.72) -- (5.8,3.8) node[midway,right]{$y$};
  \draw[cota] (3.65,2.4) -- (4.35,2.4) node[midway,above]{$2r_0$};
  \draw[cota] (0.4,1.0) -- (4.0,1.0) node[midway,below]{$D$};
  \draw[cota] (4.0,1.0) -- (7.6,1.0) node[midway,below]{$D$};
  \draw[cota] (8.2,0) -- (8.2,1.8) node[midway,right]{$d$};
\end{tikzpicture}}

\section{Solución gráfica de Moody (1964)}
Resolvió la ecuación anterior mediante diferencias finitas y con los
resultados construyó gráficos adimensionales de $y/y_0$ y de
$\dfrac{qD}{K\pi y_0}$ en función de $\alpha$, para distintos valores de $m$.

\begin{figure}[H]
\centering
\drenesimperm
\caption{Notación para drenes cerrados en régimen impermanente.}
\end{figure}

Cambio de variable:
\[
y=h-d,\qquad y_0=H_0-d,\qquad m=\frac{y_0}{H_0},\qquad\alpha=\frac{KH_0t}{SD^2}
\]
Gráficas válidas para $d\neq0$.

\begin{figure}[H]
\centering
\begin{subfigure}{0.45\textwidth}\centering
\begin{tikzpicture}[font=\small]
  \draw[flecha] (0,0) -- (4.5,0) node[right]{$\alpha$};
  \draw[flecha] (0,0) -- (0,3.3) node[left]{$y/y_0$};
  \foreach \s in {0,0.35,0.7}
    \draw[azulusm,thick] (0,3) .. controls (1.2+\s,3) and (1.3+\s,0.2) .. (4.2,0.05);
  \draw[<-] (2.1,1.5) -- (2.7,2.2) node[above right]{$m$};
\end{tikzpicture}
\end{subfigure}
\begin{subfigure}{0.45\textwidth}\centering
\begin{tikzpicture}[font=\small]
  \draw[flecha] (0,0) -- (4.5,0) node[right]{$\alpha$};
  \draw[flecha] (0,0) -- (0,3.3) node[left]{$\dfrac{qD}{K\pi y_0}$};
  \foreach \a/\b in {2.8/0.9,2.1/0.6,1.5/0.4}
    \draw[azulusm,thick] (0.3,\a) .. controls (2,\a-0.5) and (3,\b+0.3) .. (4.2,\b);
  \draw[<-] (2.8,1.6) -- (3.3,2.2) node[above right]{$m$};
\end{tikzpicture}
\end{subfigure}
\caption{Esquema de las curvas de Moody.}
\end{figure}

\section{Solución de Glover--Dumm (1954)}
\begin{itemize}
  \item El suelo parte totalmente saturado (condición inicial).
  \item Entrega la separación entre drenes para que el nivel baje de $y_0$ a
        $y$ en un tiempo $t$.
\end{itemize}
Con el cambio de variable $y=h-d$, $y_0=H_0-d$:
\begin{formula}
$\displaystyle y(x,t)=\frac{4y_0}{\pi}\sum_{n=1}^\infty\frac{1}{2n-1}\sin\left(\frac{(2n-1)\pi x}{D}\right)\exp\left(-(2n-1)^2\alpha t\right),
\qquad\alpha=\frac{\pi^2Kd}{SD^2}$
\end{formula}
Para $x=D/2$ y $\alpha t\geq0.2$ basta el primer término:
\begin{formula}
$y=1.16\,y_0\,e^{-\alpha t}\quad\Longrightarrow\quad
\displaystyle D=\pi\sqrt{\frac{K\cdot d\cdot t}{S\cdot\ln\left(\frac{1.16\cdot y_0}{y}\right)}}$
\qquad válida para $d\neq0$
\end{formula}

\section{Corrección de Hooghoudt}
Al incluir simplificaciones, las soluciones anteriores producen algunos
errores en la estimación de $D$. Se reemplaza $d$ por una profundidad
equivalente $d'$:
\begin{formula}
Si $\dfrac dD<0.25$:\quad $\displaystyle\frac{d'}{d}=\frac{1}{1+\frac8\pi\frac dD\ln\left(\frac d\chi\right)}$\\[10pt]
Si $\dfrac dD\geq0.25$:\quad $\displaystyle\frac{d'}{d}=\frac{\pi}{2\left(\ln\left(\frac D\chi\right)+0.18\right)}$\\[8pt]
$\chi=\pi r_0$: perímetro mojado del dren
\end{formula}
Procedimiento iterativo:
\begin{enumerate}
  \item Partir con $d$ y obtener $D$.
  \item Con $D$ se corrige $d$ (se obtiene $d'$).
  \item Se calcula nuevamente $D$.
  \item Iterar hasta que $D$ converja.
\end{enumerate}

\section{Diseño en equilibrio dinámico}
Si se conoce el ciclo anual de recargas (infiltración), por ejemplo un
histograma mensual de recarga $R$ [mm], el sistema se encuentra en su
\textbf{equilibrio dinámico} si se cumple:
\begin{formula}
$y_f\leq y_0$
\end{formula}
con $y_f$: altura de agua en el último mes de medición; $y_0$: altura de agua
en el primer mes de medición. Además se exige un \textbf{índice de saturación}:
\begin{formula}
$\displaystyle I_s=\int(y>y_a)\,\dd t\leq200\ [\text{mm}\cdot\text{día}]$
\end{formula}
donde $Z_a$ es la profundidad de raíces ($y_a$ la altura correspondiente).

\begin{figure}[H]
\centering
\begin{tikzpicture}[font=\small]
  \draw[flecha] (0,0) -- (9,0) node[right]{$t$};
  \draw[flecha] (0,0) -- (0,3.2) node[left]{$y$};
  \draw[rojo,dashed,thick] (0,2.0) node[left,black]{$y_a$} -- (8.8,2.0);
  \draw[azulusm,thick] (0,1.6)
    \foreach \i in {0,...,9}{ -- ({0.3+0.85*\i},{2.8-0.08*\i}) .. controls ({0.5+0.85*\i},{1.6-0.05*\i}) .. ({0.85*(\i+1)},{1.4-0.05*\i}) };
  \node[right] at (8.6,0.9) {$y_f$};
  \node[above] at (1.5,2.9) {$I_s$: área sobre $y_a$};
\end{tikzpicture}
\caption{Evolución de la napa con recargas sucesivas; el área sobre $y_a$ define el índice de saturación.}
\end{figure}

Para cada mes se debe calcular la altura de agua al inicio y final del mes:
\begin{itemize}
  \item Al inicio del mes $i$ se tendrá una altura de agua:
  \[
  y_{0,i}=y_{f,i-1}+\Delta y_i,\qquad\Delta y_i=\frac{R_i}{S}\quad(S:\ \text{coeficiente de almacenamiento})
  \]
  \item Al final del mes $i$ se tendrá una altura de agua:
  \[
  y_{f,i}=1.16\,y_{0,i}\,e^{-\alpha t_i},\qquad\alpha=\frac{\pi^2Kd}{SD^2}
  \]
\end{itemize}


\part{Drenaje superficial}
% =====================================================================
%  CLASE 09 -- Drenaje longitudinal de caminos y carreteras
% =====================================================================
\clase{09}{Drenaje longitudinal de caminos y carreteras}

\section{Drenaje superficial de caminos y carreteras}
Se debe diseñar un correcto drenaje superficial del agua que fluye por los
caminos y carreteras porque:
\begin{itemize}
  \item El agua afecta al pavimento provocando problemas de socavación y
        otros de mecánica de suelos.
  \item No se deben interrumpir cursos de agua existentes, como quebradas.
\end{itemize}
\paragraph{Lectura de referencia:} Manual de Carreteras (MOP), Volumen N°~3
\emph{Instrucciones y Criterios de Diseño} (Capítulo 3.700: Diseño del
drenaje, saneamiento, mecánica e hidráulica fluvial) y Volumen N°~2
\emph{Procedimientos de Estudios Viales} (Capítulo 2.400), disponibles en
\url{https://mc.mop.gob.cl/}.

Existen dos tipos de drenaje:
\begin{itemize}
  \item Drenaje longitudinal.
  \item Drenaje transversal.
\end{itemize}

\section{Drenaje longitudinal}
Las obras hidráulicas son el \textbf{contrafoso} (a media ladera, intercepta
el agua que escurre por la ladera) y la \textbf{cuneta} (junto a la calzada).

\begin{figure}[H]
\centering
\begin{tikzpicture}[font=\small]
  % ladera en planta
  \fill[suelo!80] (0,0) rectangle (7,4);
  \fill[aguaclara] (4.2,0) rectangle (4.6,4);
  \fill[aguaclara] (6.9,0) rectangle (7.05,4);
  \fill[gray!50] (7.05,0) rectangle (8.2,4);
  \foreach \y in {0.4,1.2,2.0,2.8,3.6} \draw[white,very thick] (7.62,\y) -- (7.62,\y+0.4);
  \foreach \x/\l in {0.8/h_7,1.7/h_6,2.6/h_5,3.5/h_4,5.3/h_3,6.2/h_2}
    \draw (\x,0) .. controls (\x+0.2,1.3) and (\x-0.2,2.7) .. (\x+0.1,4) node[pos=0,below]{$\l$};
  \node[below] at (6.9,0) {$h_1$};
  \draw (0.5,2.0) ellipse (0.3 and 0.8); \node at (0.5,2.0) {$h_9$};
  \foreach \x in {1.1,2.0,2.9,3.7,4.9,5.8} \draw[-{Stealth},cyan!60!blue,line width=2pt] (\x,2) -- (\x+0.5,2);
  \draw[flecha,azulusm] (4.4,3) -- (4.4,1);
  \draw[flecha,azulusm] (6.97,3) -- (6.97,1);
  \node[above,azulusm] at (1.5,4.1) {Curvas de nivel};
  \node[above,azulusm] at (4.4,4.1) {Contrafoso};
  \node[above,azulusm] at (7.1,4.1) {Cuneta};
  \node[azulusm,right] at (8.3,0.2) {Vista en planta};
\end{tikzpicture}
\caption{Contrafoso y cuneta en vista en planta: el agua de la ladera es
interceptada antes de llegar a la calzada.}
\end{figure}

\paragraph{Objetivo:} evitar que las aguas que escurren por una ladera
lleguen a la plataforma vial (carretera). El agua que escurre por el
contrafoso y la cuneta se descarga a una quebrada, río, lago o mar.

\begin{figure}[H]
\centering
\begin{tikzpicture}[font=\small]
  \fill[left color=suelo!90!black,right color=suelo!60] (0,0) -- (0,4) -- (1,3.8) -- (1.8,3.5) -- (2.4,3.3)
     -- (2.4,2.6) -- (3.0,2.6) -- (3.1,3.2) -- (3.6,3.0) -- (4.2,2.4) -- (4.8,1.6) -- (5.3,0.8) -- (5.6,0.3) -- (5.6,0) -- cycle;
  \fill[gray!50] (5.6,0) -- (5.6,0.4) -- (5.75,0.4) -- (6.0,0.05) -- (7.6,0.35) -- (8.6,0.25) -- (8.6,0) -- cycle;
  \draw[dashed] (6.0,0.05) -- (8.4,0.05);
  \node at (7.3,0.15) {\scriptsize $i_2$};
  \draw[<-] (2.7,2.6) -- (3.3,3.8) node[above,azulusm]{Contrafoso};
  \draw[<-] (5.85,0.15) -- (5.6,1.3) node[above,azulusm]{Cuneta};
  \draw[<-] (7.6,0.4) -- (8.0,1.1) node[above,azulusm]{Bombeo};
  \node[azulusm] at (6.5,4.5) {Perfil transversal};
\end{tikzpicture}
\caption{Perfil transversal de un camino a media ladera.}
\end{figure}

\subsection{Recomendaciones de diseño de la cuneta}
\begin{itemize}
  \item Ancho: $w\geq0.5$ [m].
  \item Pendiente transversal:
  \begin{itemize}
    \item $i_1\geq8\%$ para cunetas con ancho superior a 0.5 [m].
    \item 30\% máximo para cunetas de ancho de 0.5 [m], para aprovechar de
          mejor forma la capacidad de la cuneta y la eficiencia de los sumideros.
  \end{itemize}
  \item Pendiente longitudinal mínima de la cuneta:
  \begin{itemize}
    \item 0.12\% para cunetas revestidas.
    \item 0.25\% para cunetas sin revestir.
  \end{itemize}
  \item Bombeo o pendiente de la calzada: $i_2\approx2\sim4\%$.
\end{itemize}

\begin{figure}[H]
\centering
\begin{tikzpicture}[font=\small]
  \fill[gray!50] (0,0) rectangle (0.5,2);
  \fill[gray!40] (0.5,0) -- (0.5,0.1) -- (2.2,1.0) -- (6.5,1.5) -- (7.5,1.35) -- (7.5,0) -- cycle;
  \draw (0.5,0.1) -- (2.2,1.0) -- (6.5,1.5) -- (7.5,1.35);
  \draw[dashed] (2.2,0) -- (2.2,2.2);
  \draw[cota] (0.5,1.8) -- (2.2,1.8) node[midway,above]{$w$};
  \draw[dashed] (2.2,1.0) -- (6.8,1.0);
  \draw (1.3,0.1) -- (1.8,0.1); \draw (1.3,0.1) arc (0:28:0.8); \node at (1.5,0.4) {$i_1$};
  \draw (5.8,1.0) arc (180:174:1.2); \node at (5.4,1.15) {$i_2$};
  \draw[<-] (6.5,1.55) -- (6.9,2.1) node[above,azulusm]{Bombeo};
\end{tikzpicture}
\caption{Geometría de la cuneta.}
\end{figure}

Para el correcto diseño del contrafoso y cuneta se debe considerar el cálculo de:
\begin{itemize}
  \item Caudal de diseño.
  \item Altura de agua (velocidad de escurrimiento).
  \item Periodo de retorno.
\end{itemize}

\section{Caudal máximo}
Cuando no se cuenta con registros de caudales, se deben usar fórmulas de
precipitación--escorrentía. En general se trabaja con cuencas pequeñas
(área $<10$ [km$^2$]), lo que se traduce en tiempos de concentración cortos:
\begin{center}
\colorbox{azulusm}{\parbox{0.5\textwidth}{\centering\color{white}\large Duración tormenta $(t_d)>t_c$}}
\end{center}
Por lo tanto, es válido el uso de la \textbf{fórmula racional} (área $<25$ [km$^2$]):
\begin{formula}
$\displaystyle Q_{max}=\frac{C\cdot\bar\imath(t_c)\cdot A}{3.6}\ \ [\text{m}^3/\text{s}]$\\[4pt]
$\bar\imath(t_c)$ en mm/h;\quad $A$: área aportante pluvial en km$^2$
\end{formula}
Para el cálculo del tiempo de concentración se pueden utilizar las fórmulas
vistas en hidrología, pero se debe verificar $t_c\geq10$ [min].

\section{Tiempo de concentración}
Recomendaciones del Volumen 3, Manual de Carreteras:
\begin{table}[H]
\centering\small
\resizebox{\linewidth}{!}{%
\begin{tabular}{lll}
\toprule
\textbf{Autor} & \textbf{Expresión} & \textbf{Observaciones}\\
\midrule
Normas Españolas & $T_c=18\,L^{0.76}/S^{0.19}$ & \\
California Culverts Practice (1942) & $T_c=57\,(L^3/H)^{0.385}$ & Cuencas de montaña\\
Giandotti & $T_c=60\,\dfrac{4A^{0.5}+1.5L}{0.8\,H_m^{0.5}}$ & Cuencas pequeñas con pendiente\\
SCS (1975) & $T_c=3.42\,L^{0.8}\left(\frac{1000}{CN}-9\right)^{0.7}/S^{0.5}$ & Cuencas rurales\\
\bottomrule
\end{tabular}}
\caption{Tabla 3.702.501.A: tiempos de concentración para cuencas.}
\end{table}
Notación: $T_c$: tiempo de concentración (min); $L$: longitud del cauce
(km); $S$: pendiente (\%); $A$: área de la cuenca (km$^2$); $H_m$: diferencia
de nivel (m) entre la cota media de la cuenca y la salida; $H$: diferencia de
nivel total entre cotas extremas de la cuenca (m); $CN$: curva número.

\begin{table}[H]
\centering\small
\resizebox{\linewidth}{!}{%
\begin{tabular}{lll}
\toprule
\textbf{Autor} & \textbf{Expresión} & \textbf{Observaciones}\\
\midrule
Federal Aviation Agency (1970) & $T_c=3.26\,(1.1-C)\,L^{0.5}/(100S)^{0.33}$ & Aeropuertos\\
Izzard (1946) & $T_c=525.28\,(0.0000276\,i+C)\,L_s^{0.33}/(i^{0.667}S^{0.333})$ & Experimentos de laboratorio\\
Morgali y Linsley (1965) & $T_c=7\,L_s^{0.6}\,n^{0.6}/(i^{0.4}S^{0.3})$ & Flujo superficial\\
\bottomrule
\end{tabular}}
\caption{Tabla 3.702.501.B: tiempos de concentración para áreas planas.}
\end{table}
Notación: $L_s$: longitud de escurrimiento superficial (m); $i$: intensidad
de lluvia (mm/h); $C$: coeficiente de escurrimiento; $n$: rugosidad
superficial de Manning. ¡Revisar clases de hidrología!

\section{Precipitación de diseño}
\begin{itemize}
  \item Si se tienen datos horarios se pueden obtener las curvas IDF.
  \item Si sólo se cuenta con registros diarios, se pueden utilizar:
  \begin{itemize}
    \item La fórmula de Grunsky:
    \begin{formula}
    $\displaystyle i_t=i_{24}\sqrt{\frac{24}{t}}\ \ \left[\frac{\text{mm}}{\text{h}}\right]$
    \end{formula}
    \item Coeficientes de duración y frecuencia (V3, Manual de Carreteras):
    \begin{formula}
    $P_t^T=K\cdot CD_t\cdot CF_T\cdot P_D^{10}$
    \end{formula}
  \end{itemize}
\end{itemize}
donde $P_t^T$: lluvia con período de retorno de $T$ años y duración $t$
horas; $P_D^{10}$: lluvia diaria (8 AM a 8 AM) con 10 años de período de
retorno obtenida de una estación pluviométrica; $CD_t$: coeficiente de
duración para $t$ horas; $CF_T$: coeficiente de frecuencia para $T$ años de
período de retorno; $K$: coeficiente de corrección para la lluvia máxima
$P_D^{10}$ medida entre 8 AM y 8 AM respecto de las 24 h más lluviosas de la
tormenta, para el que se ha adoptado un valor $K=1.1$.

Para duraciones inferiores a una hora:
\begin{itemize}
  \item La fórmula de Grunsky.
  \item La fórmula de Bell (1969, V3 Manual de Carreteras):
  \begin{formula}
  $P_t^T=\left(0.54\,t^{0.25}-0.50\right)\left(0.21\ln T+0.52\right)P_1^{10}$
  \end{formula}
  con $P_t^T$: precipitación (mm) con período de retorno $T$ años y duración
  $t$ minutos; $P_1^{10}$: precipitación (mm) con $T=10$ años y duración de 1 hora.
\end{itemize}

\section{Coeficiente de escorrentía}
\begin{table}[H]
\centering
\begin{tabular}{lc}
\toprule
\textbf{Tipo de terreno} & \textbf{Coeficiente de escurrimiento}\\
\midrule
Pavimentos de adoquín & 0.50 -- 0.70\\
Pavimentos asfálticos & 0.70 -- 0.95\\
Pavimentos en concreto & 0.80 -- 0.95\\
Suelo arenoso con vegetación y pendiente 2\% -- 7\% & 0.15 -- 0.20\\
Suelo arcilloso con pasto y pendiente 2\% -- 7\% & 0.25 -- 0.65\\
Zonas de cultivo & 0.20 -- 0.40\\
\bottomrule
\end{tabular}
\caption{Tabla 3.702.503.A: coeficientes de escurrimiento $C$.}
\end{table}

\begin{table}[H]
\centering\scriptsize
\begin{tabularx}{\textwidth}{>{\raggedright}p{1.8cm}XXXX}
\toprule
\textbf{Factor} & \textbf{Extremo} & \textbf{Alto} & \textbf{Normal} & \textbf{Bajo}\\
\midrule
Relieve & 0.28--0.35: escarpado con pendientes mayores que 30\% & 0.20--0.28: montañoso con pendientes entre 10 y 30\% & 0.14--0.20: con cerros y pendientes entre 5 y 10\% & 0.08--0.14: relativamente plano con pendientes menores al 5\%\\
\addlinespace
Infiltración & 0.12--0.16: suelo rocoso, o arcilloso con capacidad de infiltración despreciable & 0.08--0.12: suelos arcillosos o limosos con baja capacidad de infiltración, mal drenados & 0.06--0.08: normales, bien drenados, textura mediana, limos arenosos, suelos arenosos & 0.04--0.06: suelos profundos de arena u otros suelos bien drenados con alta capacidad de infiltración\\
\addlinespace
Cobertura vegetal & 0.12--0.16: cobertura escasa, terreno sin vegetación o escasa cobertura & 0.08--0.12: poca vegetación, terrenos cultivados o naturales, menos del 20\% del área con buena cobertura vegetal & 0.06--0.08: regular a buena; 50\% del área con praderas o bosques, no más del 50\% cultivado & 0.04--0.06: buena a excelente; 90\% del área con praderas, bosques o cobertura equivalente\\
\addlinespace
Almacenamiento superficial & 0.10--0.12: despreciable, pocas depresiones superficiales, sin zonas húmedas & 0.08--0.10: baja, sistema de cauces superficiales pequeños bien definidos, sin zonas húmedas & 0.06--0.08: normal; posibilidad de almacenamiento buena, zonas húmedas, pantanos, lagunas y lagos & 0.04--0.06: capacidad alta, sistema hidrográfico poco definido, buenas planicies de inundación o gran cantidad de zonas húmedas, lagunas o pantanos\\
\bottomrule
\end{tabularx}
\caption{Tabla 3.702.503.B: coeficientes de escorrentía $C$ para $T=10$ años
(se suman los aportes de cada factor). Si $T>10$ años, amplificar el
resultado por: $T=25$: $C\times1.10$; $T=50$: $C\times1.20$; $T=100$: $C\times1.25$.}
\end{table}

\section{Altura de agua $h$}
Por tratarse de flujo a superficie libre, con el ancho o algún otro dato
geométrico dado se utiliza la fórmula de Manning:
\begin{formula}
$\displaystyle\frac{Q\cdot n}{\sqrt s}=\frac{\Omega^{5/3}}{\Psi^{2/3}}$\qquad
$\Omega$: área de escurrimiento;\quad $\Psi$: perímetro mojado
\end{formula}
\begin{itemize}
  \item Generalmente los contrafosos están construidos de tierra: $n=0.023\sim0.025$.
  \item Se espera que el escurrimiento sea subcrítico; sin embargo, se debe
        verificar con la altura crítica.
\end{itemize}

\subsection{Rugosidad de Manning}
Se pueden usar rugosidades características, dependiendo del tipo de
revestimiento del canal (extracto de la Tabla 3.705.1.A, valores mínimo,
normal y máximo):
\begin{table}[H]
\centering\small
\begin{tabular}{llccc}
\toprule
\textbf{Revestimiento} & \textbf{Descripción} & \textbf{Mín.} & \textbf{Normal} & \textbf{Máx.}\\
\midrule
Hormigón & Platachado & 0.011 & 0.013 & 0.015\\
 & Alisado con regla & 0.013 & 0.015 & 0.016\\
 & Alisado con ripio a la vista en el fondo & 0.015 & 0.017 & 0.020\\
 & Sin alisar & 0.014 & 0.017 & 0.020\\
 & Gunita, sección regular & 0.016 & 0.019 & 0.023\\
 & Gunita, sección ondulada & 0.018 & 0.022 & 0.025\\
Fondo de hormigón alisado & Piedra acomodada en mortero & 0.015 & 0.017 & 0.020\\
\quad con lados de: & Piedra distribuida al azar en mortero & 0.017 & 0.020 & 0.024\\
 & Empedrado o enrocado (rip rap) & 0.020 & 0.030 & 0.035\\
Fondo de grava con lados de: & Hormigón (con moldaje) & 0.017 & 0.020 & 0.025\\
 & Empedrado o rip rap & 0.023 & 0.033 & 0.036\\
Ladrillo & Barnizado o vidriado & 0.011 & 0.013 & 0.015\\
 & En mortero de cemento & 0.012 & 0.015 & 0.018\\
Albañilería & Empedrado cementado & 0.017 & 0.025 & 0.030\\
 & Empedrado libre & 0.023 & 0.032 & 0.035\\
Asfalto & Liso & 0.013 & 0.013 & --\\
 & Rugoso & 0.016 & 0.016 & --\\
Cubierto con vegetación & & 0.030 & -- & 0.500\\
\bottomrule
\end{tabular}
\caption{Coeficientes de rugosidad de Manning en canales.}
\end{table}

\subsection{Cálculo de altura normal y crítica}
Utilizar lo aprendido en Hidráulica Teórica. Se puede ayudar de software
(p.~ej.\ HCANALES) para determinar rápidamente $h_n$ y $h_c$.

\section{Velocidades de escurrimiento}
\paragraph{Velocidades máximas} (evitar problemas de socavación):
\begin{table}[H]
\centering
\begin{tabular}{cc}
\toprule
$V_{max}$ [m/s] & Tipo de suelo\\
\midrule
0.6 & Limo\\
0.9 & Arena\\
1.2 & Arcilla\\
3.0 & Asfalto\\
4.5 & Hormigón\\
\bottomrule
\end{tabular}
\end{table}
\paragraph{Velocidades mínimas} (evitar problemas de sedimentación):
\begin{formula}
$v_{min}=0.4\sim0.5$ [m/s]
\end{formula}
La condición de velocidad mínima también se puede controlar por la pendiente
de diseño (fondo del contrafoso):
\begin{table}[H]
\centering
\begin{tabular}{cc}
\toprule
$S_{min}$ [\%] & Tipo de suelo\\
\midrule
0.25 & Tierra\\
0.15 & Asfalto\\
\bottomrule
\end{tabular}
\end{table}

\paragraph{Revancha de seguridad ($R$):} evitar el rebalse del contrafoso
debido a alguna condición no prevista en el diseño.
\begin{formula}
$R=0.15\sim0.2\,h$
\end{formula}

\section{Geometría y periodo de retorno}
\paragraph{Geometría}
\begin{itemize}
  \item En caso de contrafosos de suelo natural, se recomiendan secciones de
        diseño trapezoidales.
  \item Además, se debe verificar que $\theta<\phi$, con $\phi$: ángulo de
        fricción del suelo.
\end{itemize}

\begin{figure}[H]
\centering
\begin{tikzpicture}[font=\small]
  \fill[suelo!80] (-0.8,0) -- (0.4,0) -- (-0.3,2.6) -- (-0.8,2.6) -- cycle;
  \fill[suelo!80] (4.4,0) -- (5.6,0) -- (5.6,2.6) -- (5.1,2.6) -- cycle;
  \fill[gray!30] (-0.3,2.6) -- (5.1,2.6) -- (4.85,1.7) -- (-0.05,1.7) -- cycle;
  \fill[top color=aguaclara,bottom color=agua] (-0.05,1.7) -- (4.85,1.7) -- (4.4,0) -- (0.4,0) -- cycle;
  \draw[thick] (-0.8,2.6) -- (-0.3,2.6) -- (0.4,0) -- (4.4,0) -- (5.1,2.6) -- (5.6,2.6);
  \pic at (0.4,1.72) {nivelagua};
  \draw[dashed] (4.4,0) -- (5.6,0);
  \draw (5.1,0) arc (0:75:0.7); \node at (5.25,0.35) {$\theta$};
  \draw[cota] (5.9,0) -- (5.9,1.7) node[midway,right]{$h$};
  \draw[cota] (5.9,1.7) -- (5.9,2.6) node[midway,right]{$R=0.15\sim0.2\,h$};
\end{tikzpicture}
\caption{Sección trapezoidal del contrafoso con revancha $R$.}
\end{figure}

\paragraph{Periodo de retorno de diseño}
\begin{itemize}
  \item Dependerá del tipo de camino y su importancia.
  \item También se puede diseñar considerando el riesgo de falla:
\end{itemize}
\begin{formula}
$\displaystyle R=1-\left[1-\frac1T\right]^n$
\end{formula}
\begin{table}[H]
\centering
\begin{tabular}{lccccc}
\toprule
& \multicolumn{2}{c}{$T$ [años]} & Vida útil & \multicolumn{2}{c}{Riesgo de falla $R$ [\%]}\\
\cmidrule(lr){2-3}\cmidrule(lr){5-6}
\textbf{Tipo de obra} & Diseño & Verificación & supuesta $n$ [años] & Diseño & Verificación\\
\midrule
Carreteras & 10 & 25 & 10 & 65 & 34\\
Caminos & 5 & 10 & 5 & 67 & 41\\
\bottomrule
\end{tabular}
\end{table}

\section{Bajadas de agua: rápidos o canal en gradería}
\begin{itemize}
  \item Su función es proteger de la erosión a los taludes de terraplenes y
        cortes, conduciendo el agua hasta un sitio seguro (cuneta al pie del talud).
  \item Diámetro mínimo de 200 mm para $H/V\geq4$.
  \item Gradería para razones $H/V$ 4:1 o menores.
  \item Pueden ser canaletas semicirculares en metal corrugado, $H/V$ 2:1.
  \item Velocidades de 5 m/s para el diseño hidráulico.
\end{itemize}

\begin{figure}[H]
\centering
\begin{tikzpicture}[font=\small]
  \fill[gray!25] (0,0) -- (0,0.4) -- (6,3.4) -- (6,0) -- cycle;
  \draw (0,0.4) -- (6,3.4);
  \draw[thick,azulusm] (0,0.4) \foreach \i in {0,...,6}{ -- ++(0,0.43) -- ++(0.86,0) };
  \foreach \i in {0,...,6} \draw[flecha,cyan!70!blue] ({0.86*\i+0.7},{0.43*\i+0.95}) -- ++(-0.5,-0.25);
  \node[align=center] at (4.5,1.0) {Canal en gradería\\(corte longitudinal)};
\end{tikzpicture}
\caption{Bajada de agua escalonada (canal en gradería) sobre un talud.}
\end{figure}

% =====================================================================
%  CLASE 10 -- Drenaje transversal de caminos y carreteras: alcantarillas
% =====================================================================
\clase{10}{Drenaje transversal de caminos y carreteras: alcantarillas}

\section{Drenaje transversal de caminos y carreteras}
El objetivo es no interrumpir cursos de agua existentes, por ejemplo: ríos,
quebradas, esteros, etc. Existen distintas soluciones:
\begin{enumerate}
  \item \textbf{Tubos de hormigón o HDPE} (polietileno de alta densidad):
  \begin{itemize}
    \item Si se tiene más de un tubo, la distancia libre entre tubos debe ser $D/2$.
    \item Solución utilizada para caudales pequeños.
    \item Más recomendaciones constructivas en el Manual de Carreteras V4.
  \end{itemize}
  \item \textbf{Cajones de hormigón armado:}
  \begin{itemize}
    \item Recomendaciones constructivas en el Manual de Carreteras V4.
    \item La solución de cajones múltiples compite con puentes pequeños.
  \end{itemize}
  \item \textbf{Puentes} (estribos y cepas): solución utilizada cuando se
        tienen caudales de magnitud importante, por ejemplo ríos o grandes quebradas.
\end{enumerate}
Tanto para cajones como para tubos se deben considerar \textbf{muros guía} a
la entrada y salida. También se conocen como \textbf{atraviesos o alcantarillas}.

\begin{figure}[H]
\centering
\begin{tikzpicture}[font=\small]
  \fill[suelo!80] (0,0) rectangle (9,1.8);
  \fill[gray!25] (0,1.8) .. controls (2.5,1.8) and (3.3,0.4) .. (4.5,0.4) .. controls (5.7,0.4) and (6.5,1.8) .. (9,1.8) -- cycle;
  \draw[thick] (0,1.8) -- (9,1.8);
  \draw[thick,fill=white] (4.5,1.0) circle (0.65);
  \begin{scope} \clip (4.5,1.0) circle (0.65); \fill[aguaclara] (3.8,0.3) rectangle (5.2,1.2); \end{scope}
  \draw[thick] (4.5,1.0) circle (0.65);
  \pic at (4.5,1.22) {nivelagua};
  \draw[dashed] (3.85,1.0) -- (3.85,-0.4) (5.15,1.0) -- (5.15,-0.4);
  \draw[cota] (3.85,-0.3) -- (5.15,-0.3) node[midway,below]{$D$};
  \draw[<-] (6.3,1.4) -- (7.2,2.3) node[above,azulusm]{Terraplén};
  \draw[<-] (5.4,0.7) -- (6.6,-0.3) node[right,azulusm]{Material de relleno};
\end{tikzpicture}
\caption{Alcantarilla de tubo bajo un terraplén (corte transversal al cauce).}
\end{figure}

\section{Alcantarillas}
\begin{itemize}
  \item Son conductos relativamente cortos a través de los cuales se cruza el
        agua por debajo del camino o carretera, desde un lado al otro de éste.
        Se entiende por alcantarilla una \textbf{estructura de drenaje} cuya
        luz, medida paralela al eje de la carretera, sea de \textbf{hasta seis
        metros (ancho)}; losas mayores se tratarán como puentes (MC, 3.703.101).
  \item Se consideran estructuras menores en el diseño y construcción de una
        vía, pero deben ser diseñadas y construidas para soportar las cargas
        del tránsito, el peso de tierra sobre ellas y las cargas durante la construcción.
  \item El diseño consiste en determinar el \textbf{diámetro más económico}
        que permita pasar el caudal de diseño sin exceder la carga máxima de
        entrada y atendiendo criterios de arrastre de sedimentos (velocidades
        de escurrimiento).
  \item Criterios para este tipo de obras:
  \begin{enumerate}[label=\alph*)]
    \item En planta, las obras se dispondrán en la dirección del cauce.
    \item En perfil, se deberá ajustar al perfil de la obra.
    \item La sección debe respetar las dimensiones del cauce natural y no
          provocar fuertes estrechamientos.
  \end{enumerate}
\end{itemize}

\subsection{Algunas recomendaciones (Horacio Mery)}
\begin{itemize}
  \item El nivel aguas arriba sea similar al que se tiene en el cauce para la
        crecida de diseño. ¿Para qué?
  \item El nivel de aguas arriba debe resguardar la seguridad del terraplén.
  \item La mayoría de las alcantarillas se diseñan para operar con
        escurrimiento libre cercano al crítico, logrando $Fr$ entre 0.7 y 0.9.
  \item Se busca maximizar el caudal pasante, reduciendo la sección
        transversal y por tanto el costo.
\end{itemize}

\paragraph{Tragedia del estero Minte (07/05/1995).} El ancho de la
alcantarilla era algo estrecho frente a las precipitaciones (192 mm de agua
en 48 horas); la entrada de la alcantarilla se tapó con árboles y barro,
generando la falla de la autopista, que causó la muerte de 27 personas, 15
de ellas niños. La justicia dictaminó que el Estado fue responsable como
consecuencia de un mal diseño.

\section{Diseño hidráulico de alcantarillas}
\subsection{Recomendaciones de diseño}
\begin{itemize}
  \item La tubería o cajón debe seguir la línea de flujo del agua (en planta,
        alineada con el cauce; no forzar el cauce a cruzar perpendicular al camino).
  \item En un cauce amplio es mejor usar tubos múltiples que dispersen el
        flujo, en lugar de concentrarlo en un solo tubo.
  \item La tubería o cajón debe seguir la pendiente del terreno: no instalarla
        muy profunda ni muy alta; no alterar la elevación del fondo del canal.
\end{itemize}

\begin{figure}[H]
\centering
\begin{tikzpicture}[font=\small]
  \foreach \X/\ok in {0/0,5/1}{
  \begin{scope}[xshift=\X cm]
    \fill[gray!50] (-0.3,1.2) rectangle (3.7,2.0);
    \ifnum\ok=0
      \fill[aguaclara,draw=azulusm] (0.9,3.2) -- (1.5,3.2) -- (1.6,2.0) -- (1.0,2.0) -- cycle;
      \fill[aguaclara,draw=azulusm] (1.0,1.2) -- (1.6,1.2) -- (1.5,0) -- (0.9,0) -- cycle;
      \fill[gray!20,draw] (1.0,1.2) rectangle (1.6,2.0);
      \draw[rojo,line width=3pt] (2.3,0.3) -- (3.1,1.0) (2.3,1.0) -- (3.1,0.3);
    \else
      \fill[aguaclara,draw=azulusm] (1.8,3.2) -- (2.4,3.2) -- (1.9,2.0) -- (1.3,2.0) -- cycle;
      \fill[aguaclara,draw=azulusm] (1.0,1.2) -- (1.6,1.2) -- (1.1,0) -- (0.5,0) -- cycle;
      \fill[gray!20,draw] (1.3,2.0) -- (1.9,2.0) -- (1.6,1.2) -- (1.0,1.2) -- cycle;
      \draw[azulusm,line width=3pt] (2.4,0.7) -- (2.7,0.3) -- (3.4,1.1);
    \fi
  \end{scope}}
\end{tikzpicture}
\caption{La alcantarilla debe seguir la línea de flujo del cauce (derecha), no
forzar un cruce perpendicular (izquierda).}
\end{figure}

% ---------------------------------------------------------------------
% Macro: perfil longitudinal de una alcantarilla
% #1: contenido extra
\newcommand{\perfilalcantarilla}[1]{%
\begin{tikzpicture}[font=\small,scale=1.0]
  % fondo inclinado
  \draw[thick,gray!70!black] (0,0.6) -- (9,0);
  \draw[dashed] (0,0.2) -- (8,0.2);
  % alcantarilla (cajón) y terraplén
  \fill[suelo!80] (3.1,2.6) -- (3.4,2.8) -- (7.6,2.5) -- (7.9,2.0) -- (7.9,1.4) -- (3.1,1.8) -- cycle;
  \draw[thick] (3.0,2.8) -- (3.0,1.85) -- (8.0,1.5) -- (8.0,2.0);
  #1
\end{tikzpicture}}

\subsection{Criterio 1: diseño con control de entrada}
Se diseña la alcantarilla para que el escurrimiento en ella sea
\textbf{torrencial}:
\begin{center}
\colorbox{azulusm}{\parbox{0.7\textwidth}{\centering\color{white}\large
Independencia aguas arriba de lo que ocurre aguas abajo}}
\end{center}
Así, asumiendo la condición más crítica, es decir aguas arriba de la
alcantarilla se tiene escurrimiento de río:
\begin{center}
\colorbox{azulusm}{\parbox{0.7\textwidth}{\centering\color{white}\large
Caída hidráulica en la entrada de la alcantarilla}}
\end{center}

\begin{figure}[H]
\centering
\perfilalcantarilla{
  \fill[aguaclara] (0,0.6) -- (0,2.4) .. controls (2,2.3) and (2.6,1.4) .. (3.1,1.1) .. controls (5,0.7) and (7,0.4) .. (8,0.35) -- (8,0.06) -- cycle;
  \draw[azulusm,thick] (0,2.4) .. controls (2,2.3) and (2.6,1.4) .. (3.1,1.1) .. controls (5,0.7) and (7,0.4) .. (8,0.35);
  \draw[cota] (0.3,0.58) -- (0.3,2.4) node[midway,right]{$H=h_n$};
  \draw[cota] (3.2,0.4) -- (3.2,1.08) node[midway,right]{$h_c$};
  \draw[cota] (7.3,0.12) -- (7.3,1.55) node[midway,left]{$D$};
  \draw[cota] (8.2,0.05) -- (8.2,0.35) node[midway,right]{$h_T$};
  \node[azulusm] at (0.8,2.8) {Río}; \node[azulusm,rotate=-5] at (5.4,0.9) {Torrente};
  \node at (1.5,0.35) {$S_D$};
}
\caption{Control de entrada: escurrimiento de río aguas arriba y torrente en la alcantarilla.}
\end{figure}

\paragraph{Sección rectangular.} Con $H=h+\frac{v^2}{2g}$:
\[
H_c=\frac32h_c,\qquad h_c=\sqrt[3]{\frac{Q^2}{gb^2}}
\]
Despreciando pérdidas de energía, $H_c=H$; y considerando $D=H_c$:
\begin{formula}
$\displaystyle D=\frac32\frac{1}{\sqrt[3]g}\left(\frac Qb\right)^{2/3}$
\end{formula}
Para determinar el ancho $b$ de la alcantarilla se impone $D=H$. Luego,
considerando la ecuación de Manning para $h=h_c=\frac23H$:
\begin{formula}
$\displaystyle S_c=\frac23g\cdot H\cdot n^2\left(\frac{b+\frac43H}{\frac23H\cdot b}\right)^{4/3}$
\end{formula}
La pendiente de diseño debe cumplir con
\[
S_D\geq S_c\ \Rightarrow\ \text{torrente};\qquad\text{si } S_D\leq S_c\ \Rightarrow\ S_D=S_c
\]
Una vez definida la geometría se debe verificar que $v<5$ [m/s].

\paragraph{Sección circular} (considerando $h_c=0.7D$, con $D$ en [m] y $Q$ en [m$^3$/s]):
\begin{formula}
$Q=1.425\,D^{5/2}\qquad\qquad\displaystyle S_c=\frac{31.14\,n^2}{D^{1/3}}$
\end{formula}
Por seguridad se debe diseñar con $D>H$.

\begin{table}[H]
\centering\small
\renewcommand{\arraystretch}{2.0}
\begin{tabular}{l>{$}c<{$}>{$}c<{$}>{$}c<{$}>{$}c<{$}}
\toprule
\textbf{Sección} & \text{Área } A & \text{Perímetro } P & \text{Radio hidráulico } R_h & \text{Ancho superficial } b_s\\
\midrule
Rectangular & by & b+2y & \dfrac{by}{b+2y} & b\\
Trapezoidal & (b+ky)y & b+2y\sqrt{1+k^2} & \dfrac{(b+ky)y}{b+2y\sqrt{1+k^2}} & b+2ky\\
Triangular & ky^2 & 2y\sqrt{1+k^2} & \dfrac{ky}{2\sqrt{1+k^2}} & 2ky\\
Circular & \dfrac{(\theta-\operatorname{sen}\theta)D^2}{8} & \dfrac{\theta D}{2} & \left(1-\dfrac{\operatorname{sen}\theta}{\theta}\right)\dfrac D4 & \operatorname{sen}\left(\dfrac\theta2\right)D\ \text{ó}\ 2\sqrt{y(D-y)}\\
Parabólica & \dfrac23 yb_s & b_s+\dfrac{8y^2}{3b_s} & \dfrac{2b_s^2y}{3b_s^2+8y^2} & \dfrac{3A}{2y}\\
\bottomrule
\end{tabular}
\caption{Elementos geométricos de secciones típicas; en la sección circular
$\theta=2\cos^{-1}\left(1-\frac{2y}{D}\right)$.}
\end{table}

\subsection{Criterio 2: diseño y verificación}
Se toman en cuenta los caudales correspondientes a 2 períodos de retorno:
\begin{itemize}
  \item $T_1$: diseño con control de entrada, p.~ej.\ $T_1=25$ años.
  \item $T_2$: verificación a sección completa, p.~ej.\ $T_2=50$ años
        (entrada ahogada).
\end{itemize}
Se busca que el flujo no cause sobrepresiones en la cara superior del cajón.

\paragraph{Verificación a sección completa.} Se calcula la pendiente normal
para $Q_{2D}$ ($T_2$) imponiendo $h_n=D$:
\begin{formula}
Sección circular: $\displaystyle S_n=\frac{10.3\,Q_{2D}^2\cdot n^2}{D^{16/3}}$\qquad\qquad
Sección rectangular: $\displaystyle\frac{Q_{2D}\cdot n}{\sqrt{S_n}}=\Omega\cdot R_H^{2/3}$
\end{formula}
La comparación con la pendiente de diseño ($S_D$) genera tres casos:
\begin{enumerate}
  \item $S_D>S_n$
  \item $S_D=S_n$
  \item $S_D<S_n$
\end{enumerate}

\paragraph{Caso 1: $S_D>S_n$.} La alcantarilla actúa como compuerta
(contracción a la entrada). Se debe calcular $H'$ para verificar que no se
inunde el camino. Se conserva la energía entre (1) y (2):
\begin{formula}
$\displaystyle H'=c_c\cdot D+\frac{1}{2g}\left(\frac{Q_{2D}}{b\cdot c_c\cdot D}\right)^2,\qquad c_c\approx0.611$
\end{formula}

\begin{figure}[H]
\centering
\perfilalcantarilla{
  \fill[aguaclara] (0,0.6) -- (0,2.7) -- (3.0,2.7) -- (3.0,1.85) .. controls (3.4,1.2) and (4,1.0) .. (5,0.85) .. controls (6.5,0.55) and (7.5,0.4) .. (8,0.35) -- (8,0.06) -- cycle;
  \draw[azulusm,thick] (0,2.7) -- (3.0,2.7);
  \draw[azulusm,thick] (3.0,1.85) .. controls (3.4,1.2) and (4,1.0) .. (5,0.85) .. controls (6.5,0.55) and (7.5,0.4) .. (8,0.35);
  \pic at (2.5,2.72) {nivelagua};
  \draw[cota] (1.3,0.52) -- (1.3,2.7) node[midway,right]{$H'$};
  \draw[cota] (4.3,0.33) -- (4.3,0.95) node[midway,right]{$h_2=c_cD$};
  \draw[cota] (7.3,0.12) -- (7.3,1.55) node[midway,left]{$D$};
  \node at (1.3,0.2) {(1)}; \node at (4.3,0) {(2)};
  \node[azulusm] at (0.6,2.4) {Río}; \node[azulusm,rotate=-5] at (6,0.95) {Torrente};
}
\caption{Caso 1 ($S_D>S_n$): la alcantarilla actúa como compuerta.}
\end{figure}

\paragraph{Caso 2: $S_D=S_n$.} Se puede realizar un balance de energía entre
(1) y (2) considerando la pérdida singular a la entrada:
\begin{formula}
$\displaystyle H'=D+\frac1{2g}\,(k_e+1)\cdot\frac{Q_{2D}^2}{(D\cdot b)^2}$
\end{formula}
El valor de $k_e$ dependerá de la forma del flujo de entrada:
\begin{center}
\begin{tabular}{lc}
\toprule
Tipo de entrada & $k_e$\\
\midrule
Tubo que sobresale del muro (entrada reentrante) & 0.9\\
Entrada abocinada (redondeada) & 0.2\\
Entrada a ras del muro (arista viva) & 0.5\\
\bottomrule
\end{tabular}
\end{center}

\paragraph{Caso 3: $S_D<S_n$.} Se producen sobrepresiones sobre la cara
posterior de la alcantarilla; el escurrimiento aguas arriba se ve
influenciado por lo que ocurre aguas abajo. Se realiza un balance de energía
entre (1) y (2) considerando las pérdidas singulares y por fricción:
\begin{formula}
$\displaystyle H'=D+\frac{v_2^2}{2g}+\Lambda_e+J\cdot L-S_D\cdot L,\qquad
\Lambda_e=k_e\frac{v_2^2}{2g},\qquad J=\frac{v_2^2\cdot n^2}{R_H^{4/3}}$
\end{formula}

\begin{figure}[H]
\centering
\perfilalcantarilla{
  \fill[aguaclara] (0,0.6) -- (0,2.7) -- (3.0,2.7) -- (3.0,1.85) -- (8.0,1.5) -- (8.0,0.0) -- (0,0.6);
  \draw[azulusm,thick] (0,2.7) -- (3.0,2.7);
  \pic at (2.5,2.72) {nivelagua};
  \foreach \y in {0.3,0.6,0.9,1.2} \draw[->,azulusm] (8.0,\y) .. controls (8.4,\y) .. (8.8,\y+0.2);
  \draw[cota] (1.3,0.52) -- (1.3,2.7) node[midway,right]{$H'$};
  \draw[cota] (7.3,0.06) -- (7.3,1.55) node[midway,left]{$D$};
  \draw[dashed] (3.0,0.37) -- (3.0,-0.6) (8.0,0.0) -- (8.0,-0.6);
  \draw[cota] (3.0,-0.5) -- (8.0,-0.5) node[midway,below]{$L$};
  \node at (1.3,0.2) {(1)}; \node[right] at (8.0,-0.2) {(2)};
}
\caption{Caso 3 ($S_D<S_n$): alcantarilla trabajando a sección llena.}
\end{figure}

\paragraph{Adicional: control de salida.} La pendiente aguas abajo es tan
pequeña que el escurrimiento de río ahoga al atravieso. Se realiza un
balance de energía entre (1) y (3) considerando las pérdidas singulares
(contracción a la entrada y expansión a la salida) y por fricción:
\begin{formula}
$\displaystyle H'=D+\frac{v_2^2}{2g}+\Lambda_e+\Lambda_s+J\cdot L-S_D\cdot L$
\end{formula}
donde
\[
\Lambda_s=k_s\frac{(v_2-v_3)^2}{2g},\qquad\Lambda_e=k_e\frac{v_2^2}{2g},\qquad J=\frac{v_2^2\cdot n^2}{R_H^{4/3}},
\]
\[
k_s=0.22-0.44\ \ (k_s=0.4\ \text{recomendado})
\]
y $h_3=h_{n3}$ es la altura normal aguas abajo. Referencia: FHWA, \emph{Hydraulic
Design of Highway Culverts},
\url{https://www.fhwa.dot.gov/engineering/hydraulics/library_arc.cfm?id=13&pub_number=7}.

\subsection{Períodos de retorno}
\begin{table}[H]
\centering\small
\resizebox{\linewidth}{!}{%
\begin{tabular}{llcccccc}
\toprule
& & \multicolumn{2}{c}{$T$ [años]} & Vida útil & \multicolumn{2}{c}{Riesgo de falla [\%]}\\
\cmidrule(lr){3-4}\cmidrule(lr){6-7}
\textbf{Tipo de obra} & \textbf{Ruta} & Diseño & Verif. & $n$ [años] & Diseño & Verif.\\
\midrule
Puentes y viaductos & Carreteras & 200 & 300 & 50 & 22 & 15\\
 & Caminos & 100 & 150 & 50 & 40 & 28\\
Alcantarillas ($S>1.75$ m$^2$) o $H$ terrap.\ $\geq10$ m & Carreteras & 100 & 150 & 50 & 40 & 28\\
\quad y estructuras enterradas & Caminos & 50 & 100 & 30 & 45 & 26\\
Alcantarillas ($S<1.75$ m$^2$) & Carreteras & 50 & 100 & 50 & 64 & 40\\
 & Caminos & 25 & 50 & 30 & 71 & 45\\
\bottomrule
\end{tabular}}
\caption{Períodos de retorno (Manual de Carreteras, Volumen 3). $S$ es el área útil de la alcantarilla.}
\end{table}

\subsection{Alturas y velocidades}
\begin{table}[H]
\centering\small
\begin{tabular}{llll}
\toprule
\textbf{Tipo de cauce} & \textbf{Tubos} & \textbf{Cajones} & \textbf{Losas} ($L\leq6.0$ m)$^*$\\
\midrule
Canales & $D$ (diámetro) & $H$ (altura total) & $H-0.10$ m\\
Diseño cauces naturales & $D+0.3$ m & $H+0.3$ m & $H-0.10$ m\\
Verificación cauces naturales & $D+0.6$ m & $H+0.6$ m & $H$\\
\multicolumn{4}{l}{\quad Pero $H_e$ máximo no puede sobrepasar la cota exterior del SAP $-0.3$ m}\\
\bottomrule
\end{tabular}
\caption{Tabla 3.703.301.A: carga hidráulica de diseño ($H_e$, m).
$^*$Si $L>6.0$ m, revancha como en puentes.}
\end{table}

\begin{table}[H]
\centering\small
\begin{tabular}{lcc}
\toprule
\textbf{Tipo de terreno} & \textbf{Flujo intermitente} [m/s] & \textbf{Flujo permanente} [m/s]\\
\midrule
Arena fina (no coloidal) & 0.75 & 0.75\\
Arcilla arenosa (no coloidal) & 0.75 & 0.75\\
Arcilla limosa (no coloidal) & 0.90 & 0.90\\
Arcilla fina & 1.00 & 1.00\\
Ceniza volcánica & 1.20 & 1.00\\
Grava fina & 1.50 & 1.20\\
Arcilla dura (coloidal) & 1.80 & 1.40\\
\multicolumn{3}{l}{\emph{Material graduado (no coloidal):}}\\
\quad desde arcilla a grava & 2.00 & 1.50\\
\quad desde limo a grava & 2.10 & 1.70\\
Grava & 2.30 & 1.80\\
Grava gruesa & 2.40 & 2.00\\
Desde grava a piedras (bajo 15 cm) & 2.70 & 2.10\\
Desde grava a piedras (sobre 20 cm) & 3.00 & 2.40\\
\bottomrule
\end{tabular}
\caption{Tabla 3.703.301.B: velocidades máximas admisibles en canales no
revestidos (fuente: Manual de Carreteras de California).}
\end{table}

% =====================================================================
%  CLASE 11 -- Drenaje transversal: puentes
% =====================================================================
\clase{11}{Drenaje transversal de caminos y carreteras: puentes}

\section{Diseño hidráulico de puentes}
Cuando el caudal a atravesar es muy grande, la única alternativa es la
construcción de un puente.
\begin{itemize}
  \item \textbf{Opción 1: terraplén y estribos} (más económico). \emph{Método
        fusible}: el objetivo es que el río, en una crecida, se lleve el
        terraplén y no la estructura.
  \item \textbf{Opción 2: puente que cubre todo el cauce, con cepas} (¡alto costo!).
\end{itemize}

\begin{figure}[H]
\centering
\begin{subfigure}{0.48\textwidth}\centering
\begin{tikzpicture}[font=\small,scale=0.8]
  \fill[suelo!80] (0,3) .. controls (0.8,3) and (0.8,1.5) .. (2,1.3) -- (2.4,0.6) -- (4.4,0.6) -- (4.8,1.3) .. controls (6,1.5) and (6.2,3) .. (7,3) -- (7,0) -- (0,0) -- cycle;
  \fill[sueloclaro] (2,1.4) -- (2.4,0.6) -- (3.0,0.6) -- (3.0,1.4) -- cycle;
  \fill[sueloclaro] (3.9,1.4) -- (3.9,0.6) -- (4.4,0.6) -- (4.8,1.4) -- cycle;
  \fill[aguaclara] (3.0,0.6) rectangle (3.9,0.9);
  \draw[very thick,gray] (1.2,1.45) -- (5.8,1.45);
  \draw[thick] (3.0,0.6) -- (3.0,1.45) (3.9,0.6) -- (3.9,1.45);
  \node[azulusm] at (3.45,2.4) {Terraplén};
  \draw[->,azulusm] (3.1,2.2) -- (2.6,1.1); \draw[->,azulusm] (3.8,2.2) -- (4.3,1.1);
\end{tikzpicture}
\caption{Opción 1: terraplén y estribos.}
\end{subfigure}\hfill
\begin{subfigure}{0.48\textwidth}\centering
\begin{tikzpicture}[font=\small,scale=0.8]
  \fill[suelo!80] (0,3) .. controls (0.3,1.5) and (1.2,1.5) .. (1.8,0.6) -- (5.2,0.6) .. controls (5.8,1.5) and (6.7,1.5) .. (7,3) -- (7,0) -- (0,0) -- cycle;
  \fill[aguaclara] (3.1,0.6) rectangle (3.9,0.8);
  \fill[gray!60] (0,3) rectangle (7,3.2);
  \foreach \x in {1.2,5.8} \fill[gray!60] (\x-0.1,1.5) rectangle (\x+0.1,3);
\end{tikzpicture}
\caption{Opción 2: puente con cepas.}
\end{subfigure}
\caption{Alternativas para cruzar un cauce de gran caudal.}
\end{figure}

\subsection{Recomendaciones de diseño}
\begin{itemize}
  \item No generar resaltos.
  \item Evitar que el peraltamiento del escurrimiento aguas arriba produzca
        inundaciones que antes no ocurrían.
  \item Que no se generen velocidades excesivas que puedan provocar problemas
        de socavación.
\end{itemize}
(Ejemplo reciente: desborde del Canal San Carlos, 16 de junio de 2024.)

\section{Opción 1: terraplén y estribos}
Considerando escurrimiento de río en el cauce (ancho $b$) que se contrae al
ancho $L$ del puente. Se definen las secciones (1) aguas abajo, (2) en la
contracción y (3) aguas arriba.

\begin{figure}[H]
\centering
\begin{tikzpicture}[font=\small]
  \fill[top color=aguaclara,bottom color=agua] (0,0) rectangle (8,2.6);
  \fill[sueloclaro,draw] (3.2,1.9) rectangle (4.4,2.6);
  \fill[sueloclaro,draw] (3.2,0) rectangle (4.4,0.7);
  \fill[gray!60] (3.2,1.8) rectangle (4.4,1.9); \fill[gray!60] (3.2,0.7) rectangle (4.4,0.8);
  \draw[cota] (0.9,0) -- (0.9,2.6) node[midway,left]{$b$};
  \draw[cota] (3.8,0.8) -- (3.8,1.8) node[midway,left]{$L$};
  \foreach \x/\n in {2.6/3,4.4/2,5.4/1} {\draw[rojo,dashed,thick] (\x,-0.4) -- (\x,3.0); \node[rojo,below] at (\x,-0.4) {(\n)};}
  \node[azulusm] at (6.6,3.3) {Terraplén};
  \draw[->,azulusm] (6.2,3.1) -- (4.3,2.4); \draw[->,azulusm] (6.9,3.1) .. controls (7.8,-0.6) and (5,-0.6) .. (4.3,0.3);
  \draw[flecha,azulusm] (7.6,1.3) -- (6.9,1.3);
\end{tikzpicture}
\caption{Contracción del cauce por terraplenes y estribos (vista en planta;
el flujo va de (3) hacia (1)).}
\end{figure}

La momenta (fuerza específica) en cada sección es:
\begin{formula}
$\displaystyle M_1=\frac{Q^2}{g(b\cdot h_1)}+\frac{h_1}{2}(b\cdot h_1)\qquad
M_2=\frac{Q^2}{g(L\cdot h_2)}+\frac{h_2}{2}(b\cdot h_2)$
\end{formula}

\begin{figure}[H]
\centering
\begin{tikzpicture}
\begin{axis}[width=8cm,height=6cm,xmin=0,xmax=10,ymin=0,ymax=5,axis lines=left,
  xlabel={$M$},ylabel={$h$},xtick=\empty,ytick=\empty,clip=false,font=\small]
  \addplot[thick,azulusm,domain=0.35:4.6,samples=60] ({6/x+x^2/2*0.9},{x});
  \addplot[thick,rojo,domain=0.45:4.6,samples=60] ({9/x+x^2/2*0.9},{x});
  \addplot[thick,green!60!black,domain=0.5:4.6,samples=60] ({11/x+x^2/2*0.9},{x});
  \addplot[thick,orange,domain=0.55:4.6,samples=60] ({13/x+x^2/2*0.9},{x});
  \draw[dashed] (axis cs:0,3.4) node[left]{$h_{n1}$} -- (axis cs:{6/3.4+3.4^2*0.45},3.4) -- (axis cs:{6/3.4+3.4^2*0.45},0) node[below]{$M_1(h_{n1})$};
\end{axis}
\end{tikzpicture}
\caption{Curvas de momenta: sección 1 (azul) y sección 2 con $M_{c2}<M_1$
(rojo), $M_{c2}=M_1$ (verde) y $M_{c2}>M_1$ (naranja, se produce resalto).}
\end{figure}

Para la sección (2) se tiene:
\[
\frac{\partial M_2}{\partial h_2}=0\quad\Longrightarrow\quad h_{c2,M}=\sqrt[3]{\frac{Q^2}{g\cdot b\cdot L}}
\]
Por lo tanto, para definir el largo mínimo del puente ($L$), tal que no se
produzca resalto, se considera la condición crítica, resultando el siguiente
sistema de ecuaciones:
\begin{formula}
$M_1=M_{2c}\quad(1)\qquad\qquad\displaystyle h_{c2,M}=\sqrt[3]{\frac{Q^2}{g\cdot b\cdot L}}\quad(2)$
\end{formula}
donde
\[
M_{2c}=\frac{Q^2}{g(L\cdot h_{2c,M})}+\frac{h_{2c,M}}{2}(b\cdot h_{2c,M}),\qquad
M_1=\frac{Q^2}{g(b\cdot h_1)}+\frac{h_1}{2}(b\cdot h_1)\ \ \text{(conocido)}
\quad\Longrightarrow\quad\boxed{L\geq L_{min}}
\]

Con $L$ definido para que no se produzca resalto, se puede realizar un
balance de energía entre (1) y (2):
\begin{formula}
$\displaystyle h_2+\frac{v_2^2}{2g}=h_{n1}+\frac{v_1^2}{2g}+\Lambda_s$
\end{formula}
donde las pérdidas singulares a la salida se calculan a través de la
ecuación de Borda--Carnot:
\begin{formula}
$\displaystyle\Lambda_s=\frac{(v_2-v_1)^2}{2g}-\frac{(h_{n1}-h_2)^2}{2h_{n1}}$
\end{formula}
Además, $v_1=\dfrac{Q}{h_{n1}b}$ y $v_2=\dfrac{Q}{h_2L}$, de donde se obtiene $h_2$.

Luego, realizando un balance de energía entre (2) y (3):
\begin{formula}
$\displaystyle h_3+\frac{v_3^2}{2g}=h_2+\frac{v_2^2}{2g}+\Lambda_e,\qquad
v_3=\frac{Q}{h_3b},\qquad\Lambda_e=k_e\frac{v_2^2}{2g},\quad k_e=0.1\sim0.3$
\end{formula}
de donde se obtiene $h_3$. Se debe verificar que
\[
h_3<h_{max}\qquad\qquad v_2<v_{max}
\]
Si no se cumple alguna de las condiciones, se debe elegir un $L$ mayor y
recalcular todo hasta que se cumplan; o bien, a partir de dichas condiciones
se determinan $L$ y $h_2$ y se verifica que no haya resalto.

\paragraph{Conservación del momentum lineal entre (3) y (2).} También se
puede aplicar, con:
\[
n=\frac{b_3}{b_2},\qquad X_3=\frac{h_3}{h_{c2}},\qquad X_2=\frac{h_2}{h_{c2}},\qquad h_{c2}=\sqrt[3]{\frac{Q^2}{g\,b_2^2}}
\]
\begin{formula}
$\displaystyle\frac{1}{X_2}-\frac{1}{nX_3}=\frac12\left(X_3^2-X_2^2\right)$\qquad
($X_2$: 1 -- 6;\quad $n$: 1 -- 4.5)
\end{formula}

\section{Opción 2: cepas --- ecuación de Yarnell}
\begin{itemize}
  \item Basada en miles de experimentos realizados entre 1927 y 1936 en
        cauces obstruidos por el efecto de las cepas de diferentes puentes.
  \item Método válido si no se producen resaltos a la salida del puente.
\end{itemize}
\begin{formula}
$\displaystyle\frac{h_3-h_1}{h_1}=\frac{\Delta h}{h_1}=K\cdot Fr_1^2\left(K+5Fr_1^2-0.6\right)\left(\sigma+1.5\sigma^4\right)$
\end{formula}
donde $n$: número de cepas; $b$: ancho del puente [m]; $\sigma=1-\frac{b_e}{b}$:
constante de ancho efectivo (fracción obstruida); $b_e=b-n\cdot D$: ancho
efectivo del puente [m]; $K$: factor de forma de la cepa (sección
transversal), $K=0.9\sim1.15$; $Fr_1$: número de Froude de la sección 1
(aguas abajo).

\begin{figure}[H]
\centering
\begin{tikzpicture}[font=\small,rotate=-20]
  \fill[top color=aguaclara,bottom color=agua] (0,0) rectangle (6,3);
  \fill[gray!40] (2.4,-0.6) rectangle (3.6,3.6);
  \foreach \y in {0.5,1.3,2.1} \fill[gray!70,rounded corners=3pt,draw] (2.3,\y) rectangle (3.7,\y+0.45);
  \foreach \y in {0.35,1.1,1.9,2.75} \draw[azulusm,thick] (0.5,\y) .. controls (2,\y) and (4,\y) .. (5.5,\y);
  \draw[rojo,dashed,thick] (1.8,-0.5) -- (1.8,3.5) node[above]{(3)};
  \draw[rojo,dashed,thick] (4.2,-0.5) -- (4.2,3.5) node[above]{(1)};
  \draw[cota] (5.7,0) -- (5.7,3) node[midway,right]{$b$};
  \draw[cota] (2.1,1.3) -- (2.1,1.75) node[midway,left]{$D$};
\end{tikzpicture}
\caption{Cepas en el cauce (vista en planta).}
\end{figure}

\begin{table}[H]
\centering
\begin{tabular}{lc}
\toprule
\textbf{Forma de la cepa del puente} & \textbf{Coeficiente de Yarnell, $K$}\\
\midrule
Nariz y cola semicircular & 0.90\\
Cepas de doble cilindro con diafragma de conexión & 0.95\\
Cepas de doble cilindro sin diafragma & 1.05\\
Nariz y cola triangular de 90° & 1.05\\
Nariz y cola cuadrada & 1.25\\
Estructura tipo \emph{trestle bent} de diez pilotes & 2.50\\
\bottomrule
\end{tabular}
\caption{Coeficiente de pila de Yarnell.}
\end{table}

\section{Períodos de retorno}
Se utilizan los mismos períodos de retorno del Manual de Carreteras,
Volumen 3, presentados en la clase anterior: para puentes y viaductos,
$T=200$ años (diseño) y 300 años (verificación) en carreteras, y $T=100$ y
150 años en caminos, con vida útil de 50 años.

% =====================================================================
%  CLASE 12 -- Sistemas de drenaje de aguas lluvias
% =====================================================================
\clase{12}{Sistemas de drenaje urbano de aguas lluvias}

\section{Motivación}
Las inundaciones de calles en temporales son frecuentes en ciudades con
sistemas de drenaje insuficientes o mal mantenidos (sumideros tapados, calles
anegadas). Ver: \emph{La historia del alcantarillado del Santiago del
Centenario y su resistencia ante la lluvia}.

\section{Sistemas de drenaje de aguas lluvias}
El objetivo de estos sistemas es \textbf{evitar la inundación de las calles}:
el agua escurre por las calles (con pendientes $S_x$, $S_y$), es captada por
sumideros y conducida por colectores hasta un cauce receptor (río).

\begin{figure}[H]
\centering
\begin{tikzpicture}[font=\small]
  \fill[gray!30,rounded corners=8pt] (-0.3,-0.3) rectangle (6.9,6.9);
  \foreach \i in {0,1,2} \foreach \j in {0,1,2}
    \fill[suelo!80,rounded corners=4pt] ({0.2+2.2*\i},{0.6+2.2*\j}) rectangle ({1.8+2.2*\i},{2.2+2.2*\j});
  \foreach \i in {0,1,2} \foreach \j in {0,1,2}{
    \draw[fill=white] ({2.0+2.2*\i},{0.4+2.2*\j}) circle (0.17);
    \draw[flecha,azulusm] ({0.3+2.2*\i},{0.4+2.2*\j}) -- ({1.75+2.2*\i},{0.4+2.2*\j});
    \draw[flecha,azulusm] ({2.0+2.2*\i},{2.1+2.2*\j}) -- ({2.0+2.2*\i},{0.65+2.2*\j});}
  \draw[flecha,azulusm,line width=2pt] (6.4,0.4) -- (7.3,-0.4) node[right]{Río};
  \draw[flecha] (3,7.1) -- (4,7.1) node[midway,above]{$S_x$};
  \draw[flecha] (-0.5,5.5) -- (-0.5,4.5) node[midway,left]{$S_y$};
\end{tikzpicture}
\caption{Esquema de una red urbana: escurrimiento por las calles, captación
en sumideros/cámaras (círculos) y descarga al río.}
\end{figure}

\begin{itemize}
  \item Las obras de drenaje están diseñadas para captar, retener, almacenar,
        infiltrar, transportar y descargar aguas lluvias.
  \item Las obras hidráulicas de drenaje son para redes domiciliaria,
        secundaria y primaria.
  \item En la red domiciliaria y secundaria se da mayor importancia a la
        infiltración y al almacenamiento.
  \item En la red primaria se da relevancia al transporte y descarga en las
        redes de aguas abajo.
\end{itemize}
Referencia: \emph{Manual de Drenaje Urbano (MDU)}, DOH: Cap.~4 Estudios
Básicos; Cap.~5 Técnicas de Diseño de Redes; Cap.~6 Diseño Hidráulico de Obras.

\begin{figure}[H]
\centering
\begin{tikzpicture}[font=\scriptsize,node distance=0.5cm,
  caja/.style={draw,fill=#1,minimum width=1.7cm,minimum height=0.7cm,align=center}]
  \node[caja=gray!20] (r) {Redes};
  \node[caja=rojo!20,right=of r] (d) {Red\\domiciliaria};
  \node[caja=suelo!40,right=of d] (s) {Red\\secundaria};
  \node[caja=green!25,right=of s] (p) {Red\\primaria};
  \node[caja=aguaclara,right=of p] (n) {Red\\natural};
  \node[caja=gray!20,below=of r] (e) {Elementos\\propios};
  \node[caja=rojo!20,below=of d] (d2) {Techos\\pavimentos};
  \node[caja=suelo!40,below=of s] (s2) {Calles\\sumideros};
  \node[caja=green!25,below=of p] (p2) {Cauces\\descargas};
  \node[caja=aguaclara,below=of n] (n2) {Zonas\\inundables};
  \draw[->] (d) -- (s); \draw[->] (s) -- (p); \draw[->] (p) -- (n);
  \foreach \a/\b in {d2/d,s2/s,p2/p,n2/n} \draw[->] (\a) -- (\b);
\end{tikzpicture}
\caption{Organización de las redes de drenaje urbano (MDU, Fig.~5.1.2).
Complementos: infiltración, almacenamiento y conducción.}
\end{figure}

Estos sistemas están compuestos por: \textbf{cunetas}, \textbf{sumideros},
\textbf{cámaras de observación} y \textbf{colectores de aguas lluvias}.

\begin{figure}[H]
\centering
\begin{tikzpicture}[font=\small]
  \fill[gray!50] (0,0) -- (0,1.8) -- (0.3,1.8) -- (0.3,0.9) -- (4,1.8) -- (7.7,0.9) -- (7.7,1.8) -- (8,1.8) -- (8,0) -- cycle;
  \fill[aguaclara] (0.3,0.9) -- (0.3,1.25) -- (1.7,1.25) -- cycle;
  \fill[aguaclara] (7.7,0.9) -- (7.7,1.25) -- (6.3,1.25) -- cycle;
  \fill[sueloclaro] (0.5,-1.3) rectangle (7.5,0);
  \draw[thick,fill=white] (4,-0.6) circle (0.5);
  \begin{scope}\clip (4,-0.6) circle (0.5); \fill[aguaclara] (3.4,-1.2) rectangle (4.6,-0.75);\end{scope}
  \draw[thick] (4,-0.6) circle (0.5);
  \draw (1.0,0.9) -- (1.8,0.9); \draw (1.8,0.9) arc (0:14:0.8); \node at (2.1,1.0) {$i$};
  \draw[<-] (0.7,1.1) -- (-0.8,2.2) node[above,azulusm,align=center]{Pendiente transversal\\o bombeo};
  \draw[<-] (4.5,-0.8) -- (8.2,-1.0) node[right,azulusm,align=left]{Colector de\\aguas lluvias};
  \draw[->,azulusm,thick] (1.2,1.0) .. controls (2.2,0.2) .. (3.5,-0.5);
  \draw[->,azulusm,thick] (6.8,1.0) .. controls (5.8,0.2) .. (4.5,-0.5);
\end{tikzpicture}
\caption{Perfil transversal de una calle: cunetas a ambos lados (bombeo $i$)
que descargan, vía sumideros, al colector.}
\end{figure}

\section{Tipos de obras}
Las obras dependen de la capacidad de conducción ($Q$) requerida:
\begin{itemize}
  \item Cuneta simple.
  \item Cuneta compuesta.
  \item Colector de aguas lluvias (con sumideros y cámaras).
\end{itemize}

% Macro sección de cuneta: #1=1 compuesta
\newcommand{\seccioncuneta}[1]{%
\begin{tikzpicture}[font=\small]
  \fill[gray!55] (0,0) rectangle (0.3,1.7);
  \ifnum#1=1
    \fill[gray!40] (0.3,0) -- (0.3,0.1) -- (1.8,0.9) -- (5.5,1.4) -- (6.3,1.25) -- (6.3,0) -- cycle;
    \fill[aguaclara] (0.3,0.1) -- (0.3,1.1) -- (3.5,1.1) -- (1.8,0.9) -- cycle;
    \draw (0.3,0.1) -- (1.8,0.9) -- (5.5,1.4) -- (6.3,1.25);
    \draw[dashed] (1.8,0.9) -- (5.8,0.9) (1.8,-0.2) -- (1.8,0.9);
    \draw[cota] (0.3,-0.1) -- (1.8,-0.1) node[midway,below]{$w$};
    \node at (1.3,0.3) {$i_1$}; \node at (4.9,1.1) {$i_2$};
    \draw[cota] (0.3,1.35) -- (3.5,1.35) node[midway,above]{$b$};
  \else
    \fill[gray!40] (0.3,0) -- (0.3,0.4) -- (5.5,1.2) -- (6.3,1.05) -- (6.3,0) -- cycle;
    \fill[aguaclara] (0.3,0.4) -- (0.3,0.85) -- (3.2,0.85) -- cycle;
    \draw (0.3,0.4) -- (5.5,1.2) -- (6.3,1.05);
    \draw[dashed] (0.3,0.4) -- (5.8,0.4);
    \node at (4.9,0.6) {$i$};
    \draw[cota] (0.3,1.1) -- (3.2,1.1) node[midway,above]{$b$};
    \draw[cota] (0.1,0.4) -- (0.1,0.85) node[midway,left]{$y$};
    \draw[cota] (-0.4,0.4) -- (-0.4,1.7) node[midway,left]{$h_s$};
  \fi
  \draw[<-] (5.5,1.45) -- (5.9,1.9) node[above,azulusm]{Bombeo};
\end{tikzpicture}}

\section{Cuneta simple}
Conformada por la calle misma y la solera.
\begin{itemize}
  \item Destinada para caudales pequeños.
  \item Para su diseño se debe determinar la pendiente transversal de la
        calle ($i$, bombeo) en función de:
  \begin{itemize}
    \item Ancho máximo permitido ($b_{max}=1-2$ [m]).
    \item Pendiente longitudinal de la calle ($S=3\%-10\%$).
    \item Revisar Tabla 5.3.4 del MDU.
  \end{itemize}
\end{itemize}

\begin{figure}[H]
\centering
\seccioncuneta{0}
\caption{Corte transversal de una cuneta simple: $y=b\cdot i\leq h_s\approx15$ [cm].}
\end{figure}

Utilizando la ecuación de Manning:
\[
Q=\frac{\sqrt S}{n}\Omega\cdot R_H^{2/3},\qquad
\Omega=\frac{b\cdot y}{2}=\frac{b^2\cdot i}{2},\qquad R_H=\frac\Omega\Psi,\qquad
\Psi=b\cdot i+b\sqrt{i^2+1}\approx b(i+1)
\]
\begin{formula}
$\displaystyle Q=\frac{\sqrt S}{n}\,\frac{b^2\cdot i}{2}\left(\frac{b\cdot i}{2(i+1)}\right)^{2/3}$\qquad $i$ varía entre 2 -- 4\%
\end{formula}

\section{Cuneta compuesta}
Se usan cuando el ancho de inundación está en el rango de 1 a 2 [m] para una
cuneta simple y se requiere mayor capacidad.
\begin{itemize}
  \item Generalmente utilizadas en carreteras.
  \item El ancho superficial ($b$) se calcula con la fórmula de Manning,
        teniendo en cuenta que $b\leq1-2$ [m].
  \item Revisar Tabla 5.3.5 del MDU.
  \item Si $b<w$, el cálculo se debe realizar como cuneta simple.
\end{itemize}

\begin{figure}[H]
\centering
\seccioncuneta{1}
\caption{Corte transversal de una cuneta compuesta (con zarpa de ancho $w$).}
\end{figure}

\begin{formula}
$\displaystyle Q=0.315\frac{\sqrt S}{n}b^{8/3}\left[i_2+(i_1-i_2)\left(\frac wb\right)^2\right]^{5/3}$
\end{formula}
con $w$: ancho de la zarpa ($w\approx0.6$ [m]); $S$: pendiente longitudinal
de la calle; $i_2$: bombeo de la calle (2 -- 4\%); $i_1$: pendiente
transversal de la zarpa (2 -- 10\%).

\section{Sumideros}
Obras de arte que se utilizan para captar el agua de la calle hacia el colector.
\begin{table}[H]
\centering\small
\begin{tabularx}{\textwidth}{lX}
\toprule
\textbf{Tipo} & \textbf{Características}\\
\midrule
Horizontal de fondo (tipo S3 y S4, SERVIU) & Usado para un amplio rango de pendientes. Inconvenientes: peatones, ciclistas, obstrucción.\\
Lateral (tipo S2, SERVIU) & Funciona admitiendo basuras; su capacidad decrece para pendientes longitudinales mayores al 3\%, donde se recomienda cuneta con pendiente longitudinal para mayor captación.\\
Mixto (tipo S1 y S2, SERVIU) & Funciona admitiendo basuras; capacidad de captación aumenta 10\% con respecto al sumidero horizontal.\\
\bottomrule
\end{tabularx}
\end{table}

\begin{figure}[H]
\centering
\begin{tikzpicture}[font=\small,x={(1cm,0cm)},y={(0.5cm,0.35cm)},z={(0cm,1cm)}]
  % losa
  \draw[fill=gray!15] (0,0,0) -- (5,0,0) -- (5,4,0) -- (0,4,0) -- cycle;
  \draw[fill=gray!35] (0,4,0) -- (5,4,0) -- (5,4,0.6) -- (0,4,0.6) -- cycle;
  % rejilla
  \draw[fill=gray!60] (1.2,2.4,0) -- (3.8,2.4,0) -- (3.8,3.6,0) -- (1.2,3.6,0) -- cycle;
  \foreach \x in {1.45,1.7,...,3.6} \draw[white] (\x,2.45,0) -- (\x,3.55,0);
  \draw[cota] (1.2,2.1,0) -- (3.8,2.1,0) node[midway,below]{$L$};
  \draw[cota] (4.1,2.4,0) -- (4.1,3.6,0) node[midway,right]{$w$};
\end{tikzpicture}
\caption{Sumidero horizontal de fondo junto a la solera: reja de largo $L$ y
ancho $w$ ($w=0.4\sim0.7$ [m], $L=1\sim2$ [m]).}
\end{figure}

\subsection{Sumidero horizontal de fondo}
El caudal máximo de captación depende del funcionamiento:
\begin{enumerate}
  \item \textbf{Actuando como vertedero} (alturas de agua pequeñas): un
        vertedero es un punto de control; produce un peralte del
        escurrimiento aguas arriba.
  \begin{formula}
  $\displaystyle Q_m=1.66\,(L_e+2\cdot w_e)\,h^{1.5}\ [\text{m}^3/\text{s}]\qquad\text{si } h<1.6\frac{A_e}{L_e+2w_e}$
  \end{formula}
  \item \textbf{Actuando como orificio} (alturas de agua mayores): el
        sumidero se ahoga.
  \begin{formula}
  $\displaystyle Q_m=2.66\cdot A_e\cdot h^{0.5}\ [\text{m}^3/\text{s}]\qquad\text{si } h\geq1.6\frac{A_e}{L_e+2w_e}$
  \end{formula}
\end{enumerate}
donde $w_e=w-e\cdot n_L$: ancho efectivo [m]; $L_e=L-e\cdot n_T$: largo
efectivo [m]; $A_e=L_e\cdot w_e$: área efectiva [m$^2$]; $n_L$: número de
barras longitudinales; $n_T$: número de barras transversales; $e$: diámetro o
espesor de las barras [m]; $h$: altura de agua en la calle frente al sumidero [m].

\paragraph{Eficiencia.} El sumidero no siempre logra captar toda el agua que
circula por la cuneta:
\begin{formula}
$Q_s=\begin{cases}\eta_HQ & \eta_HQ<Q_m\\ Q_m & \eta_HQ\geq Q_m\end{cases}$
\end{formula}
\[
\eta_H=R_f\cdot\eta_0+R_s(1-\eta_0),\qquad
\eta_0=1-\left(1-\frac w{b_s}\right)^{2.67},
\]
\[
R_s=\frac{1}{1+\dfrac{0.0828\,v^{1.8}}{i\cdot L^{2.3}}},\qquad
R_f=1-0.295(v-v_0)
\]
todos acotados entre 0 y 1; $\eta_H$: eficiencia del sumidero de fondo;
$Q_s$: caudal captado [m$^3$/s]; $Q$: caudal que circula por la cuneta
[m$^3$/s]; $v$: velocidad del agua en la calle frente al sumidero [m/s];
$v_0$: velocidad a la que ocurre el primer fenómeno de salpicadura en la reja [m/s].

\subsection{Sumidero lateral}
\begin{formula}
$Q_m=\begin{cases}1.27\cdot L\cdot h^{1.5} & h<a\quad\text{(vertedero)}\\ 2.66\cdot L\cdot a\cdot h^{0.5} & h\geq a\quad\text{(ahogado)}\end{cases}$
\end{formula}
Recomendable su uso sólo para pendientes longitudinales $S<3\%$.

\paragraph{Eficiencia}
\begin{formula}
$Q_s=\begin{cases}\eta_LQ & \eta_LQ<Q_m\\ Q_m & \eta_LQ\geq Q_m\end{cases}\qquad
\eta_L=\begin{cases}1 & h>a\\ 1-\left(1-\dfrac{L}{L_T}\right)^{1.8} & h\leq a\end{cases}$
\end{formula}
\[
\begin{aligned}
&\text{Cuneta compuesta:} && L_T=0.81\,Q^{0.42}\cdot S^{0.3}(n\cdot i)^{-0.6}\ [\text{m}]\\
&\text{Cuneta simple } (w=0): && L_T=0.81\,Q^{0.42}\cdot S^{0.3}(n\cdot i_2)^{-0.6}\ [\text{m}]
\end{aligned}
\]
con $n$: rugosidad de Manning; $S$: pendiente longitudinal de la calle;
$L_T$: largo necesario para captar todo el caudal ($L_{T,min}=L$).

\begin{figure}[H]
\centering
\begin{tikzpicture}[font=\small]
  \fill[gray!30] (0,0) rectangle (2.5,5);
  \fill[gray!55] (-0.2,0) rectangle (0,5); \fill[gray!55] (2.5,0) rectangle (2.7,1.6); \fill[gray!55] (2.5,3.0) rectangle (2.7,5);
  \draw[dashed,thick] (1.2,0) -- (1.2,5);
  \draw[flecha,azulusm] (2.0,4.9) -- (2.0,4.0) node[pos=0,right]{$Q$};
  \foreach \y in {3.6,3.1,2.6,2.1} \draw[->,azulusm] (2.0,\y+0.3) .. controls (2.3,\y) .. (2.5,\y-0.1);
  \draw[flecha,azulusm,line width=2pt] (2.5,2.3) -- (3.3,2.3) node[right]{$\eta_LQ$};
  \draw[cota] (2.9,1.6) -- (2.9,3.0) node[pos=0.8,right]{$L$};
  \draw[cota] (4.2,1.2) -- (4.2,3.6) node[midway,right]{$L_T$};
  \draw[dashed] (2.7,1.2) -- (4.3,1.2) (2.7,3.6) -- (4.3,3.6);
  \node[below] at (1.8,0) {$(1-\eta_L)Q$};
\end{tikzpicture}
\caption{Sumidero lateral (vista en planta): capta $\eta_LQ$ y deja pasar $(1-\eta_L)Q$.}
\end{figure}

\subsection{Sumideros mixtos}
\begin{itemize}
  \item Combinación de ambos tipos (de fondo y lateral).
  \item La experiencia indica que sólo existe un pequeño aumento del caudal
        de captación respecto a sumideros de fondo ($\approx10\%$).
  \item Se recomienda su uso en zonas con muchos árboles de hoja caduca, ya
        que de emergencia puede funcionar el sumidero lateral.
\end{itemize}
\begin{formula}
$Q_s=\begin{cases}(\eta_H+\eta_L)Q & (\eta_H+\eta_L)Q<Q_m\\ Q_m & (\eta_H+\eta_L)Q\geq Q_m\end{cases}$
\end{formula}

\subsection{Factor de seguridad de eficiencia}
Las eficiencias empíricas o calculadas suponen que el sumidero está limpio.
En la práctica esta situación rara vez ocurre; por lo tanto, es importante
incluir un factor de seguridad $F$:
\begin{formula}
$\eta_d=F\cdot\eta$
\end{formula}
\begin{table}[H]
\centering
\begin{tabular}{lcc}
\toprule
\textbf{Caudal aportante} & \textbf{Mantención y limpieza} & \textbf{Factor de seguridad}\\
\midrule
Inferior a 30 l/s & Buena & 1.0\\
 & Mala & 0.8\\
Superior a 30 l/s & Buena & 0.7\\
 & Mala & 0.6\\
\bottomrule
\end{tabular}
\caption{Estimación de un factor de seguridad para sumideros (MDU, Tabla 6.5.13).}
\end{table}
Caudales altos tienden a llevar más basura y por lo tanto a taparse; por ello
se recomienda que siempre se haga con ventana lateral, así ésta no se tenga
en cuenta para el cálculo de la eficiencia.

\subsection{Distancia entre sumideros}
\begin{itemize}
  \item Se debe garantizar que el flujo no pase las condiciones de seguridad.
  \item Para la red primaria, el caudal máximo en la calle debe calcularse
        mediante el modelo lluvia--escorrentía SWMM.
  \item Para la red secundaria, el flujo que llega a la calle se puede
        estimar por el método racional.
\end{itemize}
\begin{formula}
$\displaystyle Q_{captado}=\frac{C\cdot i\cdot A}{3600}=\eta_dQ_{max}
\quad\Longrightarrow\quad
L=\frac{A}{b_c}=\frac{3600\,(F\cdot\eta)\,Q_{max}}{C\cdot i\cdot b_c}$
\end{formula}
\[
\text{Calzada con doble bombeo: } b_c=0.5\cdot e+w_c;\qquad\text{calle con bombeo único: } b_c=e+w_c
\]
donde $L$: separación entre sumideros [m]; $A$: área aportante [m$^2$];
$w_c$: ancho de los terrenos aportantes [m]; $e$: ancho de la calle [m]; $C$:
coeficiente de escorrentía; $i$: intensidad de diseño [mm/h] ($T_r=2$ años,
$D_{min}=30$ minutos).

\subsection{Consideraciones de diseño (cunetas y sumideros)}
\begin{itemize}
  \item Para tormentas menores ($T_r=2$ a 10 años), la calzada no podrá
        inundarse más del ancho recomendado. Para el cálculo del tiempo de
        concentración se pueden utilizar las fórmulas vistas en hidrología,
        pero se debe verificar $t_c\geq30$ [min].
  \item Para tormentas mayores ($T_r=50$ a 100 años), la calzada no podrá
        inundarse con altura mayor a la solera y en todos los casos $h<15$ cm.
\end{itemize}

\section{Colectores de aguas lluvias}
\begin{itemize}
  \item Se utilizan en el caso de que $b>1-2$ [m] (excede lo permitido para una cuneta).
  \item Su diseño considera caudales ($Q$) pequeños y medianos.
  \item Se utilizan tuberías de diámetro mínimo de 300 [mm] y máximo 2 [m]:
        tubos o cajones de cemento comprimido (C.C.C.) o tubos de HDPE.
  \item MDU 6.7.2: colectores de diámetro pequeño ($D<1.2$ m), $T_r=2-10$ años.
  \item MDU 6.7.2: colectores de diámetro mayor ($D>1.2$ m), $T_r=5-10$ años.
\end{itemize}

La elección del tamaño queda definida por el caudal y los requisitos de
velocidad del flujo. La altura se puede calcular con la fórmula de Manning:
\[
Q=\frac{\sqrt S}{n}A\cdot R_H^{2/3}
\]
\begin{itemize}
  \item Bajo condiciones normales se debe verificar que $h\leq0.8D$.
  \item Para el diseño se considera la condición crítica ($h=0.8D$).
  \item Bajo condiciones excepcionales podría exceder el límite de $0.8D$
        (p.~ej.\ $T_r=50$ años), pero siempre sin presión.
\end{itemize}

\begin{figure}[H]
\centering
\begin{minipage}{0.35\textwidth}\centering
\begin{tikzpicture}[font=\small]
  \begin{scope}\clip (0,0) circle (1.2); \fill[aguaclara] (-1.3,-1.3) rectangle (1.3,0.72);\end{scope}
  \draw[thick] (0,0) circle (1.2);
  \draw (-0.96,0.72) -- (0.96,0.72);
  \draw (-0.96,0.72) -- (0,0) -- (0.96,0.72);
  \draw[->] (0,0) ++(-127:0.4) arc (-127:127:0.4) node[pos=0.5,right]{$\theta$};
  \pic at (0,0.74) {nivelagua};
  \draw[cota] (-1.6,-1.2) -- (-1.6,1.2) node[midway,left]{$D$};
  \draw[cota] (1.5,-1.2) -- (1.5,0.72) node[midway,right]{$h$};
\end{tikzpicture}
\end{minipage}
\begin{minipage}{0.55\textwidth}
\[
\begin{gathered}
A=\frac{(\theta-\operatorname{sen}\theta)\,D^2}{8},\qquad
R_h=\left(1-\frac{\operatorname{sen}\theta}{\theta}\right)\frac D4\\
\theta=2\cos^{-1}\left(1-\frac{2h}{D}\right)
\end{gathered}
\]
\end{minipage}
\caption{Sección circular parcialmente llena.}
\end{figure}

\begin{table}[H]
\centering
\begin{tabular}{lccccc}
\toprule
\textbf{Situación} & $h$ [m] & $\theta$ [grados] & $A_m$ [m$^2$] & $P_m$ [m] & $R_h$ [m]\\
\midrule
Sección completa & $D$ & 360 & $0.785D^2$ & $3.142D$ & $0.250D$\\
Caudal máximo & $0.95D$ & 308 & $0.771D^2$ & $2.686D$ & $0.287D$\\
Sección diseño & $0.80D$ & 254 & $0.674D^2$ & $2.217D$ & $0.304D$\\
\bottomrule
\end{tabular}
\caption{Propiedades de secciones especiales en ductos circulares (MDU, Tabla 6.7.13).}
\end{table}

Además, se debe verificar:
\begin{itemize}
  \item \textbf{Velocidad mínima} (evitar depósito de sedimentos dentro de la
        tubería, autolavado):
  \begin{formula}
  $\displaystyle V_c=0.397\,\frac{D^{2/3}\sqrt S}{n}\geq0.6\ [\text{m/s}]$
  \end{formula}
  \item \textbf{Velocidad máxima} (evitar socavación):
\end{itemize}
\begin{table}[H]
\centering
\begin{tabular}{lc}
\toprule
\textbf{Tipo de tubería} & \textbf{Velocidad máxima} [m/s]\\
\midrule
Concreto simple hasta 45 cm de diámetro & 3.0\\
Concreto reforzado de 61 cm de diámetro o mayores & 3.5\\
Fibrocemento & 5.0\\
Policloruro de vinilo (PVC) & 5.0\\
Polietileno de alta densidad & 5.0\\
Tubería de acero corrugado & 4.5\\
Tubería de acero corrugado revestido en shotcrete & 3.0\\
Polímero reforzado con fibra de vidrio (PRFV) & 3.5\\
\bottomrule
\end{tabular}
\caption{Velocidad máxima permisible en colectores de pequeño diámetro (MDU,
Tabla 6.7.14; fuente: Comisión Nacional del Agua, México, 2009).}
\end{table}

\section{Recomendaciones de diseño de colectores}
\paragraph{Pendientes longitudinales muy grandes.} Si se tienen pendientes
$S$ muy grandes, dentro de las soluciones para reducir la velocidad se encuentran:
\begin{itemize}
  \item \textbf{Cámaras de observación:} se colocan para escalonar el
        escurrimiento, de forma tal que las tuberías no registren velocidades
        mayores a lo permitido ($v_{max}$), con $S_D<S$.
  \item \textbf{Dientes disipadores de energía:} reducen la altura de
        velocidad; se usan en cajones de hormigón; ¡alto costo!
\end{itemize}

\begin{figure}[H]
\centering
\begin{tikzpicture}[font=\small]
  \fill[gray!50] (0,3.0) -- (9,1.2) -- (9,1.4) -- (0,3.2) -- cycle;
  \node[above] at (4.5,2.3) {$S$};
  \foreach \x/\yt in {2/2.6,5.5/1.9}{
    \fill[gray!25,draw] (\x,\yt) rectangle (\x+0.6,\yt-2.0);}
  \draw[thick] (0,1.6) -- (2,1.25) (0,1.8) -- (2,1.45);
  \draw[thick] (2.6,0.8) -- (5.5,0.35) (2.6,1.0) -- (5.5,0.55);
  \draw[thick] (6.1,-0.1) -- (9,-0.5) (6.1,0.1) -- (9,-0.3);
  \node at (1,1.2) {$S_D$}; \node at (4,0.3) {$S_D$}; \node at (7.5,-0.6) {$S_D$};
  \node[azulusm,align=center] at (8,2.4) {Cámaras de\\observación};
  \draw[->,azulusm] (7.2,2.3) -- (6.1,1.4);
  \node at (4,-0.8) {$S_D<S$};
\end{tikzpicture}
\caption{Cámaras de observación para escalonar un colector en una calle de gran pendiente.}
\end{figure}

\paragraph{Trampas de arena.} Se deben colocar si existen:
\begin{itemize}
  \item Cambios de dirección.
  \item Cambios de pendiente.
  \item Tramos largos sin ninguna perturbación en el flujo (cámaras cada
        $50\sim130$ [m]).
\end{itemize}
La trampa de arena es la zona de la cámara ubicada bajo la profundidad de la
tubería; retiene los sedimentos que trae el flujo y se debe limpiar cada
cierto tiempo.

\paragraph{Profundidad de instalación.} Colector de aguas lluvias a 1.8 [m]
bajo la superficie; red de agua potable a $1\sim1.5$ [m].

\paragraph{Diámetros mínimos.} Para posibilitar la limpieza y correcto
funcionamiento. Desde el punto de vista administrativo, existen dos tipos de
colectores:
\begin{itemize}
  \item Primarios ($D>0.8$ [m]): dependen de la Dirección de Obras Hidráulicas (DOH).
  \item Secundarios ($D\leq0.8$ [m]): dependen del SERVIU.
\end{itemize}

\section{Obtención de caudales de diseño}
Se tiene una serie de cuencas conectadas a través de una red de colectores.
Como las cuencas son de tamaño pequeño, se asumen condiciones meteorológicas
homogéneas, es decir, la precipitación es la misma en todas las cuencas.
Generalmente se utiliza la fórmula racional para el cálculo de $Q$.

\begin{figure}[H]
\centering
\begin{tikzpicture}[font=\small]
  \foreach \x/\y/\n in {0/2.2/1,3.5/1.8/2,7/1.5/3}{
    \fill[green!55!black!60,draw=green!30!black] (\x,\y) .. controls (\x-0.9,\y+0.6) and (\x-0.7,\y+1.8) .. (\x+0.1,\y+1.9)
      .. controls (\x+0.9,\y+2.0) and (\x+1.1,\y+0.9) .. (\x,\y);
    \node at (\x+1.0,\y+1.7) {$A_\n$};
    \fill[white,draw] (\x,\y) circle (0.07);}
  \draw[flecha,azulusm,thick] (0,2.2) -- (3.5,1.8) node[midway,above]{$Q_{D1},L_1$} node[midway,below=4pt,draw,circle,inner sep=1pt]{1};
  \draw[flecha,azulusm,thick] (3.5,1.8) -- (7,1.5) node[midway,above]{$Q_{D2},L_2$} node[midway,below=4pt,draw,circle,inner sep=1pt]{2};
  \draw[flecha,azulusm,thick] (7,1.5) -- (9.5,1.2) node[midway,above]{$Q_{D3},L_3$} node[midway,below=4pt,draw,circle,inner sep=1pt]{3} node[right]{Río};
\end{tikzpicture}
\caption{Serie de cuencas conectadas por colectores.}
\end{figure}

\paragraph{Colector 1:} sólo aporta caudal la cuenca 1:
\[
Q_1=\frac{C_1\cdot\bar\imath(t_c)\cdot A_1}{3.6}\quad\Longrightarrow\quad Q_{D1}=Q_1
\]

\paragraph{Colector 2:} las cuencas 1 y 2 aportan caudal al colector 2. El
caudal de diseño corresponde a:
\begin{formula}
$Q_{D2}=\max(Q_{D1},Q_2,Q_{2I})$
\end{formula}
donde $Q_{D1}$: caudal de diseño del colector 1; $Q_2$: caudal máximo
aportado por la cuenca 2; $Q_{2I}$: caudal máximo aportado por las cuencas 1 y 2:
\[
Q_{2I}=\frac{\bar C_{1,2}\cdot\bar\imath(t_{c2I})\cdot(A_1+A_2)}{3.6},\qquad
\bar C_{1,2}=\frac{C_1+C_2}{2},\qquad
t_{c2I}=\max(t_{c1}+t_{t1},\,t_{c2}),\qquad t_{t1}=\frac{L_1}{v_1}
\]

\paragraph{Colector 3:} el procedimiento es similar al caso del colector 2:
\begin{formula}
$Q_{D3}=\max(Q_{D2},Q_3,Q_{3I})$
\end{formula}
\[
Q_{3I}=\frac{\bar C_{1,2,3}\cdot\bar\imath(t_{c3I})\cdot(A_1+A_2+A_3)}{3.6},\qquad
\bar C_{1,2,3}=\frac{C_1+C_2+C_3}{3},
\]
\[
t_{c3I}=\max(t_{c2I}+t_{t2},\,t_{c3}),\qquad t_{t2}=\frac{L_2}{v_2}
\]

\paragraph{Hidrograma unitario triangular del SCS.} Si se dispone de la
distribución temporal de la precipitación (hietograma), una alternativa es
utilizar el hidrograma de crecidas del SCS:
\begin{formula}
$\displaystyle t_p=\frac{t_{LL}}{2}+0.6\cdot t_c\ \ [\text{h}],\qquad t_r=1.67\cdot t_p\ \ [\text{h}],\qquad t_B=2.67\cdot t_p\ \ [\text{h}]$
\end{formula}
Igualando el volumen de escorrentía directa con el área del hidrograma:
\[
V_{ed}=P_{ef}\cdot A=\frac12Q_p\cdot t_B
\quad\Longrightarrow\quad
\boxed{Q_p=\frac{2\cdot P_{ef}\cdot A}{2.67\left(\frac{t_{LL}}{2}+0.6\cdot t_c\right)}}
\]

\begin{figure}[H]
\centering
\begin{tikzpicture}[font=\small]
  \draw[flecha] (0,0) -- (7,0) node[right]{$t$};
  \draw[flecha] (0,0) -- (0,3) node[left]{$Q$};
  \fill[azulusm] (0,3.4) rectangle (0.6,3.2); \node[right] at (0.6,3.3) {$P_{ef}$ (duración $t_{LL}$)};
  \draw[azulusm,thick] (0,0) -- (2.2,2.4) -- (6,0);
  \draw[dashed] (0,2.4) node[left]{$Q_p$} -- (2.2,2.4) -- (2.2,0);
  \draw[cota] (0,-0.3) -- (2.2,-0.3) node[midway,below]{$t_p$};
  \draw[cota] (2.2,-0.3) -- (6,-0.3) node[midway,below]{$t_r$};
  \draw[cota] (0,-0.9) -- (6,-0.9) node[midway,below]{$t_B$};
\end{tikzpicture}
\caption{Hidrograma unitario triangular del SCS.}
\end{figure}


\end{document}
